% =============================================================================
%  Cuerpo del Tema 5. Energía y fuerzas (primera parte)
% =============================================================================
\hacerportada

% ------------------------------------------------------------- Página de resumen
\thispagestyle{empty}
\begin{center}
  {\small\color{tenue}\espaciado{Electromagnetismo I}\quad$\cdot$\quad\espaciado{Tema 5}}\\[1.6em]
  {\fuentetitulo\fontsize{21}{25}\selectfont Energía y fuerzas}\\[0.4em]
  {\large\itshape Primera parte}\\[0.9em]
  {\color{acentoA}\rule{2.6cm}{0.6pt}}\\[1.1em]
  {\large Daniel Martín Guerrero}\\[0.3em]
  {\small\itshape\color{tenue} Grado en Física, Facultad de Ciencias, Universidad de Salamanca}
\end{center}
\vspace{1.6em}

\noindent
\begin{minipage}[t]{0.6\textwidth}
\small
{\color{acentoA}\espaciado{Resumen.}}\enspace
Definimos la energía electrostática de un sistema como el trabajo reversible necesario para formarlo trayendo
sus cargas desde el infinito. Para cargas puntuales obtenemos la energía de interacción
$\frac12\sum_iq_i\phi_i$; para distribuciones continuas, la energía total $\frac12\int\rho\phi\,\dd V$, que el teorema de la
divergencia transforma en $\frac{\eo}{2}\int E^2\dd V$: la energía queda almacenada en el campo, con densidad
$u_e=\frac{\eo}{2}E^2$, es siempre positiva y no es aditiva. En un sistema de conductores la energía es
$\frac12\sum_iQ_i\phi_i$, y para un condensador, $\frac12QV=\frac12CV^2=Q^2/2C$. Tras demostrar el teorema de Earnshaw,
calculamos las fuerzas sobre conductores con el método de los trabajos virtuales, a carga constante,
$\vF_e=-\grad U_e|_Q$, y a potencial constante, $\vF_e=+\grad U_e|_\phi$, y probamos que ambos conducen a la misma
fuerza. Los seminarios resuelven los problemas de esta parte del tema.
\end{minipage}\hfill
\begin{minipage}[t]{0.36\textwidth}\centering\vspace{-0.6em}
\begin{tikzpicture}
\begin{axis}[marco, width=6.0cm, xmin=-2.6, xmax=2.6, ymin=-1.8, ymax=1.8, enlargelimits=false]
  \addplot graphics[xmin=-2.6, xmax=2.6, ymin=-1.8, ymax=1.8] {t5-ue-opuestas-\modotema};
  %% Generado por scripts/tema5.py -- no editar a mano
\addplot[lineaE, flechaen=0.419] coordinates {(2.600,1.027) (2.568,1.002) (2.537,0.977) (2.505,0.952) (2.473,0.928) (2.441,0.903) (2.409,0.879) (2.376,0.855) (2.344,0.831) (2.311,0.807) (2.279,0.784) (2.246,0.760) (2.213,0.737) (2.180,0.714) (2.147,0.691) (2.114,0.668) (2.081,0.645) (2.048,0.622) (2.015,0.600) (1.981,0.578) (1.947,0.555) (1.914,0.533) (1.880,0.511) (1.846,0.490) (1.812,0.468) (1.778,0.447) (1.744,0.425) (1.710,0.404) (1.675,0.383) (1.641,0.362) (1.607,0.341) (1.572,0.321) (1.537,0.300) (1.503,0.280) (1.468,0.259) (1.433,0.239) (1.398,0.219) (1.363,0.199) (1.328,0.180) (1.293,0.160) (1.258,0.140) (1.222,0.121) (1.187,0.102) (1.152,0.082) (1.117,0.063)};
\addplot[lineaE, flechaen=0.581] coordinates {(-1.117,0.063) (-1.152,0.082) (-1.187,0.102) (-1.222,0.121) (-1.258,0.140) (-1.293,0.160) (-1.328,0.180) (-1.363,0.199) (-1.398,0.219) (-1.433,0.239) (-1.468,0.259) (-1.503,0.280) (-1.537,0.300) (-1.572,0.321) (-1.607,0.341) (-1.641,0.362) (-1.675,0.383) (-1.710,0.404) (-1.744,0.425) (-1.778,0.447) (-1.812,0.468) (-1.846,0.490) (-1.880,0.511) (-1.914,0.533) (-1.947,0.555) (-1.981,0.578) (-2.015,0.600) (-2.048,0.622) (-2.081,0.645) (-2.114,0.668) (-2.147,0.691) (-2.180,0.714) (-2.213,0.737) (-2.246,0.760) (-2.279,0.784) (-2.311,0.807) (-2.344,0.831) (-2.376,0.855) (-2.409,0.879) (-2.441,0.903) (-2.473,0.928) (-2.505,0.952) (-2.537,0.977) (-2.568,1.002) (-2.600,1.027)};
\addplot[lineaE, flechaen=0.159] coordinates {(2.112,1.800) (2.099,1.761) (2.086,1.723) (2.072,1.684) (2.058,1.646) (2.043,1.608) (2.028,1.570) (2.013,1.532) (1.997,1.495) (1.981,1.457) (1.964,1.420) (1.948,1.383) (1.930,1.346) (1.913,1.309) (1.895,1.272) (1.877,1.235) (1.858,1.199) (1.840,1.163) (1.820,1.127) (1.801,1.091) (1.781,1.055) (1.761,1.020) (1.741,0.984) (1.720,0.949) (1.699,0.914) (1.677,0.880) (1.656,0.845) (1.634,0.811) (1.611,0.777) (1.589,0.743) (1.566,0.709) (1.543,0.675) (1.519,0.642) (1.496,0.608) (1.472,0.575) (1.447,0.542) (1.423,0.510) (1.398,0.477) (1.373,0.445) (1.348,0.413) (1.323,0.381) (1.297,0.349) (1.272,0.317) (1.246,0.286) (1.220,0.254) (1.193,0.223) (1.167,0.192) (1.140,0.161) (1.114,0.130) (1.087,0.099)};
\addplot[lineaE, flechaen=0.841] coordinates {(-1.087,0.099) (-1.114,0.130) (-1.140,0.161) (-1.167,0.192) (-1.193,0.223) (-1.220,0.254) (-1.246,0.286) (-1.272,0.317) (-1.297,0.349) (-1.323,0.381) (-1.348,0.413) (-1.373,0.445) (-1.398,0.477) (-1.423,0.510) (-1.447,0.542) (-1.472,0.575) (-1.496,0.608) (-1.519,0.642) (-1.543,0.675) (-1.566,0.709) (-1.589,0.743) (-1.611,0.777) (-1.634,0.811) (-1.656,0.845) (-1.677,0.880) (-1.699,0.914) (-1.720,0.949) (-1.741,0.984) (-1.761,1.020) (-1.781,1.055) (-1.801,1.091) (-1.820,1.127) (-1.840,1.163) (-1.858,1.199) (-1.877,1.235) (-1.895,1.272) (-1.913,1.309) (-1.930,1.346) (-1.948,1.383) (-1.964,1.420) (-1.981,1.457) (-1.997,1.495) (-2.013,1.532) (-2.028,1.570) (-2.043,1.608) (-2.058,1.646) (-2.072,1.684) (-2.086,1.723) (-2.099,1.761) (-2.112,1.800)};
\addplot[lineaE, flechaen=0.500] coordinates {(1.459,1.800) (1.461,1.759) (1.463,1.719) (1.464,1.678) (1.465,1.637) (1.464,1.597) (1.463,1.556) (1.462,1.515) (1.460,1.474) (1.457,1.434) (1.454,1.393) (1.450,1.353) (1.445,1.312) (1.440,1.272) (1.434,1.232) (1.427,1.191) (1.420,1.151) (1.413,1.111) (1.405,1.071) (1.396,1.032) (1.387,0.992) (1.377,0.952) (1.367,0.913) (1.356,0.874) (1.345,0.835) (1.333,0.796) (1.321,0.757) (1.309,0.718) (1.296,0.679) (1.282,0.641) (1.268,0.603) (1.254,0.565) (1.239,0.527) (1.224,0.489) (1.209,0.451) (1.193,0.414) (1.177,0.376) (1.160,0.339) (1.144,0.302) (1.127,0.265) (1.110,0.228) (1.092,0.191) (1.075,0.154) (1.057,0.117)};
\addplot[lineaE, flechaen=0.500] coordinates {(-1.057,0.117) (-1.075,0.154) (-1.092,0.191) (-1.110,0.228) (-1.127,0.265) (-1.144,0.302) (-1.160,0.339) (-1.177,0.376) (-1.193,0.414) (-1.209,0.451) (-1.224,0.489) (-1.239,0.527) (-1.254,0.565) (-1.268,0.603) (-1.282,0.641) (-1.296,0.679) (-1.309,0.718) (-1.321,0.757) (-1.333,0.796) (-1.345,0.835) (-1.356,0.874) (-1.367,0.913) (-1.377,0.952) (-1.387,0.992) (-1.396,1.032) (-1.405,1.071) (-1.413,1.111) (-1.420,1.151) (-1.427,1.191) (-1.434,1.232) (-1.440,1.272) (-1.445,1.312) (-1.450,1.353) (-1.454,1.393) (-1.457,1.434) (-1.460,1.474) (-1.462,1.515) (-1.463,1.556) (-1.464,1.597) (-1.465,1.637) (-1.464,1.678) (-1.463,1.719) (-1.461,1.759) (-1.459,1.800)};
\addplot[lineaE, flechaen=0.500] coordinates {(0.904,1.800) (0.924,1.765) (0.943,1.728) (0.961,1.692) (0.979,1.655) (0.995,1.617) (1.010,1.579) (1.024,1.541) (1.037,1.502) (1.050,1.463) (1.061,1.424) (1.071,1.385) (1.081,1.345) (1.089,1.305) (1.097,1.265) (1.103,1.225) (1.109,1.184) (1.114,1.144) (1.119,1.103) (1.122,1.062) (1.125,1.022) (1.127,0.981) (1.128,0.940) (1.128,0.899) (1.128,0.858) (1.127,0.817) (1.125,0.777) (1.123,0.736) (1.120,0.695) (1.117,0.654) (1.112,0.614) (1.108,0.573) (1.103,0.533) (1.097,0.492) (1.091,0.452) (1.084,0.411) (1.078,0.371) (1.070,0.331) (1.063,0.291) (1.055,0.251) (1.046,0.211) (1.038,0.171) (1.029,0.131)};
\addplot[lineaE, flechaen=0.500] coordinates {(-1.029,0.131) (-1.038,0.171) (-1.046,0.211) (-1.055,0.251) (-1.063,0.291) (-1.070,0.331) (-1.078,0.371) (-1.084,0.411) (-1.091,0.452) (-1.097,0.492) (-1.103,0.533) (-1.108,0.573) (-1.112,0.614) (-1.117,0.654) (-1.120,0.695) (-1.123,0.736) (-1.125,0.777) (-1.127,0.817) (-1.128,0.858) (-1.128,0.899) (-1.128,0.940) (-1.127,0.981) (-1.125,1.022) (-1.122,1.062) (-1.119,1.103) (-1.114,1.144) (-1.109,1.184) (-1.103,1.225) (-1.097,1.265) (-1.089,1.305) (-1.081,1.345) (-1.071,1.385) (-1.061,1.424) (-1.050,1.463) (-1.037,1.502) (-1.024,1.541) (-1.010,1.579) (-0.995,1.617) (-0.979,1.655) (-0.961,1.692) (-0.943,1.728) (-0.924,1.765) (-0.904,1.800)};
\addplot[lineaE, flechaen=0.386, flechaen=0.614] coordinates {(-1.000,0.131) (-0.999,0.171) (-0.999,0.211) (-0.998,0.251) (-0.997,0.291) (-0.996,0.331) (-0.994,0.371) (-0.991,0.411) (-0.989,0.451) (-0.986,0.491) (-0.982,0.531) (-0.978,0.571) (-0.973,0.611) (-0.967,0.651) (-0.961,0.690) (-0.954,0.730) (-0.946,0.769) (-0.937,0.808) (-0.928,0.847) (-0.918,0.886) (-0.907,0.925) (-0.895,0.963) (-0.882,1.001) (-0.868,1.038) (-0.853,1.076) (-0.837,1.113) (-0.821,1.149) (-0.803,1.185) (-0.784,1.220) (-0.764,1.255) (-0.742,1.289) (-0.720,1.322) (-0.697,1.355) (-0.672,1.386) (-0.646,1.417) (-0.620,1.447) (-0.592,1.476) (-0.563,1.503) (-0.532,1.530) (-0.501,1.555) (-0.469,1.579) (-0.436,1.601) (-0.401,1.622) (-0.366,1.641) (-0.330,1.659) (-0.293,1.675) (-0.256,1.689) (-0.218,1.701) (-0.179,1.711) (-0.140,1.719) (-0.100,1.726) (-0.060,1.730) (-0.020,1.732) (0.020,1.732) (0.060,1.730) (0.100,1.726) (0.140,1.719) (0.179,1.711) (0.218,1.701) (0.256,1.689) (0.293,1.675) (0.330,1.659) (0.366,1.641) (0.401,1.622) (0.436,1.601) (0.469,1.579) (0.501,1.555) (0.532,1.530) (0.563,1.503) (0.592,1.476) (0.620,1.447) (0.646,1.417) (0.672,1.386) (0.697,1.355) (0.720,1.322) (0.742,1.289) (0.764,1.255) (0.784,1.220) (0.803,1.185) (0.821,1.149) (0.837,1.113) (0.853,1.076) (0.868,1.038) (0.882,1.001) (0.895,0.963) (0.907,0.925) (0.918,0.886) (0.928,0.847) (0.937,0.808) (0.946,0.769) (0.954,0.730) (0.961,0.690) (0.967,0.651) (0.973,0.611) (0.978,0.571) (0.982,0.531) (0.986,0.491) (0.989,0.451) (0.991,0.411) (0.994,0.371) (0.996,0.331) (0.997,0.291) (0.998,0.251) (0.999,0.211) (0.999,0.171) (1.000,0.131)};
\addplot[lineaE, flechaen=0.353, flechaen=0.647] coordinates {(-0.971,0.128) (-0.962,0.167) (-0.952,0.207) (-0.943,0.246) (-0.933,0.285) (-0.922,0.323) (-0.911,0.362) (-0.900,0.401) (-0.889,0.439) (-0.876,0.477) (-0.863,0.516) (-0.850,0.553) (-0.836,0.591) (-0.821,0.628) (-0.805,0.665) (-0.788,0.702) (-0.770,0.738) (-0.752,0.774) (-0.733,0.809) (-0.712,0.843) (-0.691,0.877) (-0.668,0.911) (-0.645,0.943) (-0.620,0.975) (-0.594,1.006) (-0.567,1.036) (-0.539,1.064) (-0.510,1.092) (-0.479,1.118) (-0.448,1.143) (-0.415,1.166) (-0.381,1.188) (-0.346,1.208) (-0.311,1.227) (-0.274,1.243) (-0.236,1.258) (-0.198,1.270) (-0.159,1.280) (-0.120,1.289) (-0.080,1.294) (-0.040,1.298) (-0.000,1.299) (0.040,1.298) (0.080,1.294) (0.120,1.289) (0.159,1.280) (0.198,1.270) (0.236,1.258) (0.274,1.243) (0.311,1.227) (0.346,1.208) (0.381,1.188) (0.415,1.166) (0.448,1.143) (0.479,1.118) (0.510,1.092) (0.539,1.064) (0.567,1.036) (0.594,1.006) (0.620,0.975) (0.645,0.943) (0.668,0.911) (0.691,0.877) (0.712,0.843) (0.733,0.809) (0.752,0.774) (0.770,0.738) (0.788,0.702) (0.805,0.665) (0.821,0.628) (0.836,0.591) (0.850,0.553) (0.863,0.516) (0.876,0.477) (0.889,0.439) (0.900,0.401) (0.911,0.362) (0.922,0.323) (0.933,0.285) (0.943,0.246) (0.952,0.207) (0.962,0.167) (0.971,0.128)};
\addplot[lineaE, flechaen=0.318, flechaen=0.682] coordinates {(-0.942,0.117) (-0.924,0.154) (-0.906,0.190) (-0.887,0.226) (-0.868,0.262) (-0.849,0.297) (-0.829,0.333) (-0.809,0.368) (-0.789,0.403) (-0.767,0.438) (-0.746,0.472) (-0.723,0.505) (-0.700,0.539) (-0.676,0.571) (-0.651,0.603) (-0.625,0.635) (-0.598,0.665) (-0.571,0.695) (-0.542,0.724) (-0.512,0.751) (-0.482,0.778) (-0.450,0.803) (-0.417,0.827) (-0.383,0.849) (-0.348,0.870) (-0.313,0.889) (-0.276,0.906) (-0.238,0.921) (-0.200,0.934) (-0.161,0.944) (-0.121,0.953) (-0.081,0.959) (-0.041,0.963) (-0.000,0.964) (0.041,0.963) (0.081,0.959) (0.121,0.953) (0.161,0.944) (0.200,0.934) (0.238,0.921) (0.276,0.906) (0.313,0.889) (0.348,0.870) (0.383,0.849) (0.417,0.827) (0.450,0.803) (0.482,0.778) (0.512,0.751) (0.542,0.724) (0.571,0.695) (0.598,0.665) (0.625,0.635) (0.651,0.603) (0.676,0.571) (0.700,0.539) (0.723,0.505) (0.746,0.472) (0.767,0.438) (0.789,0.403) (0.809,0.368) (0.829,0.333) (0.849,0.297) (0.868,0.262) (0.887,0.226) (0.906,0.190) (0.924,0.154) (0.942,0.117)};
\addplot[lineaE, flechaen=0.285, flechaen=0.715] coordinates {(-0.914,0.098) (-0.887,0.128) (-0.859,0.158) (-0.832,0.188) (-0.805,0.218) (-0.777,0.248) (-0.749,0.277) (-0.720,0.305) (-0.691,0.334) (-0.662,0.362) (-0.632,0.389) (-0.602,0.416) (-0.570,0.442) (-0.539,0.467) (-0.506,0.491) (-0.473,0.514) (-0.439,0.536) (-0.404,0.557) (-0.369,0.576) (-0.333,0.594) (-0.296,0.611) (-0.258,0.626) (-0.219,0.638) (-0.180,0.649) (-0.141,0.658) (-0.101,0.665) (-0.061,0.669) (-0.020,0.672) (0.020,0.672) (0.061,0.669) (0.101,0.665) (0.141,0.658) (0.180,0.649) (0.219,0.638) (0.258,0.626) (0.296,0.611) (0.333,0.594) (0.369,0.576) (0.404,0.557) (0.439,0.536) (0.473,0.514) (0.506,0.491) (0.539,0.467) (0.570,0.442) (0.602,0.416) (0.632,0.389) (0.662,0.362) (0.691,0.334) (0.720,0.305) (0.749,0.277) (0.777,0.248) (0.805,0.218) (0.832,0.188) (0.859,0.158) (0.887,0.128) (0.914,0.098)};
\addplot[lineaE, flechaen=0.255, flechaen=0.745] coordinates {(-0.882,0.063) (-0.847,0.082) (-0.812,0.101) (-0.776,0.119) (-0.741,0.138) (-0.705,0.156) (-0.669,0.174) (-0.633,0.191) (-0.597,0.208) (-0.560,0.225) (-0.523,0.241) (-0.486,0.257) (-0.449,0.271) (-0.412,0.285) (-0.374,0.298) (-0.335,0.311) (-0.297,0.322) (-0.258,0.332) (-0.219,0.340) (-0.180,0.348) (-0.140,0.354) (-0.100,0.358) (-0.060,0.361) (-0.020,0.363) (0.020,0.363) (0.060,0.361) (0.100,0.358) (0.140,0.354) (0.180,0.348) (0.219,0.340) (0.258,0.332) (0.297,0.322) (0.335,0.311) (0.374,0.298) (0.412,0.285) (0.449,0.271) (0.486,0.257) (0.523,0.241) (0.560,0.225) (0.597,0.208) (0.633,0.191) (0.669,0.174) (0.705,0.156) (0.741,0.138) (0.776,0.119) (0.812,0.101) (0.847,0.082) (0.882,0.063)};
\addplot[lineaE, flechaen=0.419] coordinates {(2.600,-1.027) (2.568,-1.002) (2.537,-0.977) (2.505,-0.952) (2.473,-0.928) (2.441,-0.903) (2.409,-0.879) (2.376,-0.855) (2.344,-0.831) (2.311,-0.807) (2.279,-0.784) (2.246,-0.760) (2.213,-0.737) (2.180,-0.714) (2.147,-0.691) (2.114,-0.668) (2.081,-0.645) (2.048,-0.622) (2.015,-0.600) (1.981,-0.578) (1.947,-0.555) (1.914,-0.533) (1.880,-0.511) (1.846,-0.490) (1.812,-0.468) (1.778,-0.447) (1.744,-0.425) (1.710,-0.404) (1.675,-0.383) (1.641,-0.362) (1.607,-0.341) (1.572,-0.321) (1.537,-0.300) (1.503,-0.280) (1.468,-0.259) (1.433,-0.239) (1.398,-0.219) (1.363,-0.199) (1.328,-0.180) (1.293,-0.160) (1.258,-0.140) (1.222,-0.121) (1.187,-0.102) (1.152,-0.082) (1.117,-0.063)};
\addplot[lineaE, flechaen=0.581] coordinates {(-1.117,-0.063) (-1.152,-0.082) (-1.187,-0.102) (-1.222,-0.121) (-1.258,-0.140) (-1.293,-0.160) (-1.328,-0.180) (-1.363,-0.199) (-1.398,-0.219) (-1.433,-0.239) (-1.468,-0.259) (-1.503,-0.280) (-1.537,-0.300) (-1.572,-0.321) (-1.607,-0.341) (-1.641,-0.362) (-1.675,-0.383) (-1.710,-0.404) (-1.744,-0.425) (-1.778,-0.447) (-1.812,-0.468) (-1.846,-0.490) (-1.880,-0.511) (-1.914,-0.533) (-1.947,-0.555) (-1.981,-0.578) (-2.015,-0.600) (-2.048,-0.622) (-2.081,-0.645) (-2.114,-0.668) (-2.147,-0.691) (-2.180,-0.714) (-2.213,-0.737) (-2.246,-0.760) (-2.279,-0.784) (-2.311,-0.807) (-2.344,-0.831) (-2.376,-0.855) (-2.409,-0.879) (-2.441,-0.903) (-2.473,-0.928) (-2.505,-0.952) (-2.537,-0.977) (-2.568,-1.002) (-2.600,-1.027)};
\addplot[lineaE, flechaen=0.159] coordinates {(2.112,-1.800) (2.099,-1.761) (2.086,-1.723) (2.072,-1.684) (2.058,-1.646) (2.043,-1.608) (2.028,-1.570) (2.013,-1.532) (1.997,-1.495) (1.981,-1.457) (1.964,-1.420) (1.948,-1.383) (1.930,-1.346) (1.913,-1.309) (1.895,-1.272) (1.877,-1.235) (1.858,-1.199) (1.840,-1.163) (1.820,-1.127) (1.801,-1.091) (1.781,-1.055) (1.761,-1.020) (1.741,-0.984) (1.720,-0.949) (1.699,-0.914) (1.677,-0.880) (1.656,-0.845) (1.634,-0.811) (1.611,-0.777) (1.589,-0.743) (1.566,-0.709) (1.543,-0.675) (1.519,-0.642) (1.496,-0.608) (1.472,-0.575) (1.447,-0.542) (1.423,-0.510) (1.398,-0.477) (1.373,-0.445) (1.348,-0.413) (1.323,-0.381) (1.297,-0.349) (1.272,-0.317) (1.246,-0.286) (1.220,-0.254) (1.193,-0.223) (1.167,-0.192) (1.140,-0.161) (1.114,-0.130) (1.087,-0.099)};
\addplot[lineaE, flechaen=0.841] coordinates {(-1.087,-0.099) (-1.114,-0.130) (-1.140,-0.161) (-1.167,-0.192) (-1.193,-0.223) (-1.220,-0.254) (-1.246,-0.286) (-1.272,-0.317) (-1.297,-0.349) (-1.323,-0.381) (-1.348,-0.413) (-1.373,-0.445) (-1.398,-0.477) (-1.423,-0.510) (-1.447,-0.542) (-1.472,-0.575) (-1.496,-0.608) (-1.519,-0.642) (-1.543,-0.675) (-1.566,-0.709) (-1.589,-0.743) (-1.611,-0.777) (-1.634,-0.811) (-1.656,-0.845) (-1.677,-0.880) (-1.699,-0.914) (-1.720,-0.949) (-1.741,-0.984) (-1.761,-1.020) (-1.781,-1.055) (-1.801,-1.091) (-1.820,-1.127) (-1.840,-1.163) (-1.858,-1.199) (-1.877,-1.235) (-1.895,-1.272) (-1.913,-1.309) (-1.930,-1.346) (-1.948,-1.383) (-1.964,-1.420) (-1.981,-1.457) (-1.997,-1.495) (-2.013,-1.532) (-2.028,-1.570) (-2.043,-1.608) (-2.058,-1.646) (-2.072,-1.684) (-2.086,-1.723) (-2.099,-1.761) (-2.112,-1.800)};
\addplot[lineaE, flechaen=0.500] coordinates {(1.459,-1.800) (1.461,-1.759) (1.463,-1.719) (1.464,-1.678) (1.465,-1.637) (1.464,-1.597) (1.463,-1.556) (1.462,-1.515) (1.460,-1.474) (1.457,-1.434) (1.454,-1.393) (1.450,-1.353) (1.445,-1.312) (1.440,-1.272) (1.434,-1.232) (1.427,-1.191) (1.420,-1.151) (1.413,-1.111) (1.405,-1.071) (1.396,-1.032) (1.387,-0.992) (1.377,-0.952) (1.367,-0.913) (1.356,-0.874) (1.345,-0.835) (1.333,-0.796) (1.321,-0.757) (1.309,-0.718) (1.296,-0.679) (1.282,-0.641) (1.268,-0.603) (1.254,-0.565) (1.239,-0.527) (1.224,-0.489) (1.209,-0.451) (1.193,-0.414) (1.177,-0.376) (1.160,-0.339) (1.144,-0.302) (1.127,-0.265) (1.110,-0.228) (1.092,-0.191) (1.075,-0.154) (1.057,-0.117)};
\addplot[lineaE, flechaen=0.500] coordinates {(-1.057,-0.117) (-1.075,-0.154) (-1.092,-0.191) (-1.110,-0.228) (-1.127,-0.265) (-1.144,-0.302) (-1.160,-0.339) (-1.177,-0.376) (-1.193,-0.414) (-1.209,-0.451) (-1.224,-0.489) (-1.239,-0.527) (-1.254,-0.565) (-1.268,-0.603) (-1.282,-0.641) (-1.296,-0.679) (-1.309,-0.718) (-1.321,-0.757) (-1.333,-0.796) (-1.345,-0.835) (-1.356,-0.874) (-1.367,-0.913) (-1.377,-0.952) (-1.387,-0.992) (-1.396,-1.032) (-1.405,-1.071) (-1.413,-1.111) (-1.420,-1.151) (-1.427,-1.191) (-1.434,-1.232) (-1.440,-1.272) (-1.445,-1.312) (-1.450,-1.353) (-1.454,-1.393) (-1.457,-1.434) (-1.460,-1.474) (-1.462,-1.515) (-1.463,-1.556) (-1.464,-1.597) (-1.465,-1.637) (-1.464,-1.678) (-1.463,-1.719) (-1.461,-1.759) (-1.459,-1.800)};
\addplot[lineaE, flechaen=0.500] coordinates {(0.904,-1.800) (0.924,-1.765) (0.943,-1.728) (0.961,-1.692) (0.979,-1.655) (0.995,-1.617) (1.010,-1.579) (1.024,-1.541) (1.037,-1.502) (1.050,-1.463) (1.061,-1.424) (1.071,-1.385) (1.081,-1.345) (1.089,-1.305) (1.097,-1.265) (1.103,-1.225) (1.109,-1.184) (1.114,-1.144) (1.119,-1.103) (1.122,-1.062) (1.125,-1.022) (1.127,-0.981) (1.128,-0.940) (1.128,-0.899) (1.128,-0.858) (1.127,-0.817) (1.125,-0.777) (1.123,-0.736) (1.120,-0.695) (1.117,-0.654) (1.112,-0.614) (1.108,-0.573) (1.103,-0.533) (1.097,-0.492) (1.091,-0.452) (1.084,-0.411) (1.078,-0.371) (1.070,-0.331) (1.063,-0.291) (1.055,-0.251) (1.046,-0.211) (1.038,-0.171) (1.029,-0.131)};
\addplot[lineaE, flechaen=0.500] coordinates {(-1.029,-0.131) (-1.038,-0.171) (-1.046,-0.211) (-1.055,-0.251) (-1.063,-0.291) (-1.070,-0.331) (-1.078,-0.371) (-1.084,-0.411) (-1.091,-0.452) (-1.097,-0.492) (-1.103,-0.533) (-1.108,-0.573) (-1.112,-0.614) (-1.117,-0.654) (-1.120,-0.695) (-1.123,-0.736) (-1.125,-0.777) (-1.127,-0.817) (-1.128,-0.858) (-1.128,-0.899) (-1.128,-0.940) (-1.127,-0.981) (-1.125,-1.022) (-1.122,-1.062) (-1.119,-1.103) (-1.114,-1.144) (-1.109,-1.184) (-1.103,-1.225) (-1.097,-1.265) (-1.089,-1.305) (-1.081,-1.345) (-1.071,-1.385) (-1.061,-1.424) (-1.050,-1.463) (-1.037,-1.502) (-1.024,-1.541) (-1.010,-1.579) (-0.995,-1.617) (-0.979,-1.655) (-0.961,-1.692) (-0.943,-1.728) (-0.924,-1.765) (-0.904,-1.800)};
\addplot[lineaE, flechaen=0.386, flechaen=0.614] coordinates {(-1.000,-0.131) (-0.999,-0.171) (-0.999,-0.211) (-0.998,-0.251) (-0.997,-0.291) (-0.996,-0.331) (-0.994,-0.371) (-0.991,-0.411) (-0.989,-0.451) (-0.986,-0.491) (-0.982,-0.531) (-0.978,-0.571) (-0.973,-0.611) (-0.967,-0.651) (-0.961,-0.690) (-0.954,-0.730) (-0.946,-0.769) (-0.937,-0.808) (-0.928,-0.847) (-0.918,-0.886) (-0.907,-0.925) (-0.895,-0.963) (-0.882,-1.001) (-0.868,-1.038) (-0.853,-1.076) (-0.837,-1.113) (-0.821,-1.149) (-0.803,-1.185) (-0.784,-1.220) (-0.764,-1.255) (-0.742,-1.289) (-0.720,-1.322) (-0.697,-1.355) (-0.672,-1.386) (-0.646,-1.417) (-0.620,-1.447) (-0.592,-1.476) (-0.563,-1.503) (-0.532,-1.530) (-0.501,-1.555) (-0.469,-1.579) (-0.436,-1.601) (-0.401,-1.622) (-0.366,-1.641) (-0.330,-1.659) (-0.293,-1.675) (-0.256,-1.689) (-0.218,-1.701) (-0.179,-1.711) (-0.140,-1.719) (-0.100,-1.726) (-0.060,-1.730) (-0.020,-1.732) (0.020,-1.732) (0.060,-1.730) (0.100,-1.726) (0.140,-1.719) (0.179,-1.711) (0.218,-1.701) (0.256,-1.689) (0.293,-1.675) (0.330,-1.659) (0.366,-1.641) (0.401,-1.622) (0.436,-1.601) (0.469,-1.579) (0.501,-1.555) (0.532,-1.530) (0.563,-1.503) (0.592,-1.476) (0.620,-1.447) (0.646,-1.417) (0.672,-1.386) (0.697,-1.355) (0.720,-1.322) (0.742,-1.289) (0.764,-1.255) (0.784,-1.220) (0.803,-1.185) (0.821,-1.149) (0.837,-1.113) (0.853,-1.076) (0.868,-1.038) (0.882,-1.001) (0.895,-0.963) (0.907,-0.925) (0.918,-0.886) (0.928,-0.847) (0.937,-0.808) (0.946,-0.769) (0.954,-0.730) (0.961,-0.690) (0.967,-0.651) (0.973,-0.611) (0.978,-0.571) (0.982,-0.531) (0.986,-0.491) (0.989,-0.451) (0.991,-0.411) (0.994,-0.371) (0.996,-0.331) (0.997,-0.291) (0.998,-0.251) (0.999,-0.211) (0.999,-0.171) (1.000,-0.131)};
\addplot[lineaE, flechaen=0.353, flechaen=0.647] coordinates {(-0.971,-0.128) (-0.962,-0.167) (-0.952,-0.207) (-0.943,-0.246) (-0.933,-0.285) (-0.922,-0.323) (-0.911,-0.362) (-0.900,-0.401) (-0.889,-0.439) (-0.876,-0.477) (-0.863,-0.516) (-0.850,-0.553) (-0.836,-0.591) (-0.821,-0.628) (-0.805,-0.665) (-0.788,-0.702) (-0.770,-0.738) (-0.752,-0.774) (-0.733,-0.809) (-0.712,-0.843) (-0.691,-0.877) (-0.668,-0.911) (-0.645,-0.943) (-0.620,-0.975) (-0.594,-1.006) (-0.567,-1.036) (-0.539,-1.064) (-0.510,-1.092) (-0.479,-1.118) (-0.448,-1.143) (-0.415,-1.166) (-0.381,-1.188) (-0.346,-1.208) (-0.311,-1.227) (-0.274,-1.243) (-0.236,-1.258) (-0.198,-1.270) (-0.159,-1.280) (-0.120,-1.289) (-0.080,-1.294) (-0.040,-1.298) (-0.000,-1.299) (0.040,-1.298) (0.080,-1.294) (0.120,-1.289) (0.159,-1.280) (0.198,-1.270) (0.236,-1.258) (0.274,-1.243) (0.311,-1.227) (0.346,-1.208) (0.381,-1.188) (0.415,-1.166) (0.448,-1.143) (0.479,-1.118) (0.510,-1.092) (0.539,-1.064) (0.567,-1.036) (0.594,-1.006) (0.620,-0.975) (0.645,-0.943) (0.668,-0.911) (0.691,-0.877) (0.712,-0.843) (0.733,-0.809) (0.752,-0.774) (0.770,-0.738) (0.788,-0.702) (0.805,-0.665) (0.821,-0.628) (0.836,-0.591) (0.850,-0.553) (0.863,-0.516) (0.876,-0.477) (0.889,-0.439) (0.900,-0.401) (0.911,-0.362) (0.922,-0.323) (0.933,-0.285) (0.943,-0.246) (0.952,-0.207) (0.962,-0.167) (0.971,-0.128)};
\addplot[lineaE, flechaen=0.318, flechaen=0.682] coordinates {(-0.942,-0.117) (-0.924,-0.154) (-0.906,-0.190) (-0.887,-0.226) (-0.868,-0.262) (-0.849,-0.297) (-0.829,-0.333) (-0.809,-0.368) (-0.789,-0.403) (-0.767,-0.438) (-0.746,-0.472) (-0.723,-0.505) (-0.700,-0.539) (-0.676,-0.571) (-0.651,-0.603) (-0.625,-0.635) (-0.598,-0.665) (-0.571,-0.695) (-0.542,-0.724) (-0.512,-0.751) (-0.482,-0.778) (-0.450,-0.803) (-0.417,-0.827) (-0.383,-0.849) (-0.348,-0.870) (-0.313,-0.889) (-0.276,-0.906) (-0.238,-0.921) (-0.200,-0.934) (-0.161,-0.944) (-0.121,-0.953) (-0.081,-0.959) (-0.041,-0.963) (-0.000,-0.964) (0.041,-0.963) (0.081,-0.959) (0.121,-0.953) (0.161,-0.944) (0.200,-0.934) (0.238,-0.921) (0.276,-0.906) (0.313,-0.889) (0.348,-0.870) (0.383,-0.849) (0.417,-0.827) (0.450,-0.803) (0.482,-0.778) (0.512,-0.751) (0.542,-0.724) (0.571,-0.695) (0.598,-0.665) (0.625,-0.635) (0.651,-0.603) (0.676,-0.571) (0.700,-0.539) (0.723,-0.505) (0.746,-0.472) (0.767,-0.438) (0.789,-0.403) (0.809,-0.368) (0.829,-0.333) (0.849,-0.297) (0.868,-0.262) (0.887,-0.226) (0.906,-0.190) (0.924,-0.154) (0.942,-0.117)};
\addplot[lineaE, flechaen=0.285, flechaen=0.715] coordinates {(-0.914,-0.098) (-0.887,-0.128) (-0.859,-0.158) (-0.832,-0.188) (-0.805,-0.218) (-0.777,-0.248) (-0.749,-0.277) (-0.720,-0.305) (-0.691,-0.334) (-0.662,-0.362) (-0.632,-0.389) (-0.602,-0.416) (-0.570,-0.442) (-0.539,-0.467) (-0.506,-0.491) (-0.473,-0.514) (-0.439,-0.536) (-0.404,-0.557) (-0.369,-0.576) (-0.333,-0.594) (-0.296,-0.611) (-0.258,-0.626) (-0.219,-0.638) (-0.180,-0.649) (-0.141,-0.658) (-0.101,-0.665) (-0.061,-0.669) (-0.020,-0.672) (0.020,-0.672) (0.061,-0.669) (0.101,-0.665) (0.141,-0.658) (0.180,-0.649) (0.219,-0.638) (0.258,-0.626) (0.296,-0.611) (0.333,-0.594) (0.369,-0.576) (0.404,-0.557) (0.439,-0.536) (0.473,-0.514) (0.506,-0.491) (0.539,-0.467) (0.570,-0.442) (0.602,-0.416) (0.632,-0.389) (0.662,-0.362) (0.691,-0.334) (0.720,-0.305) (0.749,-0.277) (0.777,-0.248) (0.805,-0.218) (0.832,-0.188) (0.859,-0.158) (0.887,-0.128) (0.914,-0.098)};
\addplot[lineaE, flechaen=0.255, flechaen=0.745] coordinates {(-0.882,-0.063) (-0.847,-0.082) (-0.812,-0.101) (-0.776,-0.119) (-0.741,-0.138) (-0.705,-0.156) (-0.669,-0.174) (-0.633,-0.191) (-0.597,-0.208) (-0.560,-0.225) (-0.523,-0.241) (-0.486,-0.257) (-0.449,-0.271) (-0.412,-0.285) (-0.374,-0.298) (-0.335,-0.311) (-0.297,-0.322) (-0.258,-0.332) (-0.219,-0.340) (-0.180,-0.348) (-0.140,-0.354) (-0.100,-0.358) (-0.060,-0.361) (-0.020,-0.363) (0.020,-0.363) (0.060,-0.361) (0.100,-0.358) (0.140,-0.354) (0.180,-0.348) (0.219,-0.340) (0.258,-0.332) (0.297,-0.322) (0.335,-0.311) (0.374,-0.298) (0.412,-0.285) (0.449,-0.271) (0.486,-0.257) (0.523,-0.241) (0.560,-0.225) (0.597,-0.208) (0.633,-0.191) (0.669,-0.174) (0.705,-0.156) (0.741,-0.138) (0.776,-0.119) (0.812,-0.101) (0.847,-0.082) (0.882,-0.063)};
\addplot[lineaE, flechaen=0.241, flechaen=0.759] coordinates {(-0.870,0.000) (0.870,0.000)};
\addplot[lineaE, flechaen=0.592] coordinates {(-1.130,0.000) (-2.600,0.000)};
\addplot[lineaE, flechaen=0.408] coordinates {(2.600,0.000) (1.130,0.000)};

  \fill[cargapos] (axis cs:-1,0) circle[radius=0.075];
  \fill[carganeg] (axis cs:1,0) circle[radius=0.075];
\end{axis}
\end{tikzpicture}\\[-0.1em]
{\scriptsize\color{tenue}\itshape Densidad de energía y líneas de campo\\ de dos cargas opuestas.}
\end{minipage}

\vspace{1.4em}
\noindent{\small{\color{acentoA}\espaciado{Palabras clave.}}\enspace energía electrostática $\cdot$ energía de interacción
$\cdot$ energía propia $\cdot$ densidad de energía $\cdot$ coeficientes de potencial $\cdot$ condensador $\cdot$ teorema de
Earnshaw $\cdot$ trabajos virtuales $\cdot$ fuerza a carga constante $\cdot$ fuerza a potencial constante}

\vspace{0.6em}
\noindent{\small{\color{acentoA}\espaciado{Clasificación PACS.}}\enspace 41.20.Cv (electrostática; ecuaciones de Poisson
y Laplace, problemas de contorno) $\cdot$ 84.32.Tt (condensadores)}

\vfill
\begin{center}
\small
$\displaystyle U_e=\frac12\int_V\rho\,\phi\,\dd V=\frac{\eo}{2}\int_{\mathbb R^3}E^2\,\dd V
\qquad U_e=\frac12\sum_{i,j}p_{ij}Q_iQ_j
\qquad F_\xi=-\frac{\partial U_e}{\partial\xi}\bigg|_{Q}=+\frac{\partial U_e}{\partial\xi}\bigg|_{\phi}$
\end{center}
\vspace{1em}
\newpage

% ------------------------------------------------------------- Índice
\thispagestyle{empty}
\tableofcontents
\thispagestyle{empty}
\newpage
\setcounter{page}{1}

% =============================================================================
\section{Introducción: por qué la energía}\label{sec:introduccion}
% =============================================================================

Hasta ahora hemos calculado campos y potenciales; en el Tema~3, además, la presión electrostática sobre la
superficie de un conductor~\cite{T3}. Calcular fuerzas entre cuerpos extensos sumando, o más bien integrando,
la fuerza de Coulomb sobre cada elemento de carga es casi siempre impracticable: las cargas de un conductor
no son un dato, sino que se redistribuyen al mover los cuerpos. En este tema introducimos la \emph{energía
electrostática}, que nos permitirá calcular esas fuerzas de un modo mucho más sencillo, derivando una sola
función escalar respecto a la coordenada que describe el movimiento.

El razonamiento tiene dos etapas. Primero (\cref{sec:puntuales,sec:continuas,sec:campo,sec:conductores}) definimos
la energía como el trabajo necesario para formar el sistema y la expresamos de tres maneras: en función de las
cargas y los potenciales, en función del campo y, para conductores, en función de sus cargas y potenciales
totales. Después (\cref{sec:fuerzas}) relacionamos las fuerzas con las variaciones de esa energía mediante el
principio de los trabajos virtuales. Esta primera parte se limita al vacío y a los conductores; la energía en
presencia de dieléctricos se tratará en la segunda.

Seguimos el orden de los apuntes manuscritos de la asignatura y nos apoyamos en los apuntes de teoría y en sus
figuras~\cite{DFA-teoria,DFA-figuras}, que hemos rehecho; para ampliar pueden consultarse los textos
clásicos~\cite{Griffiths,Wangsness,Purcell,Feynman,Jackson}.

\subsection{Notación y convenios}\label{sec:notacion}

\begin{notacion}\label{not:tema5}
Usamos la notación de los temas anteriores~\cite{T0,T1,T2,T3,T4} y, además:
\begin{itemize}
  \item $U_e$ es la energía electrostática del sistema y $u_e$ su densidad volúmica ($\si{J/m^3}$); $\Uint$, la energía de
        interacción entre partes del sistema.
  \item $\Wext$ es el trabajo realizado por un agente externo, $\Wmec=\int\vFe\cdot\dl$ el trabajo realizado por la fuerza
        electrostática $\vFe$ que actúa sobre un cuerpo, y $\Wbat$ el trabajo aportado por las baterías.
  \item En un sistema de cargas puntuales, $\vr_{ij}=\vr_i-\vr_j$ y $r_{ij}=\abs{\vr_{ij}}$; $\phi_i$ denota el potencial en $\vr_i$
        creado por todas las cargas \emph{salvo} la $i$-ésima.
  \item En un sistema de $N$ conductores, $Q_i$ y $\phi_i$ son la carga y el potencial del conductor $i$, y $p_{ij}$ y $c_{ij}$ los
        coeficientes de potencial y de capacidad e influencia del Tema~3~\cite{T3}.
  \item $\xi$ es una coordenada generalizada que describe la posición de un cuerpo (una distancia o un ángulo $\theta$), y
        $\eval{\cdot}{Q}$, $\eval{\cdot}{\phi}$ indican derivadas a cargas o a potenciales constantes.
  \item Las cargas están localizadas y $\phi(\infty)=0$; $k\equiv1/(4\pi\eo)$.
\end{itemize}
\end{notacion}

% =============================================================================
\section{Energía de un sistema de cargas puntuales}\label{sec:puntuales}
% =============================================================================

\subsection{Trabajo para trasladar una carga}

Recordemos la interpretación del potencial del Tema~1~\cite{T1}. Para trasladar una carga $q$ en un campo
electrostático sin acelerarla, un agente externo debe aplicarle en cada instante una fuerza igual y opuesta a la del
campo, $\vF_{\mathrm{ext}}=-q\vE$; decimos que el proceso es \emph{reversible} o cuasiestático, porque la carga pasa por una
sucesión de estados de equilibrio y su energía cinética es en todo momento despreciable.

\begin{lema}[Trabajo reversible y potencial]\label{lem:trabajo}
El trabajo reversible que un agente externo realiza para llevar una carga $q$ desde $A$ hasta $B$ en un campo
electrostático es
\begin{equation}\label{eq:trabajo}
  \Wext(A\to B)=q\big[\phi(B)-\phi(A)\big] ,
\end{equation}
independiente del camino. En particular, traerla desde el infinito cuesta $q\,\phi(B)$.
\end{lema}

\begin{proof}
Con $\vF_{\mathrm{ext}}=-q\vE$ y $\vE=-\grad\phi$, a lo largo de cualquier curva $\Gamma$ de $A$ a $B$,
\[
  \Wext=\int_\Gamma\vF_{\mathrm{ext}}\cdot\dl=q\int_\Gamma\grad\phi\cdot\dl=q\big[\phi(B)-\phi(A)\big],
\]
por el teorema fundamental del cálculo para integrales de línea. La segunda afirmación se sigue de $\phi(\infty)=0$.
\end{proof}

\begin{definicion}[Energía electrostática]\label{def:energia}
Se llama \emph{energía electrostática} $U_e$ de un sistema de cargas al trabajo reversible que debe realizar un agente
externo para formarlo, trayendo sus cargas desde una situación inicial en la que están en reposo, en el infinito e
infinitamente alejadas entre sí. Ese trabajo queda almacenado en el sistema como energía potencial.
\end{definicion}

\subsection{La energía de interacción}

Formemos el sistema trayendo las cargas una a una (\cref{fig:cargas}\,(a)). Para no confundir quién crea el potencial con dónde se evalúa,
llamamos
\begin{equation}\label{eq:phi-de-i}
  \phi^{(i)}(\vr)\defeq\frac{q_i}{4\pi\eo\abs{\vr-\vr_i}}
\end{equation}
al potencial creado \emph{por la carga $i$} en un punto $\vr$; en particular, $\phi^{(i)}(\vr_j)=q_i/(4\pi\eo r_{ij})$.
\begin{itemize}
  \item Traer $q_1$ no cuesta nada: aún no hay ningún campo que vencer, $W_1=0$.
  \item Traer $q_2$ hasta $\vr_2$ exige vencer el campo de $q_1$: $W_2=q_2\,\phi^{(1)}(\vr_2)=\dfrac{q_2q_1}{4\pi\eo r_{21}}$.
  \item Traer $q_3$ exige vencer los campos de $q_1$ y $q_2$: $W_3=q_3\big[\phi^{(1)}(\vr_3)+\phi^{(2)}(\vr_3)\big]$.
  \item En general, la carga $j$-ésima debe vencer el campo de las $j-1$ ya colocadas:
        \[
          W_j=q_j\sum_{i=1}^{j-1}\phi^{(i)}(\vr_j)=\sum_{i=1}^{j-1}\frac{q_jq_i}{4\pi\eo r_{ji}} .
        \]
\end{itemize}

% Figura: (a) formación de un sistema de N cargas puntuales trayéndolas una a una
% desde el infinito; (b) ejemplo de las cuatro cargas en un cuadrado.
\begin{figure}[t]
\centering
\begin{subfigure}[b]{0.5\textwidth}\centering
\begin{tikzpicture}[scale=1.0]
  % ejes
  \coordinate (O) at (0,0);
  \draw[tenue, line width=0.45pt, -{Stealth[length=1.6mm]}] (O) -- (2.9,0) node[below, etiqt] {$y$};
  \draw[tenue, line width=0.45pt, -{Stealth[length=1.6mm]}] (O) -- (0,2.4) node[left, etiqt] {$z$};
  \draw[tenue, line width=0.45pt, -{Stealth[length=1.6mm]}] (O) -- (-1.3,-1.0) node[left, etiqt] {$x$};
  % cargas ya colocadas
  \foreach \p/\n in {{(1.3,1.7)}/1,{(2.3,0.9)}/2,{(0.9,-0.9)}/3} {
    \fill[cargapos] \p circle (0.075);
    \draw[campoP, line width=0.6pt, -{Stealth[length=1.6mm]}] (O) -- \p;
  }
  \node[etiq, anchor=south west] at (1.35,1.72) {$q_1$};
  \node[etiq, anchor=west] at (2.38,0.92) {$q_2$};
  \node[etiq, anchor=north west] at (0.95,-0.95) {$q_3$};
  \node[etiq, text=campoP] at (0.42,1.05) {$\vr_1$};
  \node[etiq, text=campoP] at (1.42,0.3) {$\vr_2$};
  \node[etiq, text=campoP] at (0.25,-0.62) {$\vr_3$};
  % carga N trayéndose desde el infinito
  \fill[cargapos] (-1.4,1.55) circle (0.075);
  \node[etiq, anchor=south] at (-1.4,1.62) {$q_N$};
  \draw[campoE, line width=0.6pt, dashed, -{Stealth[length=1.8mm]}] (-2.9,2.25) -- (-1.5,1.6);
  \node[etiqt, anchor=south] at (-2.55,2.3) {desde $\infty$};
  \draw[campoP, line width=0.6pt, -{Stealth[length=1.6mm]}] (O) -- (-1.33,1.48);
  \node[etiq, text=campoP] at (-0.95,0.62) {$\vr_N$};
  \node[etiqt] at (1.9,-0.25) {$\cdots$};
\end{tikzpicture}
\caption{Formación del sistema carga a carga}
\end{subfigure}\hfill
\begin{subfigure}[b]{0.44\textwidth}\centering
\begin{tikzpicture}[scale=1.25]
  \draw[tenue, line width=0.45pt, -{Stealth[length=1.6mm]}] (-0.3,0) -- (2.2,0) node[below, etiqt] {$x$};
  \draw[tenue, line width=0.45pt, -{Stealth[length=1.6mm]}] (0,-0.3) -- (0,2.0) node[left, etiqt] {$y$};
  \draw[equipot] (0,0) -- (1.5,0) -- (1.5,1.5) -- (0,1.5) -- cycle;
  \draw[acentoB, line width=0.55pt] (0,0) -- (1.5,1.5);
  \node[etiqf, text=acentoB, font=\footnotesize, rotate=45] at (0.62,0.82) {$\sqrt2\,a$};
  \foreach \p in {(0,0),(1.5,0),(1.5,1.5),(0,1.5)} \fill[cargapos] \p circle (0.07);
  \node[etiq, anchor=north east] at (-0.03,-0.03) {$q$};
  \node[etiq, anchor=north] at (1.5,-0.08) {$q$};
  \node[etiq, anchor=south west] at (1.53,1.53) {$q$};
  \node[etiq, anchor=east] at (-0.08,1.5) {$q$};
  \draw[cota] (0,-0.55) -- (1.5,-0.55) node[midway, below, etiq] {$a$};
\end{tikzpicture}
\caption{Cuatro cargas en un cuadrado}
\end{subfigure}
\caption{Energía de un sistema de cargas puntuales. (a) El sistema se forma trayendo las cargas una a una desde el
infinito, donde están en reposo e infinitamente alejadas entre sí; la carga $q_N$ debe vencer el campo de las $N-1$ ya
colocadas. (b) Cuatro cargas iguales en los vértices de un cuadrado de lado $a$ (\cref{ej:cuadrado}): cada carga
tiene dos vecinas a distancia $a$ y una a distancia $\sqrt2\,a$.}
\label{fig:cargas}
\end{figure}


\begin{proposicion}[Energía de un sistema de cargas puntuales]\label{prop:puntuales}
La energía electrostática de un sistema de $N$ cargas puntuales $q_i$ situadas en $\vr_i$ es
\begin{equation}\label{eq:puntuales}
  U_e=\sum_{j=2}^{N}\sum_{i=1}^{j-1}\frac{q_iq_j}{4\pi\eo r_{ij}}
     =\frac12\sum_{i=1}^{N}\sum_{\substack{j=1\\ j\neq i}}^{N}\frac{q_iq_j}{4\pi\eo r_{ij}}
     =\boxed{\;\frac12\sum_{i=1}^{N}q_i\,\phi_i\;},
\end{equation}
donde
\[
  \phi_i\defeq\sum_{j\neq i}\phi^{(j)}(\vr_i)=\sum_{j\neq i}\frac{q_j}{4\pi\eo r_{ij}}
\]
es el potencial \emph{en el punto} $\vr_i$ creado por todas las cargas salvo la propia. El resultado
no depende del orden en que se traigan las cargas.
\end{proposicion}

\begin{proof}
La primera igualdad es la suma $U_e=\sum_jW_j$ de los trabajos anteriores. Su sumando es simétrico en $i$ y $j$, de modo que
la suma sobre los pares con $i<j$ es la mitad de la suma sobre todos los pares ordenados con $i\neq j$, que cuenta cada par dos
veces; esto da la segunda igualdad, en la que ya no aparece ningún orden. Finalmente, agrupando los términos según la carga $q_i$ que recibe el potencial,
\[
  U_e=\frac12\sum_{i=1}^Nq_i\sum_{j\neq i}\frac{q_j}{4\pi\eo r_{ij}}=\frac12\sum_{i=1}^Nq_i\sum_{j\neq i}\phi^{(j)}(\vr_i)=\frac12\sum_{i=1}^Nq_i\,\phi_i . \qedhere
\]
\end{proof}

\begin{observacion}[Una energía de interacción]\label{obs:interaccion}
\begin{enumerate}
  \item El factor $\frac12$ no es un descuido: en $\sum_iq_i\phi_i$ cada pareja aparece dos veces, una como $q_i\,\phi^{(j)}(\vr_i)$ y otra como
        $q_j\,\phi^{(i)}(\vr_j)$, que valen lo mismo. Con él, $\phi_i$ incluye el potencial de \emph{todas} las demás cargas, no solo el de las que ya estaban
        presentes al traer la $i$-ésima.
  \item Se excluyen los términos $i=j$: una carga no ejerce fuerza sobre sí misma. La expresión \eqref{eq:puntuales} no incluye,
        por tanto, la energía necesaria para formar cada carga puntual; las cargas «ya existían» en el infinito. Por eso se llama
        \emph{energía de interacción}.
  \item Puede ser positiva o negativa. Para dos cargas, $U_e=q_1q_2/(4\pi\eo r_{12})$ es positiva si se repelen (hubo que empujarlas
        para acercarlas) y negativa si se atraen (el campo hizo el trabajo y el sistema cedió energía).
\end{enumerate}
\end{observacion}

\begin{ejemplo}[Cuatro cargas en un cuadrado]\label{ej:cuadrado}
Sean cuatro cargas $q$ en los vértices de un cuadrado de lado $a$ (\cref{fig:cargas}\,(b)).

\emph{Primer método: trayéndolas por orden.} La primera no cuesta nada; la segunda, colocada a distancia $a$ de la primera,
$q\cdot kq/a$; la tercera, en el vértice opuesto a la primera, está a distancia $a$ de la segunda y $\sqrt2\,a$ de la primera; la
cuarta está a distancia $a$ de dos de ellas y $\sqrt2\,a$ de la otra:
\[
  U_e=0+\frac{kq^2}{a}+kq^2\Big(\frac1a+\frac{1}{\sqrt2\,a}\Big)+kq^2\Big(\frac2a+\frac{1}{\sqrt2\,a}\Big)
     =\frac{q^2}{4\pi\eo a}\Big(4+\frac{2}{\sqrt2}\Big)=\frac{q^2}{4\pi\eo a}\big(4+\sqrt2\big).
\]
\emph{Segundo método: con la fórmula simétrica.} Por simetría, las cuatro cargas ven el mismo potencial, creado por dos vecinas a
distancia $a$ y una a distancia $\sqrt2\,a$, de modo que
\[
  U_e=\frac12\cdot4\,q\,\phi_1=2q\Big(\frac{2kq}{a}+\frac{kq}{\sqrt2\,a}\Big)=\frac{q^2}{4\pi\eo a}\big(4+\sqrt2\big).
\]
El segundo método es mucho más rápido cuando hay simetría: basta un solo potencial.
\end{ejemplo}

% =============================================================================
\section{Distribuciones continuas de carga}\label{sec:continuas}
% =============================================================================

Para una distribución continua, cada elemento de carga $\dd q=\rho\,\dd V$ desempeña el papel de una carga puntual, y la suma de
\eqref{eq:puntuales} se convierte en una integral.

\begin{proposicion}[Energía de una distribución continua]\label{prop:continua}
La energía electrostática de una distribución de carga de densidad volúmica $\rho$ en una región $V$, o superficial $\sigma$ sobre una
superficie $S$, es
\begin{equation}\label{eq:continua}
  \boxed{\;U_e=\frac12\int_V\rho(\vr)\,\phi(\vr)\,\dd V\;},\qquad U_e=\frac12\int_S\sigma(\vr)\,\phi(\vr)\,\dd S ,
\end{equation}
donde $\phi$ es el potencial creado por \emph{toda} la distribución. Si coexisten carga volúmica y superficial, se suman ambas
integrales, con el potencial de ambas.
\end{proposicion}

\begin{proof}[Justificación]
Construyamos la distribución por etapas, multiplicando todas las densidades por un parámetro $\lambda$ que crece de $0$ a $1$. En
una etapa intermedia la densidad es $\lambda\rho$ y, por linealidad, el potencial es $\lambda\phi$. Añadir la carga $\dd\lambda\,\rho\,\dd V$ en cada
punto, trayéndola desde el infinito, cuesta por el \cref{lem:trabajo}
\[
  \dd\Wext=\int_V(\dd\lambda\,\rho)\,(\lambda\phi)\,\dd V=\lambda\,\dd\lambda\int_V\rho\,\phi\,\dd V ,
\]
y el trabajo total es $\int_0^1\lambda\,\dd\lambda\int_V\rho\phi\,\dd V=\frac12\int_V\rho\phi\,\dd V$. El factor $\frac12$ refleja, otra vez, que cada elemento de carga
se trae contra un potencial que, en promedio a lo largo del proceso, vale la mitad del final. El caso superficial es idéntico con
$\rho\,\dd V\to\sigma\,\dd S$.
\end{proof}

\begin{observacion}[Energía total, no solo de interacción]\label{obs:total}
\begin{enumerate}
  \item Ahora se parte de una nube de carga infinitamente diluida ($\rho=0$) y se \emph{forma} la distribución. El potencial $\phi$ de
        \eqref{eq:continua} es el de todos los elementos, incluido el propio, y el resultado incluye la energía de formación: es la
        energía \emph{total}. Para una distribución continua, la contribución de cada elemento a su propio potencial es despreciable
        (un elemento $\dd V$ de tamaño $\delta$ crea en sí mismo un potencial $\sim\rho\delta^2$, que se anula con $\delta$), y por eso no hace falta
        excluir nada.
  \item Veremos en la \cref{sec:campo} que esta energía es siempre positiva o nula, a diferencia de la de interacción.
  \item La integral puede extenderse a todo el espacio, pues fuera de la distribución $\rho=0$ y no contribuye.
  \item Para una carga puntual o una distribución lineal, la energía de formación diverge: concentrar una carga finita en un punto, o
        con densidad finita en una línea, exige un trabajo infinito. Por eso, para cargas puntuales, se trabaja solo con la energía de
        interacción \eqref{eq:puntuales}.
\end{enumerate}
\end{observacion}

\begin{ejemplo}[Corteza esférica conductora]\label{ej:corteza}
Sea una corteza esférica de radio $a$ con carga $Q$ uniformemente repartida (por ejemplo, una esfera conductora). Su superficie es
equipotencial, con $\phi=Q/(4\pi\eo a)$, de modo que
\[
  U_e=\frac12\oiint_S\sigma\,\phi\,\dd S=\frac12\,\phi\oiint_S\sigma\,\dd S=\frac12\,\frac{Q}{4\pi\eo a}\,Q=\boxed{\;\frac{Q^2}{8\pi\eo a}\;}.
\]
Lo comprobaremos con el campo en el \cref{ej:corteza-campo}.
\end{ejemplo}

% =============================================================================
\section{Energía en función del campo}\label{sec:campo}
% =============================================================================

Las expresiones anteriores asocian la energía a las cargas y la localizan donde están ellas. Igual que en el Tema~1 pasamos de la ley
de Coulomb, que habla de cargas, a la ley de Gauss, que habla de campos, vamos a expresar la energía solo en función del campo.

\subsection{El teorema de la energía del campo}

\begin{lema}[Dos herramientas]\label{lem:herramientas}
Para un campo escalar $\phi$ y un campo vectorial $\vA$ de clase $\mathcal C^1$:
\begin{enumerate}
  \item identidad del producto: $\divg(\phi\vA)=\grad\phi\cdot\vA+\phi\,\divg\vA$;
  \item teorema de la divergencia de Gauss: para una región acotada $V$ de frontera $S=\partial V$ con normal exterior,
        $\displaystyle\int_V\divg\vA\,\dd V=\oiint_S\vA\cdot\dS$.
\end{enumerate}
\end{lema}

\begin{teorema}[Energía electrostática en función del campo]\label{teo:energia-campo}
Para una distribución de carga localizada en el vacío,
\begin{equation}\label{eq:energia-campo}
  \boxed{\;U_e=\frac12\int_V\rho\,\phi\,\dd V=\frac{\eo}{2}\int_{\mathbb R^3}E^2\,\dd V\; ,}
\end{equation}
donde la segunda integral se extiende a todo el espacio.
\end{teorema}

\begin{proof}
\paso{Paso 1: eliminar $\rho$} Por la ley de Gauss diferencial, $\rho=\eo\,\divg\vE$, de modo que
$U_e=\dfrac{\eo}{2}\displaystyle\int_V\phi\,\divg\vE\,\dd V$. Como $\rho=0$ fuera de la distribución, podemos tomar como $V$ una bola $B_R$ de radio $R$
tan grande como queramos (\cref{fig:volumen}).

\paso{Paso 2: integrar por partes} Por el \cref{lem:herramientas}\,(i), $\phi\,\divg\vE=\divg(\phi\vE)-\grad\phi\cdot\vE=\divg(\phi\vE)+E^2$, ya que
$\grad\phi=-\vE$. Por tanto,
\[
  U_e=\frac{\eo}{2}\int_{B_R}\divg(\phi\vE)\,\dd V+\frac{\eo}{2}\int_{B_R}E^2\,\dd V .
\]
\paso{Paso 3: el teorema de la divergencia} Aplicándolo a la primera integral, con $S_R=\partial B_R$,
\[
  U_e=\frac{\eo}{2}\oiint_{S_R}\phi\,\vE\cdot\dS+\frac{\eo}{2}\int_{B_R}E^2\,\dd V .
\]
Esta igualdad vale para todo $R$ suficientemente grande, aunque cada término por separado dependa de $R$.

\paso{Paso 4: límite $R\to\infty$} Lejos de una distribución localizada, $\phi=\mathcal O(1/R)$ y $E=\mathcal O(1/R^2)$, mientras que el área de $S_R$ es $4\pi R^2$:
\[
  \abs{\oiint_{S_R}\phi\,\vE\cdot\dS}\leq\max_{S_R}\abs\phi\,\max_{S_R}\abs{\vE}\,4\pi R^2=\mathcal O\Big(\frac1R\cdot\frac{1}{R^2}\cdot R^2\Big)=\mathcal O\Big(\frac1R\Big)\xrightarrow{R\to\infty}0 .
\]
(Si la distribución es neutra, $\phi$ y $E$ decaen aún más deprisa.) La segunda integral converge, pues, a la integral de $E^2$ a todo el
espacio, y se obtiene \eqref{eq:energia-campo}.
\end{proof}

% Figura: paso de la energía en función de las cargas a la energía en función del
% campo. Volumen V con la carga, esfera S_R de radio R -> infinito y órdenes de magnitud.
\begin{figure}[t]
\centering
\begin{tikzpicture}[scale=1.0]
  % esfera grande S_R
  \draw[gauss] (0,0) circle (2.35);
  \node[etiq, text=acentoB] at (52:2.62) {$S_R$};
  \draw[texto!60!fondo, line width=0.4pt, -{Stealth[length=1.5mm]}] (0,0) -- (-35:2.35)
     node[pos=0.64, sloped, above, etiq, font=\footnotesize] {$R\to\infty$};
  % distribución de carga
  \fill[dielec] plot[smooth cycle, tension=0.85] coordinates
     {(-1.05,0.25) (-0.55,0.95) (0.35,0.85) (1.0,0.35) (0.85,-0.45) (0.05,-0.75) (-0.8,-0.5)};
  \draw[texto!70!fondo, line width=0.5pt] plot[smooth cycle, tension=0.85] coordinates
     {(-1.05,0.25) (-0.55,0.95) (0.35,0.85) (1.0,0.35) (0.85,-0.45) (0.05,-0.75) (-0.8,-0.5)};
  \node[etiq] at (-0.35,0.15) {$\rho(\vr)$};
  \node[etiq] at (0.55,-0.25) {$V$};
  \fill[cargapos] (0.05,0.5) circle (0.04);
  \node[etiqt, anchor=west] at (0.1,0.5) {$\dd V$};
  % elemento de superficie y campo lejano
  \coordinate (Q) at (125:2.35);
  \draw[vecn] (Q) -- ($(Q)+(125:0.75)$) node[anchor=south east, etiq] {$\dS$};
  \draw[vecE] (Q) -- ($(Q)+(95:0.65)$) node[anchor=south, etiq, text=campoE] {$\vE$};
  % tabla de escalas
  \node[anchor=west, align=left, font=\small, text=texto] at (3.05,0.35)
    {$\phi\sim\dfrac{1}{R}$,\qquad $E\sim\dfrac{1}{R^{2}}$,\qquad $\dd S\sim R^{2}$\\[3mm]
     $\displaystyle\Big|\oiint_{S_R}\phi\,\vE\cdot\dS\Big|\sim\frac1R\cdot\frac{1}{R^2}\cdot R^2=\frac1R\ \xrightarrow{R\to\infty}\ 0$};
\end{tikzpicture}
\caption{De las cargas a los campos. La energía $\frac12\int_V\rho\phi\,\dd V$ puede extenderse a cualquier volumen que contenga
a $V$, pues fuera de él $\rho=0$. Tomando una esfera $S_R$ de radio $R\to\infty$, el término de superficie que aparece al aplicar el
teorema de la divergencia decae como $1/R$ (para una distribución con carga neta; aún más deprisa si es neutra) y solo
sobrevive la integral de $E^2$ a todo el espacio (\cref{teo:energia-campo}).}
\label{fig:volumen}
\end{figure}


\begin{definicion}[Densidad de energía electrostática]\label{def:densidad}
Se llama \emph{densidad de energía electrostática} a
\begin{equation}\label{eq:densidad}
  u_e(\vr)\defeq\frac{\eo}{2}\,E^2(\vr)\qquad[\si{J/m^3}],\qquad U_e=\int_{\mathbb R^3}u_e\,\dd V .
\end{equation}
\end{definicion}

\begin{corolario}[Positividad]\label{cor:positiva}
La energía electrostática total de una distribución continua es positiva o nula, $U_e\geq0$, y solo se anula si $\vE\equiv\vcero$, es decir, si
no hay carga.
\end{corolario}

\begin{proof}
El integrando de \eqref{eq:energia-campo} es no negativo, y la integral de una función continua no negativa solo se anula si la
función es idénticamente nula.
\end{proof}

\begin{observacion}[¿Dónde está la energía?]\label{obs:localizacion}
Las dos expresiones de \eqref{eq:energia-campo} dan el mismo número, pero cuentan historias distintas: la primera localiza la energía
donde hay carga; la segunda, en todo punto donde hay campo, aunque allí no haya carga (\cref{fig:corteza}). En electrostática ambas
interpretaciones son equivalentes. La segunda, sin embargo, es la que sobrevive fuera de la electrostática: una onda
electromagnética transporta energía por regiones vacías, y $u_e=\frac{\eo}{2}E^2$ sigue siendo válida para campos variables en el
tiempo~\cite{Griffiths,Jackson}.
\end{observacion}

\begin{ejemplo}[Corteza esférica, de nuevo]\label{ej:corteza-campo}
Para la corteza del \cref{ej:corteza}, el campo es nulo dentro y vale $\vE=\frac{Q}{4\pi\eo r^2}\uv r$ fuera. Con $\dd V=4\pi r^2\dd r$ (simetría
esférica),
\[
  U_e=\frac{\eo}{2}\int_a^\infty\Big(\frac{Q}{4\pi\eo r^2}\Big)^2 4\pi r^2\,\dd r=\frac{Q^2}{8\pi\eo}\int_a^\infty\frac{\dd r}{r^2}
     =\frac{Q^2}{8\pi\eo}\Big[-\frac1r\Big]_a^\infty=\frac{Q^2}{8\pi\eo a},
\]
como antes. La energía almacenada más allá de un radio $r$ es $Q^2/(8\pi\eo r)$: la mitad está fuera de $2a$ (\cref{fig:corteza}).
\end{ejemplo}

% Figura: energía de una corteza esférica conductora de radio a y carga Q:
% densidad de energía u_e(r) y fracción de la energía almacenada fuera de r.
\begin{figure}[t]
\centering
\begin{subfigure}[b]{0.3\textwidth}\centering
\begin{tikzpicture}[scale=1.0]
  \shade[ball color=conductorluz] (0,0) circle (1.0);
  \draw[texto, line width=0.5pt] (0,0) circle (1.0);
  \foreach \t in {0,30,...,330} \node[cpos, font=\scriptsize\bfseries] at (\t:1.2) {$+$};
  \draw[texto, line width=0.4pt, -{Stealth[length=1.5mm]}] (0,0) -- (35:1.0);
  \node[font=\small, text=black] at (35:0.45) [above left=-2pt] {$a$};
  \node[etiq] at (0,-1.6) {$Q$, \ $\phi=\dfrac{Q}{4\pi\eo a}$};
\end{tikzpicture}
\caption{Corteza cargada}
\end{subfigure}\hfill
\begin{subfigure}[b]{0.33\textwidth}\centering
\begin{tikzpicture}
\begin{axis}[width=5.0cm, height=4.3cm, axis lines=left, xmin=0, xmax=4, ymin=0, ymax=1.15,
   xlabel={$r/a$}, ylabel={$u_e/u_e(a^{+})$}, xtick={0,1,2,3,4}, ytick={0,0.5,1}]
  \addplot[curvaE, domain=1:4, samples=80] {1/x^4};
  \addplot[curvaE] coordinates {(0,0) (1,0)};
  \draw[equipot] (axis cs:1,0) -- (axis cs:1,1.15);
  \node[etiqt, anchor=west] at (axis cs:1.4,0.55) {$\propto r^{-4}$};
\end{axis}
\end{tikzpicture}
\caption{Densidad de energía}
\end{subfigure}\hfill
\begin{subfigure}[b]{0.33\textwidth}\centering
\begin{tikzpicture}
\begin{axis}[width=5.0cm, height=4.3cm, axis lines=left, xmin=0, xmax=6, ymin=0, ymax=1.1,
   xlabel={$r/a$}, ylabel={$U_e(>r)/U_e$}, xtick={0,1,2,4,6}, ytick={0,0.5,1}, grid=major]
  \addplot[curvaD, domain=1:6, samples=80] {1/x};
  \addplot[curvaD] coordinates {(0,1) (1,1)};
  \addplot[only marks, mark=*, mark size=1.4pt, color=acentoC] coordinates {(2,0.5)};
  \node[etiqf, font=\scriptsize, anchor=south west] at (axis cs:2.15,0.55) {$\tfrac12$ en $r>2a$};
\end{axis}
\end{tikzpicture}
\caption{Energía fuera de $r$}
\end{subfigure}
\caption{\Cref{ej:corteza}. (a) Corteza esférica de radio $a$ con carga $Q$, equipotencial. (b) La densidad de energía
$u_e=\frac{\eo}{2}E^2$ es nula dentro y decae como $r^{-4}$ fuera. (c) La energía almacenada fuera de una esfera de radio $r>a$
es $U_e(>r)=Q^2/(8\pi\eo r)$: la mitad de la energía total está más allá de $r=2a$, lejos de la carga. Es la imagen que da
la \cref{obs:localizacion}: la energía «vive» en el campo.}
\label{fig:corteza}
\end{figure}


\subsection{La energía no es aditiva}

El campo verifica el principio de superposición; la energía, que es cuadrática en el campo, no.

\begin{proposicion}[Término cruzado]\label{prop:no-aditiva}
Si $\vE=\vE_1+\vE_2$ es el campo de dos distribuciones de carga, con energías propias $U_1$ y $U_2$,
\begin{equation}\label{eq:no-aditiva}
  U_e=\frac{\eo}{2}\int\big(E_1^2+E_2^2+2\,\vE_1\cdot\vE_2\big)\dd V=U_1+U_2+\eo\int_{\mathbb R^3}\vE_1\cdot\vE_2\,\dd V\neq U_1+U_2 .
\end{equation}
El término cruzado es la energía de interacción entre ambas distribuciones,
\begin{equation}\label{eq:cruzado}
  \Uint=\eo\int_{\mathbb R^3}\vE_1\cdot\vE_2\,\dd V=\int\rho_2\,\phi_1\,\dd V=\int\rho_1\,\phi_2\,\dd V .
\end{equation}
\end{proposicion}

\begin{proof}
La primera igualdad es el desarrollo de $(\vE_1+\vE_2)^2$. Para la segunda, con $\vE_1=-\grad\phi_1$ y la misma integración por partes
del \cref{teo:energia-campo},
\[
  \eo\int_{B_R}\vE_1\cdot\vE_2\,\dd V=-\eo\int_{B_R}\grad\phi_1\cdot\vE_2\,\dd V=-\eo\oiint_{S_R}\phi_1\vE_2\cdot\dS+\eo\int_{B_R}\phi_1\,\divg\vE_2\,\dd V ,
\]
y al hacer $R\to\infty$ el término de superficie se anula mientras que $\eo\,\divg\vE_2=\rho_2$. Intercambiando los papeles de 1 y 2 se obtiene la
otra forma.
\end{proof}

\begin{ejemplo}[Dos cargas puntuales]\label{ej:termino-cruzado}
Para dos cargas puntuales $q_1$ y $q_2$ separadas una distancia $d$, $U_1$ y $U_2$ son las energías propias, infinitas, de cada carga; el
término cruzado, en cambio, es finito y vale $\int\rho_2\phi_1\,\dd V=q_2\,\phi_1(\vr_2)=\dfrac{q_1q_2}{4\pi\eo d}$: la energía de interacción de la
\cref{prop:puntuales}. Para cargas iguales, $\vE_1\cdot\vE_2$ es negativo dentro de la circunferencia que tiene por diámetro el segmento que
las une (allí los campos forman un ángulo obtuso) y positivo fuera; el balance es positivo (\cref{fig:superposicion}).
\end{ejemplo}

% Figura (simulación): la energía no es aditiva. Densidad de energía de dos cargas
% iguales y opuestas, y término cruzado eps0 E1.E2.
% Mapas: scripts/tema5.py -> img/t5-ue-*.png ; líneas: datos/t5-dos-cargas-*.tex
\begin{figure}[t]
\centering
\pgfplotsset{mapa/.style={marco, width=5.6cm, xmin=-2.6, xmax=2.6, ymin=-1.8, ymax=1.8,
   enlargelimits=false, axis on top}}
\begin{subfigure}[b]{0.325\textwidth}\centering
\begin{tikzpicture}
\begin{axis}[mapa]
  \addplot graphics[xmin=-2.6, xmax=2.6, ymin=-1.8, ymax=1.8] {t5-ue-iguales-\modotema};
  %% Generado por scripts/tema5.py -- no editar a mano
\addplot[lineaE, flechaen=0.604] coordinates {(-1.111,0.069) (-1.146,0.090) (-1.180,0.111) (-1.214,0.132) (-1.249,0.153) (-1.283,0.174) (-1.318,0.194) (-1.353,0.215) (-1.387,0.235) (-1.422,0.256) (-1.457,0.276) (-1.492,0.296) (-1.527,0.316) (-1.562,0.336) (-1.597,0.355) (-1.632,0.375) (-1.668,0.395) (-1.703,0.414) (-1.738,0.433) (-1.774,0.453) (-1.809,0.472) (-1.845,0.491) (-1.880,0.510) (-1.916,0.529) (-1.951,0.547) (-1.987,0.566) (-2.023,0.585) (-2.059,0.603) (-2.095,0.622) (-2.130,0.640) (-2.166,0.658) (-2.202,0.676) (-2.238,0.694) (-2.274,0.712) (-2.310,0.730) (-2.346,0.748) (-2.383,0.766) (-2.419,0.784) (-2.455,0.802) (-2.491,0.819) (-2.527,0.837) (-2.564,0.854) (-2.600,0.872)};
\addplot[lineaE, flechaen=0.661] coordinates {(-1.073,0.113) (-1.095,0.146) (-1.117,0.180) (-1.139,0.213) (-1.161,0.246) (-1.184,0.280) (-1.207,0.312) (-1.229,0.345) (-1.253,0.378) (-1.276,0.411) (-1.299,0.443) (-1.323,0.475) (-1.347,0.508) (-1.371,0.539) (-1.395,0.571) (-1.420,0.603) (-1.444,0.634) (-1.469,0.666) (-1.494,0.697) (-1.520,0.728) (-1.545,0.759) (-1.571,0.790) (-1.597,0.820) (-1.623,0.851) (-1.649,0.881) (-1.675,0.911) (-1.701,0.941) (-1.728,0.971) (-1.755,1.001) (-1.781,1.031) (-1.808,1.061) (-1.835,1.090) (-1.863,1.119) (-1.890,1.149) (-1.917,1.178) (-1.945,1.207) (-1.972,1.236) (-2.000,1.265) (-2.028,1.294) (-2.056,1.322) (-2.084,1.351) (-2.112,1.380) (-2.140,1.408) (-2.168,1.437) (-2.196,1.465) (-2.225,1.493) (-2.253,1.521) (-2.282,1.549) (-2.310,1.578) (-2.339,1.606) (-2.368,1.633) (-2.396,1.661) (-2.425,1.689) (-2.454,1.717) (-2.483,1.745) (-2.512,1.772) (-2.541,1.800)};
\addplot[lineaE, flechaen=0.500] coordinates {(-1.031,0.126) (-1.040,0.166) (-1.050,0.206) (-1.061,0.245) (-1.071,0.285) (-1.082,0.324) (-1.093,0.364) (-1.105,0.403) (-1.116,0.442) (-1.128,0.481) (-1.141,0.520) (-1.154,0.559) (-1.167,0.597) (-1.181,0.636) (-1.195,0.674) (-1.209,0.713) (-1.224,0.751) (-1.239,0.789) (-1.255,0.827) (-1.271,0.864) (-1.287,0.902) (-1.303,0.939) (-1.320,0.976) (-1.337,1.014) (-1.355,1.050) (-1.373,1.087) (-1.391,1.124) (-1.409,1.161) (-1.428,1.197) (-1.446,1.233) (-1.465,1.269) (-1.485,1.306) (-1.504,1.341) (-1.524,1.377) (-1.544,1.413) (-1.564,1.449) (-1.584,1.484) (-1.604,1.520) (-1.625,1.555) (-1.646,1.590) (-1.667,1.625) (-1.688,1.660) (-1.709,1.695) (-1.730,1.730) (-1.751,1.765) (-1.773,1.800)};
\addplot[lineaE, flechaen=0.500] coordinates {(-0.990,0.131) (-0.988,0.171) (-0.985,0.212) (-0.983,0.253) (-0.981,0.293) (-0.980,0.334) (-0.978,0.374) (-0.978,0.415) (-0.977,0.456) (-0.978,0.496) (-0.978,0.537) (-0.980,0.577) (-0.982,0.618) (-0.984,0.659) (-0.987,0.699) (-0.991,0.740) (-0.995,0.780) (-1.000,0.820) (-1.005,0.861) (-1.011,0.901) (-1.018,0.941) (-1.025,0.981) (-1.032,1.021) (-1.040,1.061) (-1.049,1.100) (-1.057,1.140) (-1.067,1.180) (-1.077,1.219) (-1.087,1.258) (-1.097,1.298) (-1.108,1.337) (-1.119,1.376) (-1.131,1.415) (-1.142,1.454) (-1.155,1.493) (-1.167,1.531) (-1.179,1.570) (-1.192,1.608) (-1.205,1.647) (-1.219,1.685) (-1.232,1.724) (-1.246,1.762) (-1.260,1.800)};
\addplot[lineaE, flechaen=0.500] coordinates {(-0.950,0.121) (-0.935,0.159) (-0.919,0.197) (-0.905,0.235) (-0.890,0.273) (-0.876,0.312) (-0.862,0.350) (-0.849,0.389) (-0.836,0.428) (-0.824,0.467) (-0.813,0.506) (-0.803,0.545) (-0.793,0.585) (-0.784,0.625) (-0.776,0.665) (-0.768,0.705) (-0.762,0.745) (-0.756,0.786) (-0.752,0.826) (-0.748,0.867) (-0.745,0.908) (-0.743,0.948) (-0.741,0.989) (-0.741,1.030) (-0.741,1.071) (-0.742,1.112) (-0.743,1.153) (-0.745,1.193) (-0.748,1.234) (-0.751,1.275) (-0.755,1.316) (-0.759,1.356) (-0.763,1.397) (-0.768,1.437) (-0.773,1.478) (-0.779,1.518) (-0.785,1.559) (-0.791,1.599) (-0.798,1.639) (-0.805,1.680) (-0.812,1.720) (-0.819,1.760) (-0.827,1.800)};
\addplot[lineaE, flechaen=0.457] coordinates {(-0.909,0.095) (-0.882,0.124) (-0.854,0.154) (-0.827,0.183) (-0.800,0.213) (-0.773,0.243) (-0.746,0.274) (-0.720,0.305) (-0.695,0.336) (-0.670,0.368) (-0.645,0.400) (-0.622,0.433) (-0.599,0.466) (-0.578,0.500) (-0.557,0.535) (-0.538,0.571) (-0.520,0.607) (-0.503,0.643) (-0.488,0.681) (-0.474,0.719) (-0.462,0.757) (-0.450,0.796) (-0.440,0.835) (-0.432,0.874) (-0.424,0.914) (-0.417,0.954) (-0.412,0.994) (-0.407,1.034) (-0.403,1.074) (-0.400,1.114) (-0.398,1.155) (-0.396,1.195) (-0.394,1.235) (-0.394,1.276) (-0.393,1.316) (-0.393,1.356) (-0.394,1.397) (-0.395,1.437) (-0.396,1.478) (-0.397,1.518) (-0.399,1.558) (-0.401,1.599) (-0.403,1.639) (-0.406,1.679) (-0.408,1.719) (-0.411,1.760) (-0.414,1.800)};
\addplot[lineaE, flechaen=0.500] coordinates {(0.000,0.000) (0.000,0.041) (0.000,0.082) (0.000,0.123) (0.000,0.164) (0.000,0.205) (0.000,0.246) (0.000,0.286) (0.000,0.327) (0.000,0.368) (0.000,0.409) (0.000,0.450) (0.000,0.491) (0.000,0.532) (0.000,0.573) (0.000,0.614) (0.000,0.655) (0.000,0.696) (0.000,0.736) (0.000,0.777) (0.000,0.818) (0.000,0.859) (0.000,0.900) (0.000,0.941) (0.000,0.982) (0.000,1.023) (0.000,1.064) (0.000,1.105) (0.000,1.145) (0.000,1.186) (0.000,1.227) (0.000,1.268) (0.000,1.309) (0.000,1.350) (0.000,1.391) (0.000,1.432) (0.000,1.473) (0.000,1.514) (0.000,1.555) (0.000,1.595) (0.000,1.636) (0.000,1.677) (0.000,1.718) (0.000,1.759) (0.000,1.800)};
\addplot[lineaE, flechaen=0.457] coordinates {(0.909,0.095) (0.882,0.124) (0.854,0.154) (0.827,0.183) (0.800,0.213) (0.773,0.243) (0.746,0.274) (0.720,0.305) (0.695,0.336) (0.670,0.368) (0.645,0.400) (0.622,0.433) (0.599,0.466) (0.578,0.500) (0.557,0.535) (0.538,0.571) (0.520,0.607) (0.503,0.643) (0.488,0.681) (0.474,0.719) (0.462,0.757) (0.450,0.796) (0.440,0.835) (0.432,0.874) (0.424,0.914) (0.417,0.954) (0.412,0.994) (0.407,1.034) (0.403,1.074) (0.400,1.114) (0.398,1.155) (0.396,1.195) (0.394,1.235) (0.394,1.276) (0.393,1.316) (0.393,1.356) (0.394,1.397) (0.395,1.437) (0.396,1.478) (0.397,1.518) (0.399,1.558) (0.401,1.599) (0.403,1.639) (0.406,1.679) (0.408,1.719) (0.411,1.760) (0.414,1.800)};
\addplot[lineaE, flechaen=0.500] coordinates {(0.950,0.121) (0.935,0.159) (0.919,0.197) (0.905,0.235) (0.890,0.273) (0.876,0.312) (0.862,0.350) (0.849,0.389) (0.836,0.428) (0.824,0.467) (0.813,0.506) (0.803,0.545) (0.793,0.585) (0.784,0.625) (0.776,0.665) (0.768,0.705) (0.762,0.745) (0.756,0.786) (0.752,0.826) (0.748,0.867) (0.745,0.908) (0.743,0.948) (0.741,0.989) (0.741,1.030) (0.741,1.071) (0.742,1.112) (0.743,1.153) (0.745,1.193) (0.748,1.234) (0.751,1.275) (0.755,1.316) (0.759,1.356) (0.763,1.397) (0.768,1.437) (0.773,1.478) (0.779,1.518) (0.785,1.559) (0.791,1.599) (0.798,1.639) (0.805,1.680) (0.812,1.720) (0.819,1.760) (0.827,1.800)};
\addplot[lineaE, flechaen=0.500] coordinates {(0.990,0.131) (0.988,0.171) (0.985,0.212) (0.983,0.253) (0.981,0.293) (0.980,0.334) (0.978,0.374) (0.978,0.415) (0.977,0.456) (0.978,0.496) (0.978,0.537) (0.980,0.577) (0.982,0.618) (0.984,0.659) (0.987,0.699) (0.991,0.740) (0.995,0.780) (1.000,0.820) (1.005,0.861) (1.011,0.901) (1.018,0.941) (1.025,0.981) (1.032,1.021) (1.040,1.061) (1.049,1.100) (1.057,1.140) (1.067,1.180) (1.077,1.219) (1.087,1.258) (1.097,1.298) (1.108,1.337) (1.119,1.376) (1.131,1.415) (1.142,1.454) (1.155,1.493) (1.167,1.531) (1.179,1.570) (1.192,1.608) (1.205,1.647) (1.219,1.685) (1.232,1.724) (1.246,1.762) (1.260,1.800)};
\addplot[lineaE, flechaen=0.500] coordinates {(1.031,0.126) (1.040,0.166) (1.050,0.206) (1.061,0.245) (1.071,0.285) (1.082,0.324) (1.093,0.364) (1.105,0.403) (1.116,0.442) (1.128,0.481) (1.141,0.520) (1.154,0.559) (1.167,0.597) (1.181,0.636) (1.195,0.674) (1.209,0.713) (1.224,0.751) (1.239,0.789) (1.255,0.827) (1.271,0.864) (1.287,0.902) (1.303,0.939) (1.320,0.976) (1.337,1.014) (1.355,1.050) (1.373,1.087) (1.391,1.124) (1.409,1.161) (1.428,1.197) (1.446,1.233) (1.465,1.269) (1.485,1.306) (1.504,1.341) (1.524,1.377) (1.544,1.413) (1.564,1.449) (1.584,1.484) (1.604,1.520) (1.625,1.555) (1.646,1.590) (1.667,1.625) (1.688,1.660) (1.709,1.695) (1.730,1.730) (1.751,1.765) (1.773,1.800)};
\addplot[lineaE, flechaen=0.661] coordinates {(1.073,0.113) (1.095,0.146) (1.117,0.180) (1.139,0.213) (1.161,0.246) (1.184,0.280) (1.207,0.312) (1.229,0.345) (1.253,0.378) (1.276,0.411) (1.299,0.443) (1.323,0.475) (1.347,0.508) (1.371,0.539) (1.395,0.571) (1.420,0.603) (1.444,0.634) (1.469,0.666) (1.494,0.697) (1.520,0.728) (1.545,0.759) (1.571,0.790) (1.597,0.820) (1.623,0.851) (1.649,0.881) (1.675,0.911) (1.701,0.941) (1.728,0.971) (1.755,1.001) (1.781,1.031) (1.808,1.061) (1.835,1.090) (1.863,1.119) (1.890,1.149) (1.917,1.178) (1.945,1.207) (1.972,1.236) (2.000,1.265) (2.028,1.294) (2.056,1.322) (2.084,1.351) (2.112,1.380) (2.140,1.408) (2.168,1.437) (2.196,1.465) (2.225,1.493) (2.253,1.521) (2.282,1.549) (2.310,1.578) (2.339,1.606) (2.368,1.633) (2.396,1.661) (2.425,1.689) (2.454,1.717) (2.483,1.745) (2.512,1.772) (2.541,1.800)};
\addplot[lineaE, flechaen=0.604] coordinates {(1.111,0.069) (1.146,0.090) (1.180,0.111) (1.214,0.132) (1.249,0.153) (1.283,0.174) (1.318,0.194) (1.353,0.215) (1.387,0.235) (1.422,0.256) (1.457,0.276) (1.492,0.296) (1.527,0.316) (1.562,0.336) (1.597,0.355) (1.632,0.375) (1.668,0.395) (1.703,0.414) (1.738,0.433) (1.774,0.453) (1.809,0.472) (1.845,0.491) (1.880,0.510) (1.916,0.529) (1.951,0.547) (1.987,0.566) (2.023,0.585) (2.059,0.603) (2.095,0.622) (2.130,0.640) (2.166,0.658) (2.202,0.676) (2.238,0.694) (2.274,0.712) (2.310,0.730) (2.346,0.748) (2.383,0.766) (2.419,0.784) (2.455,0.802) (2.491,0.819) (2.527,0.837) (2.564,0.854) (2.600,0.872)};
\addplot[lineaE, flechaen=0.604] coordinates {(-1.111,-0.069) (-1.146,-0.090) (-1.180,-0.111) (-1.214,-0.132) (-1.249,-0.153) (-1.283,-0.174) (-1.318,-0.194) (-1.353,-0.215) (-1.387,-0.235) (-1.422,-0.256) (-1.457,-0.276) (-1.492,-0.296) (-1.527,-0.316) (-1.562,-0.336) (-1.597,-0.355) (-1.632,-0.375) (-1.668,-0.395) (-1.703,-0.414) (-1.738,-0.433) (-1.774,-0.453) (-1.809,-0.472) (-1.845,-0.491) (-1.880,-0.510) (-1.916,-0.529) (-1.951,-0.547) (-1.987,-0.566) (-2.023,-0.585) (-2.059,-0.603) (-2.095,-0.622) (-2.130,-0.640) (-2.166,-0.658) (-2.202,-0.676) (-2.238,-0.694) (-2.274,-0.712) (-2.310,-0.730) (-2.346,-0.748) (-2.383,-0.766) (-2.419,-0.784) (-2.455,-0.802) (-2.491,-0.819) (-2.527,-0.837) (-2.564,-0.854) (-2.600,-0.872)};
\addplot[lineaE, flechaen=0.661] coordinates {(-1.073,-0.113) (-1.095,-0.146) (-1.117,-0.180) (-1.139,-0.213) (-1.161,-0.246) (-1.184,-0.280) (-1.207,-0.312) (-1.229,-0.345) (-1.253,-0.378) (-1.276,-0.411) (-1.299,-0.443) (-1.323,-0.475) (-1.347,-0.508) (-1.371,-0.539) (-1.395,-0.571) (-1.420,-0.603) (-1.444,-0.634) (-1.469,-0.666) (-1.494,-0.697) (-1.520,-0.728) (-1.545,-0.759) (-1.571,-0.790) (-1.597,-0.820) (-1.623,-0.851) (-1.649,-0.881) (-1.675,-0.911) (-1.701,-0.941) (-1.728,-0.971) (-1.755,-1.001) (-1.781,-1.031) (-1.808,-1.061) (-1.835,-1.090) (-1.863,-1.119) (-1.890,-1.149) (-1.917,-1.178) (-1.945,-1.207) (-1.972,-1.236) (-2.000,-1.265) (-2.028,-1.294) (-2.056,-1.322) (-2.084,-1.351) (-2.112,-1.380) (-2.140,-1.408) (-2.168,-1.437) (-2.196,-1.465) (-2.225,-1.493) (-2.253,-1.521) (-2.282,-1.549) (-2.310,-1.578) (-2.339,-1.606) (-2.368,-1.633) (-2.396,-1.661) (-2.425,-1.689) (-2.454,-1.717) (-2.483,-1.745) (-2.512,-1.772) (-2.541,-1.800)};
\addplot[lineaE, flechaen=0.500] coordinates {(-1.031,-0.126) (-1.040,-0.166) (-1.050,-0.206) (-1.061,-0.245) (-1.071,-0.285) (-1.082,-0.324) (-1.093,-0.364) (-1.105,-0.403) (-1.116,-0.442) (-1.128,-0.481) (-1.141,-0.520) (-1.154,-0.559) (-1.167,-0.597) (-1.181,-0.636) (-1.195,-0.674) (-1.209,-0.713) (-1.224,-0.751) (-1.239,-0.789) (-1.255,-0.827) (-1.271,-0.864) (-1.287,-0.902) (-1.303,-0.939) (-1.320,-0.976) (-1.337,-1.014) (-1.355,-1.050) (-1.373,-1.087) (-1.391,-1.124) (-1.409,-1.161) (-1.428,-1.197) (-1.446,-1.233) (-1.465,-1.269) (-1.485,-1.306) (-1.504,-1.341) (-1.524,-1.377) (-1.544,-1.413) (-1.564,-1.449) (-1.584,-1.484) (-1.604,-1.520) (-1.625,-1.555) (-1.646,-1.590) (-1.667,-1.625) (-1.688,-1.660) (-1.709,-1.695) (-1.730,-1.730) (-1.751,-1.765) (-1.773,-1.800)};
\addplot[lineaE, flechaen=0.500] coordinates {(-0.990,-0.131) (-0.988,-0.171) (-0.985,-0.212) (-0.983,-0.253) (-0.981,-0.293) (-0.980,-0.334) (-0.978,-0.374) (-0.978,-0.415) (-0.977,-0.456) (-0.978,-0.496) (-0.978,-0.537) (-0.980,-0.577) (-0.982,-0.618) (-0.984,-0.659) (-0.987,-0.699) (-0.991,-0.740) (-0.995,-0.780) (-1.000,-0.820) (-1.005,-0.861) (-1.011,-0.901) (-1.018,-0.941) (-1.025,-0.981) (-1.032,-1.021) (-1.040,-1.061) (-1.049,-1.100) (-1.057,-1.140) (-1.067,-1.180) (-1.077,-1.219) (-1.087,-1.258) (-1.097,-1.298) (-1.108,-1.337) (-1.119,-1.376) (-1.131,-1.415) (-1.142,-1.454) (-1.155,-1.493) (-1.167,-1.531) (-1.179,-1.570) (-1.192,-1.608) (-1.205,-1.647) (-1.219,-1.685) (-1.232,-1.724) (-1.246,-1.762) (-1.260,-1.800)};
\addplot[lineaE, flechaen=0.500] coordinates {(-0.950,-0.121) (-0.935,-0.159) (-0.919,-0.197) (-0.905,-0.235) (-0.890,-0.273) (-0.876,-0.312) (-0.862,-0.350) (-0.849,-0.389) (-0.836,-0.428) (-0.824,-0.467) (-0.813,-0.506) (-0.803,-0.545) (-0.793,-0.585) (-0.784,-0.625) (-0.776,-0.665) (-0.768,-0.705) (-0.762,-0.745) (-0.756,-0.786) (-0.752,-0.826) (-0.748,-0.867) (-0.745,-0.908) (-0.743,-0.948) (-0.741,-0.989) (-0.741,-1.030) (-0.741,-1.071) (-0.742,-1.112) (-0.743,-1.153) (-0.745,-1.193) (-0.748,-1.234) (-0.751,-1.275) (-0.755,-1.316) (-0.759,-1.356) (-0.763,-1.397) (-0.768,-1.437) (-0.773,-1.478) (-0.779,-1.518) (-0.785,-1.559) (-0.791,-1.599) (-0.798,-1.639) (-0.805,-1.680) (-0.812,-1.720) (-0.819,-1.760) (-0.827,-1.800)};
\addplot[lineaE, flechaen=0.457] coordinates {(-0.909,-0.095) (-0.882,-0.124) (-0.854,-0.154) (-0.827,-0.183) (-0.800,-0.213) (-0.773,-0.243) (-0.746,-0.274) (-0.720,-0.305) (-0.695,-0.336) (-0.670,-0.368) (-0.645,-0.400) (-0.622,-0.433) (-0.599,-0.466) (-0.578,-0.500) (-0.557,-0.535) (-0.538,-0.571) (-0.520,-0.607) (-0.503,-0.643) (-0.488,-0.681) (-0.474,-0.719) (-0.462,-0.757) (-0.450,-0.796) (-0.440,-0.835) (-0.432,-0.874) (-0.424,-0.914) (-0.417,-0.954) (-0.412,-0.994) (-0.407,-1.034) (-0.403,-1.074) (-0.400,-1.114) (-0.398,-1.155) (-0.396,-1.195) (-0.394,-1.235) (-0.394,-1.276) (-0.393,-1.316) (-0.393,-1.356) (-0.394,-1.397) (-0.395,-1.437) (-0.396,-1.478) (-0.397,-1.518) (-0.399,-1.558) (-0.401,-1.599) (-0.403,-1.639) (-0.406,-1.679) (-0.408,-1.719) (-0.411,-1.760) (-0.414,-1.800)};
\addplot[lineaE, flechaen=0.500] coordinates {(0.000,-0.000) (0.000,-0.041) (0.000,-0.082) (0.000,-0.123) (0.000,-0.164) (0.000,-0.205) (0.000,-0.246) (0.000,-0.286) (0.000,-0.327) (0.000,-0.368) (0.000,-0.409) (0.000,-0.450) (0.000,-0.491) (0.000,-0.532) (0.000,-0.573) (0.000,-0.614) (0.000,-0.655) (0.000,-0.696) (0.000,-0.736) (0.000,-0.777) (0.000,-0.818) (0.000,-0.859) (0.000,-0.900) (0.000,-0.941) (0.000,-0.982) (0.000,-1.023) (0.000,-1.064) (0.000,-1.105) (0.000,-1.145) (0.000,-1.186) (0.000,-1.227) (0.000,-1.268) (0.000,-1.309) (0.000,-1.350) (0.000,-1.391) (0.000,-1.432) (0.000,-1.473) (0.000,-1.514) (0.000,-1.555) (0.000,-1.595) (0.000,-1.636) (0.000,-1.677) (0.000,-1.718) (0.000,-1.759) (0.000,-1.800)};
\addplot[lineaE, flechaen=0.457] coordinates {(0.909,-0.095) (0.882,-0.124) (0.854,-0.154) (0.827,-0.183) (0.800,-0.213) (0.773,-0.243) (0.746,-0.274) (0.720,-0.305) (0.695,-0.336) (0.670,-0.368) (0.645,-0.400) (0.622,-0.433) (0.599,-0.466) (0.578,-0.500) (0.557,-0.535) (0.538,-0.571) (0.520,-0.607) (0.503,-0.643) (0.488,-0.681) (0.474,-0.719) (0.462,-0.757) (0.450,-0.796) (0.440,-0.835) (0.432,-0.874) (0.424,-0.914) (0.417,-0.954) (0.412,-0.994) (0.407,-1.034) (0.403,-1.074) (0.400,-1.114) (0.398,-1.155) (0.396,-1.195) (0.394,-1.235) (0.394,-1.276) (0.393,-1.316) (0.393,-1.356) (0.394,-1.397) (0.395,-1.437) (0.396,-1.478) (0.397,-1.518) (0.399,-1.558) (0.401,-1.599) (0.403,-1.639) (0.406,-1.679) (0.408,-1.719) (0.411,-1.760) (0.414,-1.800)};
\addplot[lineaE, flechaen=0.500] coordinates {(0.950,-0.121) (0.935,-0.159) (0.919,-0.197) (0.905,-0.235) (0.890,-0.273) (0.876,-0.312) (0.862,-0.350) (0.849,-0.389) (0.836,-0.428) (0.824,-0.467) (0.813,-0.506) (0.803,-0.545) (0.793,-0.585) (0.784,-0.625) (0.776,-0.665) (0.768,-0.705) (0.762,-0.745) (0.756,-0.786) (0.752,-0.826) (0.748,-0.867) (0.745,-0.908) (0.743,-0.948) (0.741,-0.989) (0.741,-1.030) (0.741,-1.071) (0.742,-1.112) (0.743,-1.153) (0.745,-1.193) (0.748,-1.234) (0.751,-1.275) (0.755,-1.316) (0.759,-1.356) (0.763,-1.397) (0.768,-1.437) (0.773,-1.478) (0.779,-1.518) (0.785,-1.559) (0.791,-1.599) (0.798,-1.639) (0.805,-1.680) (0.812,-1.720) (0.819,-1.760) (0.827,-1.800)};
\addplot[lineaE, flechaen=0.500] coordinates {(0.990,-0.131) (0.988,-0.171) (0.985,-0.212) (0.983,-0.253) (0.981,-0.293) (0.980,-0.334) (0.978,-0.374) (0.978,-0.415) (0.977,-0.456) (0.978,-0.496) (0.978,-0.537) (0.980,-0.577) (0.982,-0.618) (0.984,-0.659) (0.987,-0.699) (0.991,-0.740) (0.995,-0.780) (1.000,-0.820) (1.005,-0.861) (1.011,-0.901) (1.018,-0.941) (1.025,-0.981) (1.032,-1.021) (1.040,-1.061) (1.049,-1.100) (1.057,-1.140) (1.067,-1.180) (1.077,-1.219) (1.087,-1.258) (1.097,-1.298) (1.108,-1.337) (1.119,-1.376) (1.131,-1.415) (1.142,-1.454) (1.155,-1.493) (1.167,-1.531) (1.179,-1.570) (1.192,-1.608) (1.205,-1.647) (1.219,-1.685) (1.232,-1.724) (1.246,-1.762) (1.260,-1.800)};
\addplot[lineaE, flechaen=0.500] coordinates {(1.031,-0.126) (1.040,-0.166) (1.050,-0.206) (1.061,-0.245) (1.071,-0.285) (1.082,-0.324) (1.093,-0.364) (1.105,-0.403) (1.116,-0.442) (1.128,-0.481) (1.141,-0.520) (1.154,-0.559) (1.167,-0.597) (1.181,-0.636) (1.195,-0.674) (1.209,-0.713) (1.224,-0.751) (1.239,-0.789) (1.255,-0.827) (1.271,-0.864) (1.287,-0.902) (1.303,-0.939) (1.320,-0.976) (1.337,-1.014) (1.355,-1.050) (1.373,-1.087) (1.391,-1.124) (1.409,-1.161) (1.428,-1.197) (1.446,-1.233) (1.465,-1.269) (1.485,-1.306) (1.504,-1.341) (1.524,-1.377) (1.544,-1.413) (1.564,-1.449) (1.584,-1.484) (1.604,-1.520) (1.625,-1.555) (1.646,-1.590) (1.667,-1.625) (1.688,-1.660) (1.709,-1.695) (1.730,-1.730) (1.751,-1.765) (1.773,-1.800)};
\addplot[lineaE, flechaen=0.661] coordinates {(1.073,-0.113) (1.095,-0.146) (1.117,-0.180) (1.139,-0.213) (1.161,-0.246) (1.184,-0.280) (1.207,-0.312) (1.229,-0.345) (1.253,-0.378) (1.276,-0.411) (1.299,-0.443) (1.323,-0.475) (1.347,-0.508) (1.371,-0.539) (1.395,-0.571) (1.420,-0.603) (1.444,-0.634) (1.469,-0.666) (1.494,-0.697) (1.520,-0.728) (1.545,-0.759) (1.571,-0.790) (1.597,-0.820) (1.623,-0.851) (1.649,-0.881) (1.675,-0.911) (1.701,-0.941) (1.728,-0.971) (1.755,-1.001) (1.781,-1.031) (1.808,-1.061) (1.835,-1.090) (1.863,-1.119) (1.890,-1.149) (1.917,-1.178) (1.945,-1.207) (1.972,-1.236) (2.000,-1.265) (2.028,-1.294) (2.056,-1.322) (2.084,-1.351) (2.112,-1.380) (2.140,-1.408) (2.168,-1.437) (2.196,-1.465) (2.225,-1.493) (2.253,-1.521) (2.282,-1.549) (2.310,-1.578) (2.339,-1.606) (2.368,-1.633) (2.396,-1.661) (2.425,-1.689) (2.454,-1.717) (2.483,-1.745) (2.512,-1.772) (2.541,-1.800)};
\addplot[lineaE, flechaen=0.604] coordinates {(1.111,-0.069) (1.146,-0.090) (1.180,-0.111) (1.214,-0.132) (1.249,-0.153) (1.283,-0.174) (1.318,-0.194) (1.353,-0.215) (1.387,-0.235) (1.422,-0.256) (1.457,-0.276) (1.492,-0.296) (1.527,-0.316) (1.562,-0.336) (1.597,-0.355) (1.632,-0.375) (1.668,-0.395) (1.703,-0.414) (1.738,-0.433) (1.774,-0.453) (1.809,-0.472) (1.845,-0.491) (1.880,-0.510) (1.916,-0.529) (1.951,-0.547) (1.987,-0.566) (2.023,-0.585) (2.059,-0.603) (2.095,-0.622) (2.130,-0.640) (2.166,-0.658) (2.202,-0.676) (2.238,-0.694) (2.274,-0.712) (2.310,-0.730) (2.346,-0.748) (2.383,-0.766) (2.419,-0.784) (2.455,-0.802) (2.491,-0.819) (2.527,-0.837) (2.564,-0.854) (2.600,-0.872)};
\addplot[lineaE, flechaen=0.592] coordinates {(-1.130,0.000) (-2.600,0.000)};
\addplot[lineaE, flechaen=0.592] coordinates {(1.130,0.000) (2.600,0.000)};

  \fill[cargapos] (axis cs:-1,0) circle[radius=0.075];
  \fill[cargapos] (axis cs:1,0) circle[radius=0.075];
\end{axis}
\end{tikzpicture}
\caption{$+q$ y $+q$: $u_e$}
\end{subfigure}\hfill
\begin{subfigure}[b]{0.325\textwidth}\centering
\begin{tikzpicture}
\begin{axis}[mapa]
  \addplot graphics[xmin=-2.6, xmax=2.6, ymin=-1.8, ymax=1.8] {t5-ue-opuestas-\modotema};
  %% Generado por scripts/tema5.py -- no editar a mano
\addplot[lineaE, flechaen=0.419] coordinates {(2.600,1.027) (2.568,1.002) (2.537,0.977) (2.505,0.952) (2.473,0.928) (2.441,0.903) (2.409,0.879) (2.376,0.855) (2.344,0.831) (2.311,0.807) (2.279,0.784) (2.246,0.760) (2.213,0.737) (2.180,0.714) (2.147,0.691) (2.114,0.668) (2.081,0.645) (2.048,0.622) (2.015,0.600) (1.981,0.578) (1.947,0.555) (1.914,0.533) (1.880,0.511) (1.846,0.490) (1.812,0.468) (1.778,0.447) (1.744,0.425) (1.710,0.404) (1.675,0.383) (1.641,0.362) (1.607,0.341) (1.572,0.321) (1.537,0.300) (1.503,0.280) (1.468,0.259) (1.433,0.239) (1.398,0.219) (1.363,0.199) (1.328,0.180) (1.293,0.160) (1.258,0.140) (1.222,0.121) (1.187,0.102) (1.152,0.082) (1.117,0.063)};
\addplot[lineaE, flechaen=0.581] coordinates {(-1.117,0.063) (-1.152,0.082) (-1.187,0.102) (-1.222,0.121) (-1.258,0.140) (-1.293,0.160) (-1.328,0.180) (-1.363,0.199) (-1.398,0.219) (-1.433,0.239) (-1.468,0.259) (-1.503,0.280) (-1.537,0.300) (-1.572,0.321) (-1.607,0.341) (-1.641,0.362) (-1.675,0.383) (-1.710,0.404) (-1.744,0.425) (-1.778,0.447) (-1.812,0.468) (-1.846,0.490) (-1.880,0.511) (-1.914,0.533) (-1.947,0.555) (-1.981,0.578) (-2.015,0.600) (-2.048,0.622) (-2.081,0.645) (-2.114,0.668) (-2.147,0.691) (-2.180,0.714) (-2.213,0.737) (-2.246,0.760) (-2.279,0.784) (-2.311,0.807) (-2.344,0.831) (-2.376,0.855) (-2.409,0.879) (-2.441,0.903) (-2.473,0.928) (-2.505,0.952) (-2.537,0.977) (-2.568,1.002) (-2.600,1.027)};
\addplot[lineaE, flechaen=0.159] coordinates {(2.112,1.800) (2.099,1.761) (2.086,1.723) (2.072,1.684) (2.058,1.646) (2.043,1.608) (2.028,1.570) (2.013,1.532) (1.997,1.495) (1.981,1.457) (1.964,1.420) (1.948,1.383) (1.930,1.346) (1.913,1.309) (1.895,1.272) (1.877,1.235) (1.858,1.199) (1.840,1.163) (1.820,1.127) (1.801,1.091) (1.781,1.055) (1.761,1.020) (1.741,0.984) (1.720,0.949) (1.699,0.914) (1.677,0.880) (1.656,0.845) (1.634,0.811) (1.611,0.777) (1.589,0.743) (1.566,0.709) (1.543,0.675) (1.519,0.642) (1.496,0.608) (1.472,0.575) (1.447,0.542) (1.423,0.510) (1.398,0.477) (1.373,0.445) (1.348,0.413) (1.323,0.381) (1.297,0.349) (1.272,0.317) (1.246,0.286) (1.220,0.254) (1.193,0.223) (1.167,0.192) (1.140,0.161) (1.114,0.130) (1.087,0.099)};
\addplot[lineaE, flechaen=0.841] coordinates {(-1.087,0.099) (-1.114,0.130) (-1.140,0.161) (-1.167,0.192) (-1.193,0.223) (-1.220,0.254) (-1.246,0.286) (-1.272,0.317) (-1.297,0.349) (-1.323,0.381) (-1.348,0.413) (-1.373,0.445) (-1.398,0.477) (-1.423,0.510) (-1.447,0.542) (-1.472,0.575) (-1.496,0.608) (-1.519,0.642) (-1.543,0.675) (-1.566,0.709) (-1.589,0.743) (-1.611,0.777) (-1.634,0.811) (-1.656,0.845) (-1.677,0.880) (-1.699,0.914) (-1.720,0.949) (-1.741,0.984) (-1.761,1.020) (-1.781,1.055) (-1.801,1.091) (-1.820,1.127) (-1.840,1.163) (-1.858,1.199) (-1.877,1.235) (-1.895,1.272) (-1.913,1.309) (-1.930,1.346) (-1.948,1.383) (-1.964,1.420) (-1.981,1.457) (-1.997,1.495) (-2.013,1.532) (-2.028,1.570) (-2.043,1.608) (-2.058,1.646) (-2.072,1.684) (-2.086,1.723) (-2.099,1.761) (-2.112,1.800)};
\addplot[lineaE, flechaen=0.500] coordinates {(1.459,1.800) (1.461,1.759) (1.463,1.719) (1.464,1.678) (1.465,1.637) (1.464,1.597) (1.463,1.556) (1.462,1.515) (1.460,1.474) (1.457,1.434) (1.454,1.393) (1.450,1.353) (1.445,1.312) (1.440,1.272) (1.434,1.232) (1.427,1.191) (1.420,1.151) (1.413,1.111) (1.405,1.071) (1.396,1.032) (1.387,0.992) (1.377,0.952) (1.367,0.913) (1.356,0.874) (1.345,0.835) (1.333,0.796) (1.321,0.757) (1.309,0.718) (1.296,0.679) (1.282,0.641) (1.268,0.603) (1.254,0.565) (1.239,0.527) (1.224,0.489) (1.209,0.451) (1.193,0.414) (1.177,0.376) (1.160,0.339) (1.144,0.302) (1.127,0.265) (1.110,0.228) (1.092,0.191) (1.075,0.154) (1.057,0.117)};
\addplot[lineaE, flechaen=0.500] coordinates {(-1.057,0.117) (-1.075,0.154) (-1.092,0.191) (-1.110,0.228) (-1.127,0.265) (-1.144,0.302) (-1.160,0.339) (-1.177,0.376) (-1.193,0.414) (-1.209,0.451) (-1.224,0.489) (-1.239,0.527) (-1.254,0.565) (-1.268,0.603) (-1.282,0.641) (-1.296,0.679) (-1.309,0.718) (-1.321,0.757) (-1.333,0.796) (-1.345,0.835) (-1.356,0.874) (-1.367,0.913) (-1.377,0.952) (-1.387,0.992) (-1.396,1.032) (-1.405,1.071) (-1.413,1.111) (-1.420,1.151) (-1.427,1.191) (-1.434,1.232) (-1.440,1.272) (-1.445,1.312) (-1.450,1.353) (-1.454,1.393) (-1.457,1.434) (-1.460,1.474) (-1.462,1.515) (-1.463,1.556) (-1.464,1.597) (-1.465,1.637) (-1.464,1.678) (-1.463,1.719) (-1.461,1.759) (-1.459,1.800)};
\addplot[lineaE, flechaen=0.500] coordinates {(0.904,1.800) (0.924,1.765) (0.943,1.728) (0.961,1.692) (0.979,1.655) (0.995,1.617) (1.010,1.579) (1.024,1.541) (1.037,1.502) (1.050,1.463) (1.061,1.424) (1.071,1.385) (1.081,1.345) (1.089,1.305) (1.097,1.265) (1.103,1.225) (1.109,1.184) (1.114,1.144) (1.119,1.103) (1.122,1.062) (1.125,1.022) (1.127,0.981) (1.128,0.940) (1.128,0.899) (1.128,0.858) (1.127,0.817) (1.125,0.777) (1.123,0.736) (1.120,0.695) (1.117,0.654) (1.112,0.614) (1.108,0.573) (1.103,0.533) (1.097,0.492) (1.091,0.452) (1.084,0.411) (1.078,0.371) (1.070,0.331) (1.063,0.291) (1.055,0.251) (1.046,0.211) (1.038,0.171) (1.029,0.131)};
\addplot[lineaE, flechaen=0.500] coordinates {(-1.029,0.131) (-1.038,0.171) (-1.046,0.211) (-1.055,0.251) (-1.063,0.291) (-1.070,0.331) (-1.078,0.371) (-1.084,0.411) (-1.091,0.452) (-1.097,0.492) (-1.103,0.533) (-1.108,0.573) (-1.112,0.614) (-1.117,0.654) (-1.120,0.695) (-1.123,0.736) (-1.125,0.777) (-1.127,0.817) (-1.128,0.858) (-1.128,0.899) (-1.128,0.940) (-1.127,0.981) (-1.125,1.022) (-1.122,1.062) (-1.119,1.103) (-1.114,1.144) (-1.109,1.184) (-1.103,1.225) (-1.097,1.265) (-1.089,1.305) (-1.081,1.345) (-1.071,1.385) (-1.061,1.424) (-1.050,1.463) (-1.037,1.502) (-1.024,1.541) (-1.010,1.579) (-0.995,1.617) (-0.979,1.655) (-0.961,1.692) (-0.943,1.728) (-0.924,1.765) (-0.904,1.800)};
\addplot[lineaE, flechaen=0.386, flechaen=0.614] coordinates {(-1.000,0.131) (-0.999,0.171) (-0.999,0.211) (-0.998,0.251) (-0.997,0.291) (-0.996,0.331) (-0.994,0.371) (-0.991,0.411) (-0.989,0.451) (-0.986,0.491) (-0.982,0.531) (-0.978,0.571) (-0.973,0.611) (-0.967,0.651) (-0.961,0.690) (-0.954,0.730) (-0.946,0.769) (-0.937,0.808) (-0.928,0.847) (-0.918,0.886) (-0.907,0.925) (-0.895,0.963) (-0.882,1.001) (-0.868,1.038) (-0.853,1.076) (-0.837,1.113) (-0.821,1.149) (-0.803,1.185) (-0.784,1.220) (-0.764,1.255) (-0.742,1.289) (-0.720,1.322) (-0.697,1.355) (-0.672,1.386) (-0.646,1.417) (-0.620,1.447) (-0.592,1.476) (-0.563,1.503) (-0.532,1.530) (-0.501,1.555) (-0.469,1.579) (-0.436,1.601) (-0.401,1.622) (-0.366,1.641) (-0.330,1.659) (-0.293,1.675) (-0.256,1.689) (-0.218,1.701) (-0.179,1.711) (-0.140,1.719) (-0.100,1.726) (-0.060,1.730) (-0.020,1.732) (0.020,1.732) (0.060,1.730) (0.100,1.726) (0.140,1.719) (0.179,1.711) (0.218,1.701) (0.256,1.689) (0.293,1.675) (0.330,1.659) (0.366,1.641) (0.401,1.622) (0.436,1.601) (0.469,1.579) (0.501,1.555) (0.532,1.530) (0.563,1.503) (0.592,1.476) (0.620,1.447) (0.646,1.417) (0.672,1.386) (0.697,1.355) (0.720,1.322) (0.742,1.289) (0.764,1.255) (0.784,1.220) (0.803,1.185) (0.821,1.149) (0.837,1.113) (0.853,1.076) (0.868,1.038) (0.882,1.001) (0.895,0.963) (0.907,0.925) (0.918,0.886) (0.928,0.847) (0.937,0.808) (0.946,0.769) (0.954,0.730) (0.961,0.690) (0.967,0.651) (0.973,0.611) (0.978,0.571) (0.982,0.531) (0.986,0.491) (0.989,0.451) (0.991,0.411) (0.994,0.371) (0.996,0.331) (0.997,0.291) (0.998,0.251) (0.999,0.211) (0.999,0.171) (1.000,0.131)};
\addplot[lineaE, flechaen=0.353, flechaen=0.647] coordinates {(-0.971,0.128) (-0.962,0.167) (-0.952,0.207) (-0.943,0.246) (-0.933,0.285) (-0.922,0.323) (-0.911,0.362) (-0.900,0.401) (-0.889,0.439) (-0.876,0.477) (-0.863,0.516) (-0.850,0.553) (-0.836,0.591) (-0.821,0.628) (-0.805,0.665) (-0.788,0.702) (-0.770,0.738) (-0.752,0.774) (-0.733,0.809) (-0.712,0.843) (-0.691,0.877) (-0.668,0.911) (-0.645,0.943) (-0.620,0.975) (-0.594,1.006) (-0.567,1.036) (-0.539,1.064) (-0.510,1.092) (-0.479,1.118) (-0.448,1.143) (-0.415,1.166) (-0.381,1.188) (-0.346,1.208) (-0.311,1.227) (-0.274,1.243) (-0.236,1.258) (-0.198,1.270) (-0.159,1.280) (-0.120,1.289) (-0.080,1.294) (-0.040,1.298) (-0.000,1.299) (0.040,1.298) (0.080,1.294) (0.120,1.289) (0.159,1.280) (0.198,1.270) (0.236,1.258) (0.274,1.243) (0.311,1.227) (0.346,1.208) (0.381,1.188) (0.415,1.166) (0.448,1.143) (0.479,1.118) (0.510,1.092) (0.539,1.064) (0.567,1.036) (0.594,1.006) (0.620,0.975) (0.645,0.943) (0.668,0.911) (0.691,0.877) (0.712,0.843) (0.733,0.809) (0.752,0.774) (0.770,0.738) (0.788,0.702) (0.805,0.665) (0.821,0.628) (0.836,0.591) (0.850,0.553) (0.863,0.516) (0.876,0.477) (0.889,0.439) (0.900,0.401) (0.911,0.362) (0.922,0.323) (0.933,0.285) (0.943,0.246) (0.952,0.207) (0.962,0.167) (0.971,0.128)};
\addplot[lineaE, flechaen=0.318, flechaen=0.682] coordinates {(-0.942,0.117) (-0.924,0.154) (-0.906,0.190) (-0.887,0.226) (-0.868,0.262) (-0.849,0.297) (-0.829,0.333) (-0.809,0.368) (-0.789,0.403) (-0.767,0.438) (-0.746,0.472) (-0.723,0.505) (-0.700,0.539) (-0.676,0.571) (-0.651,0.603) (-0.625,0.635) (-0.598,0.665) (-0.571,0.695) (-0.542,0.724) (-0.512,0.751) (-0.482,0.778) (-0.450,0.803) (-0.417,0.827) (-0.383,0.849) (-0.348,0.870) (-0.313,0.889) (-0.276,0.906) (-0.238,0.921) (-0.200,0.934) (-0.161,0.944) (-0.121,0.953) (-0.081,0.959) (-0.041,0.963) (-0.000,0.964) (0.041,0.963) (0.081,0.959) (0.121,0.953) (0.161,0.944) (0.200,0.934) (0.238,0.921) (0.276,0.906) (0.313,0.889) (0.348,0.870) (0.383,0.849) (0.417,0.827) (0.450,0.803) (0.482,0.778) (0.512,0.751) (0.542,0.724) (0.571,0.695) (0.598,0.665) (0.625,0.635) (0.651,0.603) (0.676,0.571) (0.700,0.539) (0.723,0.505) (0.746,0.472) (0.767,0.438) (0.789,0.403) (0.809,0.368) (0.829,0.333) (0.849,0.297) (0.868,0.262) (0.887,0.226) (0.906,0.190) (0.924,0.154) (0.942,0.117)};
\addplot[lineaE, flechaen=0.285, flechaen=0.715] coordinates {(-0.914,0.098) (-0.887,0.128) (-0.859,0.158) (-0.832,0.188) (-0.805,0.218) (-0.777,0.248) (-0.749,0.277) (-0.720,0.305) (-0.691,0.334) (-0.662,0.362) (-0.632,0.389) (-0.602,0.416) (-0.570,0.442) (-0.539,0.467) (-0.506,0.491) (-0.473,0.514) (-0.439,0.536) (-0.404,0.557) (-0.369,0.576) (-0.333,0.594) (-0.296,0.611) (-0.258,0.626) (-0.219,0.638) (-0.180,0.649) (-0.141,0.658) (-0.101,0.665) (-0.061,0.669) (-0.020,0.672) (0.020,0.672) (0.061,0.669) (0.101,0.665) (0.141,0.658) (0.180,0.649) (0.219,0.638) (0.258,0.626) (0.296,0.611) (0.333,0.594) (0.369,0.576) (0.404,0.557) (0.439,0.536) (0.473,0.514) (0.506,0.491) (0.539,0.467) (0.570,0.442) (0.602,0.416) (0.632,0.389) (0.662,0.362) (0.691,0.334) (0.720,0.305) (0.749,0.277) (0.777,0.248) (0.805,0.218) (0.832,0.188) (0.859,0.158) (0.887,0.128) (0.914,0.098)};
\addplot[lineaE, flechaen=0.255, flechaen=0.745] coordinates {(-0.882,0.063) (-0.847,0.082) (-0.812,0.101) (-0.776,0.119) (-0.741,0.138) (-0.705,0.156) (-0.669,0.174) (-0.633,0.191) (-0.597,0.208) (-0.560,0.225) (-0.523,0.241) (-0.486,0.257) (-0.449,0.271) (-0.412,0.285) (-0.374,0.298) (-0.335,0.311) (-0.297,0.322) (-0.258,0.332) (-0.219,0.340) (-0.180,0.348) (-0.140,0.354) (-0.100,0.358) (-0.060,0.361) (-0.020,0.363) (0.020,0.363) (0.060,0.361) (0.100,0.358) (0.140,0.354) (0.180,0.348) (0.219,0.340) (0.258,0.332) (0.297,0.322) (0.335,0.311) (0.374,0.298) (0.412,0.285) (0.449,0.271) (0.486,0.257) (0.523,0.241) (0.560,0.225) (0.597,0.208) (0.633,0.191) (0.669,0.174) (0.705,0.156) (0.741,0.138) (0.776,0.119) (0.812,0.101) (0.847,0.082) (0.882,0.063)};
\addplot[lineaE, flechaen=0.419] coordinates {(2.600,-1.027) (2.568,-1.002) (2.537,-0.977) (2.505,-0.952) (2.473,-0.928) (2.441,-0.903) (2.409,-0.879) (2.376,-0.855) (2.344,-0.831) (2.311,-0.807) (2.279,-0.784) (2.246,-0.760) (2.213,-0.737) (2.180,-0.714) (2.147,-0.691) (2.114,-0.668) (2.081,-0.645) (2.048,-0.622) (2.015,-0.600) (1.981,-0.578) (1.947,-0.555) (1.914,-0.533) (1.880,-0.511) (1.846,-0.490) (1.812,-0.468) (1.778,-0.447) (1.744,-0.425) (1.710,-0.404) (1.675,-0.383) (1.641,-0.362) (1.607,-0.341) (1.572,-0.321) (1.537,-0.300) (1.503,-0.280) (1.468,-0.259) (1.433,-0.239) (1.398,-0.219) (1.363,-0.199) (1.328,-0.180) (1.293,-0.160) (1.258,-0.140) (1.222,-0.121) (1.187,-0.102) (1.152,-0.082) (1.117,-0.063)};
\addplot[lineaE, flechaen=0.581] coordinates {(-1.117,-0.063) (-1.152,-0.082) (-1.187,-0.102) (-1.222,-0.121) (-1.258,-0.140) (-1.293,-0.160) (-1.328,-0.180) (-1.363,-0.199) (-1.398,-0.219) (-1.433,-0.239) (-1.468,-0.259) (-1.503,-0.280) (-1.537,-0.300) (-1.572,-0.321) (-1.607,-0.341) (-1.641,-0.362) (-1.675,-0.383) (-1.710,-0.404) (-1.744,-0.425) (-1.778,-0.447) (-1.812,-0.468) (-1.846,-0.490) (-1.880,-0.511) (-1.914,-0.533) (-1.947,-0.555) (-1.981,-0.578) (-2.015,-0.600) (-2.048,-0.622) (-2.081,-0.645) (-2.114,-0.668) (-2.147,-0.691) (-2.180,-0.714) (-2.213,-0.737) (-2.246,-0.760) (-2.279,-0.784) (-2.311,-0.807) (-2.344,-0.831) (-2.376,-0.855) (-2.409,-0.879) (-2.441,-0.903) (-2.473,-0.928) (-2.505,-0.952) (-2.537,-0.977) (-2.568,-1.002) (-2.600,-1.027)};
\addplot[lineaE, flechaen=0.159] coordinates {(2.112,-1.800) (2.099,-1.761) (2.086,-1.723) (2.072,-1.684) (2.058,-1.646) (2.043,-1.608) (2.028,-1.570) (2.013,-1.532) (1.997,-1.495) (1.981,-1.457) (1.964,-1.420) (1.948,-1.383) (1.930,-1.346) (1.913,-1.309) (1.895,-1.272) (1.877,-1.235) (1.858,-1.199) (1.840,-1.163) (1.820,-1.127) (1.801,-1.091) (1.781,-1.055) (1.761,-1.020) (1.741,-0.984) (1.720,-0.949) (1.699,-0.914) (1.677,-0.880) (1.656,-0.845) (1.634,-0.811) (1.611,-0.777) (1.589,-0.743) (1.566,-0.709) (1.543,-0.675) (1.519,-0.642) (1.496,-0.608) (1.472,-0.575) (1.447,-0.542) (1.423,-0.510) (1.398,-0.477) (1.373,-0.445) (1.348,-0.413) (1.323,-0.381) (1.297,-0.349) (1.272,-0.317) (1.246,-0.286) (1.220,-0.254) (1.193,-0.223) (1.167,-0.192) (1.140,-0.161) (1.114,-0.130) (1.087,-0.099)};
\addplot[lineaE, flechaen=0.841] coordinates {(-1.087,-0.099) (-1.114,-0.130) (-1.140,-0.161) (-1.167,-0.192) (-1.193,-0.223) (-1.220,-0.254) (-1.246,-0.286) (-1.272,-0.317) (-1.297,-0.349) (-1.323,-0.381) (-1.348,-0.413) (-1.373,-0.445) (-1.398,-0.477) (-1.423,-0.510) (-1.447,-0.542) (-1.472,-0.575) (-1.496,-0.608) (-1.519,-0.642) (-1.543,-0.675) (-1.566,-0.709) (-1.589,-0.743) (-1.611,-0.777) (-1.634,-0.811) (-1.656,-0.845) (-1.677,-0.880) (-1.699,-0.914) (-1.720,-0.949) (-1.741,-0.984) (-1.761,-1.020) (-1.781,-1.055) (-1.801,-1.091) (-1.820,-1.127) (-1.840,-1.163) (-1.858,-1.199) (-1.877,-1.235) (-1.895,-1.272) (-1.913,-1.309) (-1.930,-1.346) (-1.948,-1.383) (-1.964,-1.420) (-1.981,-1.457) (-1.997,-1.495) (-2.013,-1.532) (-2.028,-1.570) (-2.043,-1.608) (-2.058,-1.646) (-2.072,-1.684) (-2.086,-1.723) (-2.099,-1.761) (-2.112,-1.800)};
\addplot[lineaE, flechaen=0.500] coordinates {(1.459,-1.800) (1.461,-1.759) (1.463,-1.719) (1.464,-1.678) (1.465,-1.637) (1.464,-1.597) (1.463,-1.556) (1.462,-1.515) (1.460,-1.474) (1.457,-1.434) (1.454,-1.393) (1.450,-1.353) (1.445,-1.312) (1.440,-1.272) (1.434,-1.232) (1.427,-1.191) (1.420,-1.151) (1.413,-1.111) (1.405,-1.071) (1.396,-1.032) (1.387,-0.992) (1.377,-0.952) (1.367,-0.913) (1.356,-0.874) (1.345,-0.835) (1.333,-0.796) (1.321,-0.757) (1.309,-0.718) (1.296,-0.679) (1.282,-0.641) (1.268,-0.603) (1.254,-0.565) (1.239,-0.527) (1.224,-0.489) (1.209,-0.451) (1.193,-0.414) (1.177,-0.376) (1.160,-0.339) (1.144,-0.302) (1.127,-0.265) (1.110,-0.228) (1.092,-0.191) (1.075,-0.154) (1.057,-0.117)};
\addplot[lineaE, flechaen=0.500] coordinates {(-1.057,-0.117) (-1.075,-0.154) (-1.092,-0.191) (-1.110,-0.228) (-1.127,-0.265) (-1.144,-0.302) (-1.160,-0.339) (-1.177,-0.376) (-1.193,-0.414) (-1.209,-0.451) (-1.224,-0.489) (-1.239,-0.527) (-1.254,-0.565) (-1.268,-0.603) (-1.282,-0.641) (-1.296,-0.679) (-1.309,-0.718) (-1.321,-0.757) (-1.333,-0.796) (-1.345,-0.835) (-1.356,-0.874) (-1.367,-0.913) (-1.377,-0.952) (-1.387,-0.992) (-1.396,-1.032) (-1.405,-1.071) (-1.413,-1.111) (-1.420,-1.151) (-1.427,-1.191) (-1.434,-1.232) (-1.440,-1.272) (-1.445,-1.312) (-1.450,-1.353) (-1.454,-1.393) (-1.457,-1.434) (-1.460,-1.474) (-1.462,-1.515) (-1.463,-1.556) (-1.464,-1.597) (-1.465,-1.637) (-1.464,-1.678) (-1.463,-1.719) (-1.461,-1.759) (-1.459,-1.800)};
\addplot[lineaE, flechaen=0.500] coordinates {(0.904,-1.800) (0.924,-1.765) (0.943,-1.728) (0.961,-1.692) (0.979,-1.655) (0.995,-1.617) (1.010,-1.579) (1.024,-1.541) (1.037,-1.502) (1.050,-1.463) (1.061,-1.424) (1.071,-1.385) (1.081,-1.345) (1.089,-1.305) (1.097,-1.265) (1.103,-1.225) (1.109,-1.184) (1.114,-1.144) (1.119,-1.103) (1.122,-1.062) (1.125,-1.022) (1.127,-0.981) (1.128,-0.940) (1.128,-0.899) (1.128,-0.858) (1.127,-0.817) (1.125,-0.777) (1.123,-0.736) (1.120,-0.695) (1.117,-0.654) (1.112,-0.614) (1.108,-0.573) (1.103,-0.533) (1.097,-0.492) (1.091,-0.452) (1.084,-0.411) (1.078,-0.371) (1.070,-0.331) (1.063,-0.291) (1.055,-0.251) (1.046,-0.211) (1.038,-0.171) (1.029,-0.131)};
\addplot[lineaE, flechaen=0.500] coordinates {(-1.029,-0.131) (-1.038,-0.171) (-1.046,-0.211) (-1.055,-0.251) (-1.063,-0.291) (-1.070,-0.331) (-1.078,-0.371) (-1.084,-0.411) (-1.091,-0.452) (-1.097,-0.492) (-1.103,-0.533) (-1.108,-0.573) (-1.112,-0.614) (-1.117,-0.654) (-1.120,-0.695) (-1.123,-0.736) (-1.125,-0.777) (-1.127,-0.817) (-1.128,-0.858) (-1.128,-0.899) (-1.128,-0.940) (-1.127,-0.981) (-1.125,-1.022) (-1.122,-1.062) (-1.119,-1.103) (-1.114,-1.144) (-1.109,-1.184) (-1.103,-1.225) (-1.097,-1.265) (-1.089,-1.305) (-1.081,-1.345) (-1.071,-1.385) (-1.061,-1.424) (-1.050,-1.463) (-1.037,-1.502) (-1.024,-1.541) (-1.010,-1.579) (-0.995,-1.617) (-0.979,-1.655) (-0.961,-1.692) (-0.943,-1.728) (-0.924,-1.765) (-0.904,-1.800)};
\addplot[lineaE, flechaen=0.386, flechaen=0.614] coordinates {(-1.000,-0.131) (-0.999,-0.171) (-0.999,-0.211) (-0.998,-0.251) (-0.997,-0.291) (-0.996,-0.331) (-0.994,-0.371) (-0.991,-0.411) (-0.989,-0.451) (-0.986,-0.491) (-0.982,-0.531) (-0.978,-0.571) (-0.973,-0.611) (-0.967,-0.651) (-0.961,-0.690) (-0.954,-0.730) (-0.946,-0.769) (-0.937,-0.808) (-0.928,-0.847) (-0.918,-0.886) (-0.907,-0.925) (-0.895,-0.963) (-0.882,-1.001) (-0.868,-1.038) (-0.853,-1.076) (-0.837,-1.113) (-0.821,-1.149) (-0.803,-1.185) (-0.784,-1.220) (-0.764,-1.255) (-0.742,-1.289) (-0.720,-1.322) (-0.697,-1.355) (-0.672,-1.386) (-0.646,-1.417) (-0.620,-1.447) (-0.592,-1.476) (-0.563,-1.503) (-0.532,-1.530) (-0.501,-1.555) (-0.469,-1.579) (-0.436,-1.601) (-0.401,-1.622) (-0.366,-1.641) (-0.330,-1.659) (-0.293,-1.675) (-0.256,-1.689) (-0.218,-1.701) (-0.179,-1.711) (-0.140,-1.719) (-0.100,-1.726) (-0.060,-1.730) (-0.020,-1.732) (0.020,-1.732) (0.060,-1.730) (0.100,-1.726) (0.140,-1.719) (0.179,-1.711) (0.218,-1.701) (0.256,-1.689) (0.293,-1.675) (0.330,-1.659) (0.366,-1.641) (0.401,-1.622) (0.436,-1.601) (0.469,-1.579) (0.501,-1.555) (0.532,-1.530) (0.563,-1.503) (0.592,-1.476) (0.620,-1.447) (0.646,-1.417) (0.672,-1.386) (0.697,-1.355) (0.720,-1.322) (0.742,-1.289) (0.764,-1.255) (0.784,-1.220) (0.803,-1.185) (0.821,-1.149) (0.837,-1.113) (0.853,-1.076) (0.868,-1.038) (0.882,-1.001) (0.895,-0.963) (0.907,-0.925) (0.918,-0.886) (0.928,-0.847) (0.937,-0.808) (0.946,-0.769) (0.954,-0.730) (0.961,-0.690) (0.967,-0.651) (0.973,-0.611) (0.978,-0.571) (0.982,-0.531) (0.986,-0.491) (0.989,-0.451) (0.991,-0.411) (0.994,-0.371) (0.996,-0.331) (0.997,-0.291) (0.998,-0.251) (0.999,-0.211) (0.999,-0.171) (1.000,-0.131)};
\addplot[lineaE, flechaen=0.353, flechaen=0.647] coordinates {(-0.971,-0.128) (-0.962,-0.167) (-0.952,-0.207) (-0.943,-0.246) (-0.933,-0.285) (-0.922,-0.323) (-0.911,-0.362) (-0.900,-0.401) (-0.889,-0.439) (-0.876,-0.477) (-0.863,-0.516) (-0.850,-0.553) (-0.836,-0.591) (-0.821,-0.628) (-0.805,-0.665) (-0.788,-0.702) (-0.770,-0.738) (-0.752,-0.774) (-0.733,-0.809) (-0.712,-0.843) (-0.691,-0.877) (-0.668,-0.911) (-0.645,-0.943) (-0.620,-0.975) (-0.594,-1.006) (-0.567,-1.036) (-0.539,-1.064) (-0.510,-1.092) (-0.479,-1.118) (-0.448,-1.143) (-0.415,-1.166) (-0.381,-1.188) (-0.346,-1.208) (-0.311,-1.227) (-0.274,-1.243) (-0.236,-1.258) (-0.198,-1.270) (-0.159,-1.280) (-0.120,-1.289) (-0.080,-1.294) (-0.040,-1.298) (-0.000,-1.299) (0.040,-1.298) (0.080,-1.294) (0.120,-1.289) (0.159,-1.280) (0.198,-1.270) (0.236,-1.258) (0.274,-1.243) (0.311,-1.227) (0.346,-1.208) (0.381,-1.188) (0.415,-1.166) (0.448,-1.143) (0.479,-1.118) (0.510,-1.092) (0.539,-1.064) (0.567,-1.036) (0.594,-1.006) (0.620,-0.975) (0.645,-0.943) (0.668,-0.911) (0.691,-0.877) (0.712,-0.843) (0.733,-0.809) (0.752,-0.774) (0.770,-0.738) (0.788,-0.702) (0.805,-0.665) (0.821,-0.628) (0.836,-0.591) (0.850,-0.553) (0.863,-0.516) (0.876,-0.477) (0.889,-0.439) (0.900,-0.401) (0.911,-0.362) (0.922,-0.323) (0.933,-0.285) (0.943,-0.246) (0.952,-0.207) (0.962,-0.167) (0.971,-0.128)};
\addplot[lineaE, flechaen=0.318, flechaen=0.682] coordinates {(-0.942,-0.117) (-0.924,-0.154) (-0.906,-0.190) (-0.887,-0.226) (-0.868,-0.262) (-0.849,-0.297) (-0.829,-0.333) (-0.809,-0.368) (-0.789,-0.403) (-0.767,-0.438) (-0.746,-0.472) (-0.723,-0.505) (-0.700,-0.539) (-0.676,-0.571) (-0.651,-0.603) (-0.625,-0.635) (-0.598,-0.665) (-0.571,-0.695) (-0.542,-0.724) (-0.512,-0.751) (-0.482,-0.778) (-0.450,-0.803) (-0.417,-0.827) (-0.383,-0.849) (-0.348,-0.870) (-0.313,-0.889) (-0.276,-0.906) (-0.238,-0.921) (-0.200,-0.934) (-0.161,-0.944) (-0.121,-0.953) (-0.081,-0.959) (-0.041,-0.963) (-0.000,-0.964) (0.041,-0.963) (0.081,-0.959) (0.121,-0.953) (0.161,-0.944) (0.200,-0.934) (0.238,-0.921) (0.276,-0.906) (0.313,-0.889) (0.348,-0.870) (0.383,-0.849) (0.417,-0.827) (0.450,-0.803) (0.482,-0.778) (0.512,-0.751) (0.542,-0.724) (0.571,-0.695) (0.598,-0.665) (0.625,-0.635) (0.651,-0.603) (0.676,-0.571) (0.700,-0.539) (0.723,-0.505) (0.746,-0.472) (0.767,-0.438) (0.789,-0.403) (0.809,-0.368) (0.829,-0.333) (0.849,-0.297) (0.868,-0.262) (0.887,-0.226) (0.906,-0.190) (0.924,-0.154) (0.942,-0.117)};
\addplot[lineaE, flechaen=0.285, flechaen=0.715] coordinates {(-0.914,-0.098) (-0.887,-0.128) (-0.859,-0.158) (-0.832,-0.188) (-0.805,-0.218) (-0.777,-0.248) (-0.749,-0.277) (-0.720,-0.305) (-0.691,-0.334) (-0.662,-0.362) (-0.632,-0.389) (-0.602,-0.416) (-0.570,-0.442) (-0.539,-0.467) (-0.506,-0.491) (-0.473,-0.514) (-0.439,-0.536) (-0.404,-0.557) (-0.369,-0.576) (-0.333,-0.594) (-0.296,-0.611) (-0.258,-0.626) (-0.219,-0.638) (-0.180,-0.649) (-0.141,-0.658) (-0.101,-0.665) (-0.061,-0.669) (-0.020,-0.672) (0.020,-0.672) (0.061,-0.669) (0.101,-0.665) (0.141,-0.658) (0.180,-0.649) (0.219,-0.638) (0.258,-0.626) (0.296,-0.611) (0.333,-0.594) (0.369,-0.576) (0.404,-0.557) (0.439,-0.536) (0.473,-0.514) (0.506,-0.491) (0.539,-0.467) (0.570,-0.442) (0.602,-0.416) (0.632,-0.389) (0.662,-0.362) (0.691,-0.334) (0.720,-0.305) (0.749,-0.277) (0.777,-0.248) (0.805,-0.218) (0.832,-0.188) (0.859,-0.158) (0.887,-0.128) (0.914,-0.098)};
\addplot[lineaE, flechaen=0.255, flechaen=0.745] coordinates {(-0.882,-0.063) (-0.847,-0.082) (-0.812,-0.101) (-0.776,-0.119) (-0.741,-0.138) (-0.705,-0.156) (-0.669,-0.174) (-0.633,-0.191) (-0.597,-0.208) (-0.560,-0.225) (-0.523,-0.241) (-0.486,-0.257) (-0.449,-0.271) (-0.412,-0.285) (-0.374,-0.298) (-0.335,-0.311) (-0.297,-0.322) (-0.258,-0.332) (-0.219,-0.340) (-0.180,-0.348) (-0.140,-0.354) (-0.100,-0.358) (-0.060,-0.361) (-0.020,-0.363) (0.020,-0.363) (0.060,-0.361) (0.100,-0.358) (0.140,-0.354) (0.180,-0.348) (0.219,-0.340) (0.258,-0.332) (0.297,-0.322) (0.335,-0.311) (0.374,-0.298) (0.412,-0.285) (0.449,-0.271) (0.486,-0.257) (0.523,-0.241) (0.560,-0.225) (0.597,-0.208) (0.633,-0.191) (0.669,-0.174) (0.705,-0.156) (0.741,-0.138) (0.776,-0.119) (0.812,-0.101) (0.847,-0.082) (0.882,-0.063)};
\addplot[lineaE, flechaen=0.241, flechaen=0.759] coordinates {(-0.870,0.000) (0.870,0.000)};
\addplot[lineaE, flechaen=0.592] coordinates {(-1.130,0.000) (-2.600,0.000)};
\addplot[lineaE, flechaen=0.408] coordinates {(2.600,0.000) (1.130,0.000)};

  \fill[cargapos] (axis cs:-1,0) circle[radius=0.075];
  \fill[carganeg] (axis cs:1,0) circle[radius=0.075];
\end{axis}
\end{tikzpicture}
\caption{$+q$ y $-q$: $u_e$}
\end{subfigure}\hfill
\begin{subfigure}[b]{0.325\textwidth}\centering
\begin{tikzpicture}
\begin{axis}[mapa]
  \addplot graphics[xmin=-2.6, xmax=2.6, ymin=-1.8, ymax=1.8] {t5-ue-cruzado-\modotema};
  \draw[equipot] (axis cs:0,0) circle[radius=1];
  \fill[cargapos] (axis cs:-1,0) circle[radius=0.075];
  \fill[cargapos] (axis cs:1,0) circle[radius=0.075];
\end{axis}
\end{tikzpicture}
\caption{$+q$ y $+q$: $\eo\,\vE_1\cdot\vE_2$}
\end{subfigure}
\caption{La energía electrostática no cumple el principio de superposición. (a), (b) Densidad de energía
$u_e=\frac{\eo}{2}E^2$ (escala logarítmica, más brillante cuanto mayor) con las líneas de campo superpuestas: entre dos cargas
iguales el campo se anula en el punto medio y $u_e$ es pequeña; entre cargas opuestas los campos se refuerzan.
(c) Término cruzado $\eo\,\vE_1\cdot\vE_2$ para dos cargas iguales: es positivo (rojo) fuera de la circunferencia de
trazos, que tiene por diámetro el segmento entre las cargas, y negativo (azul) dentro. Su integral a todo el espacio
es exactamente la energía de interacción, $q^2/(4\pi\eo d)$ (\cref{ej:termino-cruzado}).}
\label{fig:superposicion}
\end{figure}


\begin{trampa}[Trampa habitual]
Si se duplica la carga de un sistema, su energía se multiplica por cuatro, no por dos; y la energía de dos sistemas juntos no es la suma
de sus energías por separado. Hay que calcular el campo total y elevarlo al cuadrado, o usar $\frac12\int\rho\phi\,\dd V$ con el potencial
total.
\end{trampa}

% =============================================================================
\section{Energía de un sistema de conductores}\label{sec:conductores}
% =============================================================================

Un conductor en equilibrio es una distribución superficial de carga con una peculiaridad: su superficie es equipotencial~\cite{T3}. Eso
simplifica enormemente la energía.

\begin{proposicion}[Energía de un sistema de conductores]\label{prop:conductores}
La energía electrostática de un sistema de $N$ conductores con cargas $Q_i$ y potenciales $\phi_i$ es
\begin{equation}\label{eq:conductores}
  \boxed{\;U_e=\frac12\sum_{i=1}^NQ_i\,\phi_i\;}
  =\frac12\sum_{i,j=1}^Np_{ij}\,Q_iQ_j=\frac12\sum_{i,j=1}^Nc_{ij}\,\phi_i\phi_j .
\end{equation}
\end{proposicion}

\begin{proof}
La carga reside en las superficies $S_i$ de los conductores, y sobre cada una el potencial es constante, $\phi=\phi_i$. Por la
\cref{prop:continua},
\[
  U_e=\frac12\sum_{i=1}^N\oiint_{S_i}\sigma_i\,\phi\,\dd S=\frac12\sum_{i=1}^N\phi_i\oiint_{S_i}\sigma_i\,\dd S=\frac12\sum_{i=1}^N\phi_i\,Q_i .
\]
Las otras dos formas se obtienen sustituyendo $\phi_i=\sum_jp_{ij}Q_j$ o $Q_i=\sum_jc_{ij}\phi_j$~\cite{T3}.
\end{proof}

\begin{observacion}[Lectura de la fórmula]\label{obs:conductores}
\begin{enumerate}
  \item Formalmente se parece a la energía de cargas puntuales \eqref{eq:puntuales}, pero el significado es otro: aquí $\phi_i$ es el
        potencial del conductor $i$ creado por \emph{todas} las cargas, incluida la suya, $Q_i$. Es energía \emph{total}, no solo de
        interacción, y por el \cref{cor:positiva} es positiva o nula.
  \item En consecuencia, las formas cuadráticas $\sum p_{ij}Q_iQ_j$ y $\sum c_{ij}\phi_i\phi_j$ son definidas positivas: se anulan solo si todas las
        cargas (o todos los potenciales) son nulos. Es otra forma de ver las propiedades de signo de $p_{ij}$ y $c_{ij}$ del Tema~3.
  \item Un conductor a tierra ($\phi_i=0$) o descargado ($Q_i=0$) no contribuye explícitamente a la suma, aunque sí a través de los
        potenciales o las cargas de los demás.
\end{enumerate}
\end{observacion}

\begin{corolario}[Energía de un condensador]\label{cor:condensador}
Un condensador de capacidad $C$, con cargas $\pm Q$ en sus armaduras y diferencia de potencial $V=\phi_1-\phi_2$, almacena la energía
\begin{equation}\label{eq:condensador}
  \boxed{\;U_e=\frac12\,QV=\frac12\,CV^2=\frac{Q^2}{2C}\;}.
\end{equation}
\end{corolario}

\begin{proof}
Con $Q_1=Q$ y $Q_2=-Q$, la \cref{prop:conductores} da $U_e=\frac12\big[Q\phi_1+(-Q)\phi_2\big]=\frac12Q(\phi_1-\phi_2)=\frac12QV$, y las otras dos formas se
siguen de $Q=CV$.
\end{proof}

\begin{ejemplo}[Esfera conductora aislada]\label{ej:esfera-aislada}
Para una esfera de radio $a$ y carga $Q$, $C=4\pi\eo a$ (Tema~3) y $U_e=Q^2/2C=Q^2/(8\pi\eo a)$, como en el \cref{ej:corteza}. Con $p_{11}=1/C$, es
también $\frac12p_{11}Q^2$.
\end{ejemplo}

\begin{ejemplo}[Condensador plano]\label{ej:plano}
Con $C=\eo S/d$, la energía es $U_e=\frac12CV^2=\frac12\eo SV^2/d$. Con el campo, $E=V/d$ uniforme en el volumen $Sd$ entre las placas y nulo
fuera (despreciando los bordes): $U_e=\frac{\eo}{2}\big(V/d\big)^2Sd$, idéntico. La energía está repartida uniformemente entre las placas.
\end{ejemplo}

% =============================================================================
\section{Fuerzas electrostáticas sobre conductores}\label{sec:fuerzas}
% =============================================================================

Dos conductores cargados se ejercen fuerzas. En principio podrían calcularse con la ley de Coulomb, integrando la fuerza sobre cada
elemento de superficie, o con la presión electrostática $\sigma^2/2\eo$ del Tema~3; pero la distribución de carga cambia al mover los
conductores y el cálculo directo es casi siempre intratable. La alternativa es la conservación de la energía.

\subsection{El teorema de Earnshaw}

Si un conductor está en reposo y sobre él actúa una fuerza electrostática neta no nula, algo tiene que equilibrarla. El siguiente
teorema muestra que ese algo no puede ser otra fuerza electrostática.

\begin{teorema}[Earnshaw, 1842]\label{teo:earnshaw}
Una carga puntual no puede estar en equilibrio estable bajo la acción exclusiva de fuerzas electrostáticas~\cite{Earnshaw}. En consecuencia, ningún
sistema de cargas puede mantenerse en equilibrio estable solo por interacciones coulombianas.
\end{teorema}

\begin{proof}
Sea $q>0$ una carga situada en $\vr_0$, en una región sin otras cargas, y sea $\phi$ el potencial creado por las demás. La energía potencial de
$q$ es $q\phi(\vr)$. Para que $\vr_0$ fuese una posición de equilibrio estable, $q\phi$ debería tener allí un mínimo local estricto: cualquier
desplazamiento debería aumentar la energía, de modo que la fuerza devolviera la carga a $\vr_0$. Ahora bien, en esa región $\phi$ satisface la
ecuación de Laplace, $\lap\phi=0$, porque las cargas que lo crean están en otro sitio. Si $\phi$ tuviese un mínimo estricto en $\vr_0$, en sus
proximidades sería
\[
  \frac{\partial^2\phi}{\partial x^2}\geq0,\qquad\frac{\partial^2\phi}{\partial y^2}\geq0,\qquad\frac{\partial^2\phi}{\partial z^2}\geq0,
\]
y como su suma es nula, las tres deberían anularse. Más rigurosamente, por la propiedad del valor medio de las funciones armónicas~\cite{T3,Kellogg},
$\phi(\vr_0)$ es el promedio de $\phi$ sobre cualquier esfera centrada en $\vr_0$ contenida en la región, y un promedio no puede ser estrictamente
menor que todos los valores promediados: $\phi$ no tiene mínimos ni máximos locales estrictos donde no hay carga. Para $q<0$ se razona igual con los
máximos. Como el razonamiento vale para cada carga del sistema, el sistema tampoco puede estar en equilibrio estable.
\end{proof}

\begin{observacion}[Puntos de silla]\label{obs:silla}
Sí puede haber puntos de \emph{equilibrio}, donde $\vE=\vcero$; pero son puntos de silla del potencial: estables en unas direcciones e inestables
en otras (\cref{fig:earnshaw}). Por eso, junto a las fuerzas electrostáticas sobre un conductor en reposo, siempre debe actuar una fuerza
mecánica (sujeciones, muelles, la cohesión del propio material) que las equilibre. Ya lo vimos en el Tema~3: la presión electrostática tiende a
arrancar la carga de la superficie, y lo impide la ligadura del metal.
\end{observacion}

% Figura: teorema de Earnshaw. Potencial de cuatro cargas positivas iguales en los
% vértices de un cuadrado (z = 0, vértices (+-1, +-1)), en el plano y = 0 cerca del centro:
% mínimo a lo largo de x, máximo a lo largo de z  ->  punto de silla.
\begin{figure}[t]
\centering
\begin{subfigure}[b]{0.56\textwidth}\centering
\begin{tikzpicture}
\begin{axis}[width=7.6cm, height=6.0cm, view={-32}{24}, xmin=-0.8, xmax=0.8, ymin=-0.8, ymax=0.8,
   xlabel={$x/a$}, ylabel={$z/a$}, zlabel={$\phi/\phi_0$}, xtick={-0.5,0,0.5}, ytick={-0.5,0,0.5},
   ztick={2.6,2.8,3.0}, colormap={tema}{color=(acentoB) color=(acentoA) color=(acentoC)},
   z buffer=sort, clip=false, label style={font=\footnotesize},
   x label style={sloped}, y label style={sloped}]
  \addplot3[surf, shader=faceted interp, faceted color=fondo!70!texto, line width=0.1pt,
     samples=29, domain=-0.8:0.8, y domain=-0.8:0.8, opacity=0.95]
     {1/sqrt((x-1)^2+1+y^2)+1/sqrt((x+1)^2+1+y^2)+1/sqrt((x-1)^2+1+y^2)+1/sqrt((x+1)^2+1+y^2)};
  \addplot3[only marks, mark=*, mark size=1.6pt, color=texto] coordinates {(0,0,2.8284)};
\end{axis}
\end{tikzpicture}
\caption{$\phi$ en el plano $y=0$ cerca del centro}
\end{subfigure}\hfill
\begin{subfigure}[b]{0.42\textwidth}\centering
\begin{tikzpicture}
\begin{axis}[width=6.0cm, height=5.0cm, xmin=-0.8, xmax=0.8, ymin=2.45, ymax=3.15,
   xlabel={desplazamiento$/a$}, ylabel={$\phi/\phi_0$}, grid=major, ytick={2.5,2.75,3.0}]
  \addplot[curvaE, domain=-0.8:0.8, samples=101]
     {2/sqrt((x-1)^2+1)+2/sqrt((x+1)^2+1)};
  \addplot[curvaD, domain=-0.8:0.8, samples=101] {4/sqrt(2+x^2)};
  \node[etiqf, font=\scriptsize, text=campoE] at (axis cs:0,3.02) {según $x$: mínimo};
  \node[etiqf, font=\scriptsize, text=campoD] at (axis cs:0,2.6) {según $z$: máximo};
  \addplot[only marks, mark=*, mark size=1.5pt, color=texto] coordinates {(0,2.8284)};
\end{axis}
\end{tikzpicture}
\caption{Cortes por el centro}
\end{subfigure}
\caption{Teorema de Earnshaw. Cuatro cargas positivas iguales en los vértices de un cuadrado de lado $2a$ en el
plano $z=0$, con $\phi_0=q/(4\pi\eo a)$. En el centro del cuadrado el campo es nulo, pero el potencial no tiene allí un mínimo:
crece al desplazarse según $x$ (hacia dos de las cargas) y decrece según $z$ (al alejarse del plano). Es un punto de
silla, como exige $\lap\phi=0$: una carga positiva colocada en el centro está en equilibrio, pero inestable frente a
desplazamientos perpendiculares al plano.}
\label{fig:earnshaw}
\end{figure}


\subsection{El método de los trabajos virtuales}

Imaginemos que uno de los conductores sufre un desplazamiento infinitesimal $\dl$ (o un giro $\dd\theta$), lento y reversible, mientras los
demás permanecen fijos (\cref{fig:desplazamiento}). La fuerza electrostática sobre él realiza el trabajo
\begin{equation}\label{eq:trabajo-mec}
  \dd\Wmec=\vFe\cdot\dl \qquad(\text{o } \dd\Wmec=\vtau_e\cdot\dd\vec\theta \text{ en un giro}),
\end{equation}
y el balance de energía depende de qué se mantiene constante durante el desplazamiento. Estudiamos los dos casos más sencillos.

% Figura: método de los trabajos virtuales. Desplazamiento infinitesimal de un
% conductor (a) aislado (cargas constantes) y (b) conectado a baterías (potenciales constantes).
\begin{figure}[t]
\centering
\tikzset{
  cond/.style={fill=conductor, draw=texto, line width=0.6pt},
  condv/.style={draw=texto, dashed, line width=0.5pt, dash pattern=on 2.5pt off 1.8pt},
  bateria/.pic={
    \draw[texto, line width=0.6pt] (-0.35,0) -- (-0.08,0) (0.08,0) -- (0.35,0);
    \draw[texto, line width=1.1pt] (-0.08,-0.22) -- (-0.08,0.22);
    \draw[texto, line width=0.6pt] (0.08,-0.13) -- (0.08,0.13);}
}
\begin{subfigure}[b]{0.47\textwidth}\centering
\begin{tikzpicture}[scale=0.88]
  \draw[cond] plot[smooth cycle, tension=0.9] coordinates {(-2.3,0.1) (-1.9,0.75) (-1.2,0.6) (-1.0,-0.1) (-1.6,-0.6) (-2.2,-0.45)};
  \node[font=\small, text=black] at (-1.65,0.05) {$Q_1$};
  \draw[cond] plot[smooth cycle, tension=0.9] coordinates {(0.6,0.25) (1.0,0.85) (1.6,0.7) (1.75,0.0) (1.3,-0.55) (0.75,-0.4)};
  \node[etiq, anchor=south] at (1.2,0.9) {$Q_2$};
  \draw[condv] plot[smooth cycle, tension=0.9] coordinates {(1.05,0.25) (1.45,0.85) (2.05,0.7) (2.2,0.0) (1.75,-0.55) (1.2,-0.4)};
  \draw[vecP] (1.75,0.0) -- (2.2,0.0);
  \node[etiq, text=campoP, anchor=west] at (2.25,0.0) {$\dl$};
  \node[etiqt, anchor=north] at (-0.3,-1.0) {cargas aisladas: $\dd Q_i=0$};
  \node[etiq, anchor=north] at (-0.3,-1.45) {$\dd W_{\mathrm{mec}}+\dd U_e=0$};
\end{tikzpicture}
\caption{Sistema aislado}
\end{subfigure}\hfill
\begin{subfigure}[b]{0.47\textwidth}\centering
\begin{tikzpicture}[scale=0.88]
  \draw[cond] plot[smooth cycle, tension=0.9] coordinates {(-2.3,0.1) (-1.9,0.75) (-1.2,0.6) (-1.0,-0.1) (-1.6,-0.6) (-2.2,-0.45)};
  \node[font=\small, text=black] at (-1.65,0.05) {$\phi_1$};
  \draw[cond] plot[smooth cycle, tension=0.9] coordinates {(0.6,0.25) (1.0,0.85) (1.6,0.7) (1.75,0.0) (1.3,-0.55) (0.75,-0.4)};
  \node[etiq, anchor=south] at (1.2,0.9) {$\phi_2$};
  \draw[condv] plot[smooth cycle, tension=0.9] coordinates {(1.05,0.25) (1.45,0.85) (2.05,0.7) (2.2,0.0) (1.75,-0.55) (1.2,-0.4)};
  \draw[vecP] (1.75,0.0) -- (2.2,0.0);
  % cables y baterías
  \draw[texto, line width=0.55pt] (-1.65,-0.55) -- (-1.65,-1.0) -- (-0.75,-1.0);
  \pic at (-0.4,-1.0) {bateria};
  \draw[texto, line width=0.55pt] (-0.05,-1.0) -- (0.25,-1.0);
  \draw[texto, line width=0.55pt] (1.15,-0.5) -- (1.15,-1.0) -- (0.55,-1.0);
  \pic at (0.4,-1.0) {bateria};
  \node[etiq, text=campoP, anchor=west] at (2.25,0.0) {$\dl$};
  \node[etiqt, anchor=north] at (-0.1,-1.35) {potenciales fijos: $\dd\phi_i=0$};
  \node[etiq, anchor=north] at (-0.1,-1.8) {$\dd W_{\mathrm{mec}}+\dd U_e=\dd W_{\mathrm{bat}}$};
\end{tikzpicture}
\caption{Sistema a potencial constante}
\end{subfigure}
\caption{Método de los trabajos virtuales. Se imagina un desplazamiento infinitesimal $\dl$ de uno de los conductores
(trazos) y se hace el balance de energía. (a) Si el sistema está aislado, las cargas no cambian y el trabajo mecánico
$\dd W_{\mathrm{mec}}=\vF_e\cdot\dl$ sale de la energía electrostática. (b) Si los conductores están unidos a baterías, sus potenciales
no cambian, pero las baterías deben suministrar las cargas $\dd Q_i$, con un trabajo $\dd W_{\mathrm{bat}}=\sum_i\phi_i\,\dd Q_i$.}
\label{fig:desplazamiento}
\end{figure}


\begin{teorema}[Fuerza a carga constante]\label{teo:fuerza-Q}
Si el sistema de conductores está aislado, de modo que las cargas $Q_i$ permanecen constantes, la fuerza electrostática y el momento sobre un
conductor son
\begin{equation}\label{eq:fuerza-Q}
  \boxed{\;\vFe=-\grad U_e\big|_{Q}\;},\qquad F_\xi=-\frac{\partial U_e}{\partial\xi}\bigg|_{Q},\qquad \tau_\theta=-\frac{\partial U_e}{\partial\theta}\bigg|_{Q},
\end{equation}
donde el gradiente se toma respecto a las coordenadas que describen la posición del conductor. En función de los coeficientes de potencial,
\begin{equation}\label{eq:fuerza-pij}
  F_\xi=-\frac12\sum_{i,j}Q_iQ_j\,\frac{\partial p_{ij}}{\partial\xi} .
\end{equation}
\end{teorema}

\begin{proof}
Al estar aislado, no entra energía de fuera: el trabajo mecánico que realiza la fuerza electrostática se obtiene a costa de la energía
almacenada,
\[
  \dd\Wmec+\dd U_e=0\quad\Longrightarrow\quad \vFe\cdot\dl=-\dd U_e\big|_Q .
\]
Como $U_e$ es función de la configuración, $\dd U_e=\grad U_e\cdot\dl$ y, al ser $\dl$ arbitrario, $\vFe=-\grad U_e|_Q$. Para un giro, $\tau_\theta\,\dd\theta=-\dd U_e|_Q$. Por
último, con $U_e=\frac12\sum p_{ij}Q_iQ_j$ y las $Q_i$ constantes, solo dependen de la posición los coeficientes $p_{ij}$, que son puramente geométricos, y se obtiene
\eqref{eq:fuerza-pij}.
\end{proof}

\begin{teorema}[Fuerza a potencial constante]\label{teo:fuerza-phi}
Si cada conductor está unido a una batería que mantiene fijo su potencial $\phi_i$, la fuerza y el momento sobre un conductor son
\begin{equation}\label{eq:fuerza-phi}
  \boxed{\;\vFe=+\grad U_e\big|_{\phi}\;},\qquad F_\xi=+\frac{\partial U_e}{\partial\xi}\bigg|_{\phi},\qquad \tau_\theta=+\frac{\partial U_e}{\partial\theta}\bigg|_{\phi},
  \qquad F_\xi=\frac12\sum_{i,j}\phi_i\phi_j\,\frac{\partial c_{ij}}{\partial\xi}.
\end{equation}
\end{teorema}

\begin{proof}
Ahora el sistema no está aislado. Para mantener los potenciales al desplazar el conductor, las baterías deben suministrar cargas $\dd Q_i$, trayéndolas
desde un potencial nulo hasta $\phi_i$; por el \cref{lem:trabajo}, su trabajo es
\[
  \dd\Wbat=\sum_i\phi_i\,\dd Q_i .
\]
La variación de energía electrostática, con los $\phi_i$ fijos, es $\dd U_e=\dd\big(\frac12\sum_i\phi_iQ_i\big)=\frac12\sum_i\phi_i\,\dd Q_i=\frac12\dd\Wbat$. El balance de energía es
\[
  \begin{aligned}
  \dd\Wmec+\dd U_e=\dd\Wbat\quad\Longrightarrow\quad \dd\Wmec&=\dd\Wbat-\dd U_e=\sum_i\phi_i\,\dd Q_i-\frac12\sum_i\phi_i\,\dd Q_i\\
  &=\frac12\sum_i\phi_i\,\dd Q_i=\dd U_e\big|_\phi .
  \end{aligned}
\]
Así $\vFe\cdot\dl=+\dd U_e|_\phi$ y $\vFe=+\grad U_e|_\phi$. La última forma se obtiene de $U_e=\frac12\sum c_{ij}\phi_i\phi_j$ con los $\phi_i$ constantes.
\end{proof}

\begin{observacion}[El reparto de la energía de las baterías]\label{obs:baterias}
A potencial constante, de la energía que aportan las baterías, la mitad se convierte en trabajo mecánico y la otra mitad aumenta la energía
almacenada en el campo. Por eso el signo cambia: el sistema evoluciona hacia estados de \emph{mayor} energía electrostática, porque las baterías
pagan el doble de lo que se almacena. A carga constante, en cambio, el sistema tiende al mínimo de $U_e$.
\end{observacion}

\subsection{Las dos fórmulas dan la misma fuerza}

El signo opuesto de \eqref{eq:fuerza-Q} y \eqref{eq:fuerza-phi} podría hacer pensar que la fuerza depende del proceso. No es así: la fuerza sobre un
conductor depende solo de la configuración en un instante, de las cargas y de los campos que hay en ese momento, no de lo que vaya a ocurrir en el
desplazamiento imaginario.

\begin{proposicion}[Equivalencia de ambos métodos]\label{prop:equivalencia}
En una misma configuración, con cargas $Q_i$ y potenciales $\phi_i$ relacionados por $\boldsymbol\phi=\mathsf P\boldsymbol Q$,
\begin{equation}\label{eq:equivalencia}
  -\frac{\partial U_e}{\partial\xi}\bigg|_{Q}=+\frac{\partial U_e}{\partial\xi}\bigg|_{\phi} .
\end{equation}
\end{proposicion}

\begin{proof}
Sea $\mathsf P(\xi)$ la matriz de coeficientes de potencial y $\mathsf C(\xi)=\mathsf P(\xi)^{-1}$ la de capacidad. Derivando $\mathsf C\mathsf P=\mathsf 1$ respecto a $\xi$ se obtiene
$\mathsf C'\mathsf P+\mathsf C\mathsf P'=\mathsf 0$, es decir, $\mathsf C'=-\mathsf C\mathsf P'\mathsf C$. Entonces, con $\boldsymbol Q=\mathsf C\boldsymbol\phi$ y la simetría de $\mathsf C$,
\[
  \frac{\partial U_e}{\partial\xi}\bigg|_{\phi}=\frac12\,\boldsymbol\phi^{\mathsf T}\mathsf C'\boldsymbol\phi=-\frac12\,\boldsymbol\phi^{\mathsf T}\mathsf C\mathsf P'\mathsf C\boldsymbol\phi
  =-\frac12\,\boldsymbol Q^{\mathsf T}\mathsf P'\boldsymbol Q=-\frac{\partial U_e}{\partial\xi}\bigg|_{Q}. \qedhere
\]
\end{proof}

\begin{corolario}[Fuerza sobre una armadura de condensador]\label{cor:fuerza-condensador}
Si la capacidad de un condensador depende de una coordenada $\xi$, la fuerza generalizada sobre la armadura móvil es
\begin{equation}\label{eq:fuerza-C}
  \boxed{\;F_\xi=\frac12\,V^2\,\frac{\dd C}{\dd\xi}=\frac{Q^2}{2C^2}\,\frac{\dd C}{\dd\xi}\;},
\end{equation}
y el momento, $\tau_\theta=\frac12V^2\,\dd C/\dd\theta$. El sistema tiende a aumentar su capacidad.
\end{corolario}

\begin{proof}
A carga constante, $F_\xi=-\dfrac{\dd}{\dd\xi}\Big(\dfrac{Q^2}{2C}\Big)=\dfrac{Q^2}{2C^2}\dfrac{\dd C}{\dd\xi}$; a potencial constante, $F_\xi=\dfrac{\dd}{\dd\xi}\Big(\dfrac12CV^2\Big)=\dfrac12V^2\dfrac{\dd C}{\dd\xi}$.
Ambas coinciden porque $Q=CV$.
\end{proof}

\begin{metodo}{Receta: fuerza por trabajos virtuales}
\begin{pasos}
  \item Elegir la coordenada $\xi$ que describe el movimiento posible (distancia, longitud introducida, ángulo).
  \item Expresar la energía en función de $\xi$: $U_e=Q^2/2C(\xi)$ o $\frac12C(\xi)V^2$ para un condensador; en general $\frac{\eo}{2}\int E^2\dd V$ o
        $\frac12\sum p_{ij}Q_iQ_j$.
  \item Derivar con el signo que corresponda: $-\partial_\xi U_e$ a carga constante, $+\partial_\xi U_e$ a potencial constante.
  \item Interpretar el signo: $F_\xi>0$ significa que la fuerza tiende a aumentar $\xi$.
\end{pasos}
\end{metodo}

\begin{ejemplo}[Comprobación con la presión electrostática]\label{ej:presion}
En un condensador plano de área $S$ con cargas $\pm Q$ y separación $\xi$, $C=\eo S/\xi$ y $F_\xi=\frac{Q^2}{2C^2}\frac{\dd C}{\dd\xi}=-\frac{Q^2}{2\eo S}$: las placas se atraen.
Con la presión del Tema~3, cada placa sufre $P=\sigma^2/2\eo$ dirigida hacia el exterior del metal, es decir, hacia la otra placa:
$F=\frac{\sigma^2}{2\eo}S=\frac{Q^2}{2\eo S}$. Coinciden (\cref{prob:fuerza-placas}).
\end{ejemplo}

\begin{ejemplo}[El acelerómetro capacitivo]\label{ej:acelerometro}
Los acelerómetros de los teléfonos móviles son condensadores diferenciales microfabricados: una placa unida a una masa sísmica se mueve
entre dos placas fijas, y la aceleración se mide por el desequilibrio $C_1-C_2$ (\cref{fig:acelerometro}). Las fuerzas electrostáticas sobre la
placa móvil se calculan exactamente como en el \cref{prob:tres-placas}; y, por el teorema de Earnshaw, son los muelles los que proporcionan un
equilibrio estable.
\end{ejemplo}

% Figura: acelerómetro capacitivo (condensador diferencial). Masa sísmica móvil
% entre dos placas fijas; capacidades C1 = eps0 S/(d - x), C2 = eps0 S/(d + x).
\begin{figure}[t]
\centering
\begin{subfigure}[b]{0.5\textwidth}\centering
\begin{tikzpicture}[scale=1.0]
  \tikzset{muelle/.style={decorate, decoration={coil, aspect=0.45, segment length=2.4mm, amplitude=1.6mm},
                          draw=texto!80!fondo, line width=0.55pt}}
  % bastidor
  \fill[conductor!55!fondo] (-2.9,-1.35) rectangle (-2.7,1.35);
  \fill[conductor!55!fondo] (2.7,-1.35) rectangle (2.9,1.35);
  \draw[texto, line width=0.5pt] (-2.9,-1.35) rectangle (-2.7,1.35) (2.7,-1.35) rectangle (2.9,1.35);
  % masa sísmica (placa central móvil) y muelles
  \fill[conductor] (-0.45,-0.55) rectangle (0.45,0.55);
  \draw[texto, line width=0.6pt] (-0.45,-0.55) rectangle (0.45,0.55);
  \node[font=\small, text=black] at (0,0) {$m$};
  \draw[muelle] (-2.7,0) -- (-0.45,0);
  \draw[muelle] (0.45,0) -- (2.7,0);
  % placas fijas del condensador diferencial (debajo)
  \draw[placa] (-1.0,-1.25) -- (-1.0,-2.25);
  \draw[placa] (1.0,-1.25) -- (1.0,-2.25);
  \draw[placa, draw=acentoC] (0,-0.55) -- (0,-2.25);
  \node[etiqt, anchor=north] at (-1.0,-2.3) {fija};
  \node[etiqt, anchor=north] at (1.0,-2.3) {fija};
  \node[etiq, text=acentoC, anchor=north] at (0,-2.3) {móvil};
  \node[etiq] at (-0.5,-1.5) {$C_1$};
  \node[etiq] at (0.5,-1.5) {$C_2$};
  \draw[cota] (-1.0,-1.95) -- (0,-1.95) node[midway, below, etiq, font=\footnotesize] {$d$};
  \draw[cota] (0,-1.95) -- (1.0,-1.95) node[midway, below, etiq, font=\footnotesize] {$d$};
  % aceleración
  \draw[vecE] (-0.6,1.0) -- (0.6,1.0);
  \node[etiq, text=campoE, anchor=south] at (0,1.05) {$\vec a$};
\end{tikzpicture}
\caption{Esquema}
\end{subfigure}\hfill
\begin{subfigure}[b]{0.46\textwidth}\centering
\begin{tikzpicture}
\begin{axis}[width=6.4cm, height=5.0cm, xmin=-0.6, xmax=0.6, ymin=-2.0, ymax=2.6, grid=major,
   xlabel={$x/d$}, ylabel={$C\,d/(\eo S)$}, xtick={-0.5,0,0.5}, ytick={-2,-1,0,1,2},
   legend style={at={(0.5,1.03)}, anchor=south, legend columns=3, font=\scriptsize}]
  \addplot[curvaE, domain=-0.6:0.6, samples=80] {1/(1-x)}; \addlegendentry{$C_1$}
  \addplot[curvaD, domain=-0.6:0.6, samples=80] {1/(1+x)}; \addlegendentry{$C_2$}
  \addplot[curvaP, domain=-0.6:0.6, samples=80] {1/(1-x)-1/(1+x)}; \addlegendentry{$C_1-C_2$}
\end{axis}
\end{tikzpicture}
\caption{Capacidades frente al desplazamiento}
\end{subfigure}
\caption{Acelerómetro capacitivo (sistema MEMS). Una masa sísmica sujeta por muelles lleva una placa que se
mueve entre dos placas fijas; en reposo, a distancia $d$ de ambas. Una aceleración la desplaza $x$ y desequilibra el
condensador diferencial: $C_1-C_2=\frac{\eo S}{d-x}-\frac{\eo S}{d+x}\simeq\frac{2\eo S}{d^2}\,x$ para $x\ll d$. Las fuerzas electrostáticas
entre las placas, que se calculan como en el \cref{prob:tres-placas}, deben ser pequeñas frente a las de los muelles,
porque el equilibrio electrostático puro es inestable (\cref{teo:earnshaw}).}
\label{fig:acelerometro}
\end{figure}


% =============================================================================
\section{Seminarios prácticos}\label{sec:seminarios}
% =============================================================================

Esta sección resuelve con detalle los problemas de la relación del tema~\cite{DFA-problemas} que corresponden a la primera parte: la energía en el
vacío y en sistemas de conductores, y las fuerzas sobre conductores y distribuciones de carga. Los problemas 5.1(b), 5.6 y 5.10 involucran la
energía en presencia de dieléctricos y se resolverán en la segunda parte. Cada problema sigue la estructura habitual: enunciado, idea, resolución
paso a paso y comentario. La guía de métodos~\cite{Guia-T5} los recoge en formato de ficha.

\subsection{Cómo plantear un problema de energía y fuerzas}\label{sec:plantear}

\begin{table}[H]
\centering
\caption{Mapa de decisión: qué método conviene según lo que pide el enunciado.}\label{tab:decision}
\small
\begin{tabularx}{\textwidth}{@{}>{\raggedright\arraybackslash}X>{\raggedright\arraybackslash}p{5.0cm}l@{}}
\toprule
Lo que pide el enunciado & Método & Problemas\\
\midrule
Energía de cargas puntuales o de una distribución con $\phi$ sencillo & M1. $\frac12\sum q_i\phi_i$,\ \ $\frac12\int\rho\phi\,\dd V$ & 5.1\\
Energía con simetría, o en una región del espacio & M2. $\frac{\eo}{2}\int E^2\dd V$ & 5.1, 5.2\\
Energía de conductores con cargas o potenciales dados & M3. $\frac12\sum Q_i\phi_i$, coeficientes $p_{ij}$ & 5.3\\
Fuerza o par sobre un conductor & M4. Trabajos virtuales: $\mp\partial_\xi U_e$ & 5.4, 5.5, 5.9\\
Fuerza sobre una distribución dada en un campo externo & M5. $\int\dd q\,\vE_{\mathrm{ext}}$ o $-\grad\Uint$ & 5.7, 5.8\\
\bottomrule
\end{tabularx}
\end{table}

\begin{metodo}{Los cinco métodos}
\begin{description}[leftmargin=0pt, itemsep=0.35em, font=\normalfont\bfseries\color{acentoA}]
  \item[M1. Cargas y potenciales.] Calcular el potencial total en cada carga (o en cada punto de la distribución) y aplicar
        $U_e=\frac12\sum_iq_i\phi_i$ o $\frac12\int\rho\phi\,\dd V$. Si hay simetría, un solo potencial basta.
  \item[M2. Campo.] Calcular $\vE$ en todo el espacio (Gauss, si hay simetría) e integrar $u_e=\frac{\eo}{2}E^2$ con el elemento de volumen adecuado:
        $4\pi r^2\dd r$ (esférica), $2\pi\rho\ell\,\dd\rho$ (cilíndrica), $S\,\dd z$ (plana), o $r^2\sen\theta\,\dd r\,\dd\theta\,\dd\varphi$ en general.
  \item[M3. Conductores.] $U_e=\frac12\sum_iQ_i\phi_i$ con los potenciales totales; con esferas lejanas, $p_{ii}\simeq k/R_i$ y $p_{ij}\simeq k/d_{ij}$ (Tema~3).
  \item[M4. Trabajos virtuales.] Escribir $U_e(\xi)$ y derivar: $F_\xi=-\partial_\xi U_e|_Q$ (aislado) o $+\partial_\xi U_e|_\phi$ (baterías); para un condensador,
        $F_\xi=\frac12V^2\,\dd C/\dd\xi$.
  \item[M5. Distribución en un campo externo.] Si el campo externo $\vE_{\mathrm{ext}}$ es conocido y no cambia al mover el cuerpo, $\vF=\int\vE_{\mathrm{ext}}\,\dd q$; o bien
        $\vF=-\grad\Uint$ con $\Uint=\int\phi_{\mathrm{ext}}\,\dd q$.
\end{description}
\end{metodo}

% -----------------------------------------------------------------------------
\subsection{Problema 5.1(a). Esfera uniformemente cargada}\label{sec:p51}

\begin{problema}[Esfera uniformemente cargada]\label{prob:esfera-cargada}
Calcula la energía electrostática asociada a una distribución uniforme de carga $Q$ en una esfera de radio $a$ cuando (a) tanto la esfera como el
resto del espacio son el vacío.
\end{problema}

\begin{resolucion}
\idea Simetría esférica: el campo sale de Gauss. Lo resolvemos con los dos métodos, M1 y M2, que deben coincidir.

\paso{Paso 1: campo} Con una esfera gaussiana de radio $r$, la carga encerrada es $Q(r/a)^3$ si $r<a$ y $Q$ si $r>a$:
\[
  \vE=\frac{Q\,r}{4\pi\eo a^3}\,\uv r\quad(r<a),\qquad \vE=\frac{Q}{4\pi\eo r^2}\,\uv r\quad(r>a).
\]
\paso{Paso 2: potencial (para M1)} Fuera, $\phi=Q/(4\pi\eo r)$. Dentro, integrando desde el infinito por los dos tramos,
\[
  \begin{aligned}
  \phi(r)&=-\int_\infty^a\frac{Q}{4\pi\eo r'^2}\,\dd r'-\int_a^r\frac{Q\,r'}{4\pi\eo a^3}\,\dd r'\\
         &=\frac{Q}{4\pi\eo a}+\frac{Q}{8\pi\eo a^3}\big(a^2-r^2\big)=\frac{Q}{8\pi\eo a^3}\big(3a^2-r^2\big).
  \end{aligned}
\]
\paso{Método 1: $\frac12\int\rho\phi\,\dd V$} Con $\rho=\dfrac{Q}{\frac43\pi a^3}=\dfrac{3Q}{4\pi a^3}$ y $\dd V=4\pi r^2\dd r$,
\[
  \begin{aligned}
  U_e&=\frac12\int_0^a\frac{3Q}{4\pi a^3}\,\frac{Q}{8\pi\eo a^3}\big(3a^2-r^2\big)\,4\pi r^2\,\dd r\\
     &=\frac{3Q^2}{16\pi\eo a^6}\int_0^a\big(3a^2r^2-r^4\big)\dd r=\frac{3Q^2}{16\pi\eo a^6}\Big(a^5-\frac{a^5}{5}\Big).
  \end{aligned}
\]
\paso{Método 2: $\frac{\eo}{2}\int E^2\dd V$} Hay que integrar en todo el espacio, dentro y fuera:
\[
  \begin{aligned}
  U_e&=\frac{\eo}{2}\Big(\frac{Q}{4\pi\eo}\Big)^2 4\pi\left[\frac{1}{a^6}\int_0^a r^4\,\dd r+\int_a^\infty\frac{\dd r}{r^2}\right]\\
     &=\frac{Q^2}{8\pi\eo}\Big(\frac{1}{5a}+\frac1a\Big)=\frac{Q^2}{8\pi\eo a}\cdot\frac65 .
  \end{aligned}
\]
Ambos dan
\begin{equation}\label{eq:p51}
  \recuadro{U_e=\frac{3Q^2}{20\pi\eo a}=\frac35\,\frac{Q^2}{4\pi\eo a}.}
\end{equation}
\end{resolucion}

\begin{observacion}[Interpretación]\label{obs:p51}
(a) El método 2 muestra el reparto: $\frac16$ de la energía está dentro de la esfera y $\frac56$ fuera (\cref{fig:p51}). (b) La esfera maciza almacena $\frac65$ de la
energía de una corteza con la misma carga y radio, $Q^2/8\pi\eo a$: tiene además el campo interior. (c) Es la energía de formación de un núcleo atómico
modelado como esfera uniforme de protones; de ella procede el término coulombiano de la fórmula semiempírica de masas~\cite{Griffiths}.
(d) Trampa: en el método 1 basta integrar donde hay carga ($r<a$); en el método 2 hay que integrar donde hay campo, es decir, en todo el espacio.
\end{observacion}

% Figura: problema 5.1(a). Esfera uniformemente cargada: campo, potencial y densidad
% de energía frente a r (unidades Q/(4 pi eps0 a^2), Q/(4 pi eps0 a) y valor en r = a).
\begin{figure}[t]
\centering
\begin{subfigure}[b]{0.28\textwidth}\centering
\begin{tikzpicture}[scale=1.0]
  \shade[inner color=acentoB!45!fondo, outer color=acentoB!18!fondo] (0,0) circle (1.05);
  \draw[texto, line width=0.5pt] (0,0) circle (1.05);
  \draw[gauss] (0,0) circle (0.55);
  \node[etiq, text=acentoB] at (-45:0.8) {$r$};
  \draw[texto, line width=0.4pt, -{Stealth[length=1.5mm]}] (0,0) -- (40:1.05);
  \node[etiq] at (40:1.32) {$a$};
  \node[etiq] at (0,-1.55) {$\rho=\dfrac{3Q}{4\pi a^{3}}$};
\end{tikzpicture}
\caption{Geometría}
\end{subfigure}\hfill
\begin{subfigure}[b]{0.34\textwidth}\centering
\begin{tikzpicture}
\begin{axis}[width=5.2cm, height=4.5cm, axis lines=left, xmin=0, xmax=3.5, ymin=0, ymax=1.6,
   xlabel={$r/a$}, xtick={0,1,2,3}, ytick={0,0.5,1,1.5}, grid=major,
   legend style={at={(0.5,1.03)}, anchor=south, legend columns=2, font=\scriptsize}]
  \addplot[curvaE, domain=0:1, samples=2] {x}; \addlegendentry{$E$}
  \addplot[curvaV, domain=0:1, samples=40] {(3-x^2)/2}; \addlegendentry{$\phi$}
  \addplot[curvaE, domain=1:3.5, samples=60, forget plot] {1/x^2};
  \addplot[curvaV, domain=1:3.5, samples=60, forget plot] {1/x};
\end{axis}
\end{tikzpicture}
\caption{$E$ y $\phi$}
\end{subfigure}\hfill
\begin{subfigure}[b]{0.34\textwidth}\centering
\begin{tikzpicture}
\begin{axis}[width=5.2cm, height=4.5cm, axis lines=left, xmin=0, xmax=3.5, ymin=0, ymax=1.15,
   xlabel={$r/a$}, ylabel={$\dd U_e/\dd r$}, xtick={0,1,2,3}, ytick={0,0.5,1}, grid=major]
  \addplot[name path=f1, curvaD, domain=0:1, samples=40] {x^4};
  \addplot[name path=f2, curvaD, domain=1:3.5, samples=60] {1/x^2};
  \addplot[name path=ej, draw=none, domain=0:3.5] {0};
  \addplot[fill=campoD, fill opacity=0.18, draw=none] fill between[of=f1 and ej, soft clip={domain=0:1}];
  \addplot[fill=campoE, fill opacity=0.18, draw=none] fill between[of=f2 and ej, soft clip={domain=1:3.5}];
  \node[etiqt] at (axis cs:0.72,0.12) {$\tfrac16$};
  \node[etiqt] at (axis cs:1.75,0.15) {$\tfrac56$};
\end{axis}
\end{tikzpicture}
\caption{Energía por unidad de radio}
\end{subfigure}
\caption{\Cref{prob:esfera-cargada}. (b) Campo (unidades $Q/4\pi\eo a^2$) y potencial (unidades $Q/4\pi\eo a$): el campo crece linealmente
dentro y decae como $1/r^2$ fuera; el potencial es parabólico dentro, con $\phi(0)=\frac32\phi(a)$. (c) Energía por unidad de radio,
$\dd U_e/\dd r=4\pi r^2u_e$: crece como $r^4$ dentro y decae como $r^{-2}$ fuera. Las áreas sombreadas dan el reparto de la energía:
una sexta parte dentro de la esfera y cinco sextas partes fuera.}
\label{fig:p51}
\end{figure}


% -----------------------------------------------------------------------------
\subsection{Problema 5.2. Energía del campo de un dipolo fuera de una esfera}\label{sec:p52}

\begin{problema}[Dipolo en el centro de una esfera]\label{prob:dipolo}
Calcula la energía electrostática almacenada en la región exterior a una esfera de radio $a$ en cuyo centro hay un dipolo de momento $\vp$.
\end{problema}

\begin{resolucion}
\idea Se pide la energía de una \emph{región}, no de todo el sistema: solo tiene sentido con la densidad de energía (método M2).

\paso{Paso 1: campo del dipolo} Con $\vp$ según el eje polar~\cite{T2},
\[
  \vE=\frac{p}{4\pi\eo r^3}\big(2\cos\theta\,\uv r+\sen\theta\,\uv\theta\big),\qquad
  E^2=\Big(\frac{p}{4\pi\eo r^3}\Big)^2\big(3\cos^2\theta+1\big),
\]
pues $4\cos^2\theta+\sen^2\theta=3\cos^2\theta+1$.
\paso{Paso 2: integral en $r>a$} Con $\dd V=r^2\sen\theta\,\dd r\,\dd\theta\,\dd\varphi$ y $u=\cos\theta$,
\[
  \begin{aligned}
  U_e&=\frac{\eo}{2}\Big(\frac{p}{4\pi\eo}\Big)^2\int_0^{2\pi}\!\dd\varphi\int_{-1}^{1}\big(3u^2+1\big)\dd u\int_a^\infty\frac{r^2\,\dd r}{r^6}\\
     &=\frac{p^2}{32\pi^2\eo}\cdot2\pi\cdot\Big[u^3+u\Big]_{-1}^{1}\cdot\Big[-\frac{1}{3r^3}\Big]_a^\infty .
  \end{aligned}
\]
Como $\big[u^3+u\big]_{-1}^{1}=2-(-2)=4$ y la integral radial vale $1/(3a^3)$,
\begin{equation}\label{eq:p52}
  \recuadro{U_e(r>a)=\frac{p^2}{32\pi^2\eo}\cdot2\pi\cdot4\cdot\frac{1}{3a^3}=\frac{p^2}{12\pi\eo a^3}.}
\end{equation}
\end{resolucion}

\begin{observacion}[Divergencia y anisotropía]\label{obs:p52}
Si $a\to0$, la energía diverge como $a^{-3}$: la energía propia de un dipolo puntual es infinita, como la de una carga puntual. Por eso no se puede
calcular con $\frac12\int\rho\phi$, que exigiría la estructura interna del dipolo. La densidad de energía es cuatro veces mayor en el eje que en el plano
ecuatorial (\cref{fig:p52}).
\end{observacion}

% Figura: problema 5.2. Dipolo puntual en el centro de una esfera de radio a:
% líneas de campo exteriores (exactas) y distribución angular de la densidad de energía.
\begin{figure}[t]
\centering
\begin{subfigure}[b]{0.52\textwidth}\centering
\begin{tikzpicture}
\begin{axis}[marco, width=7.2cm, xmin=-2.6, xmax=2.6, ymin=-2.0, ymax=2.0]
  \fill[dielec] (axis cs:0,0) circle[radius=1];
  %% Generado por scripts/tema5.py -- no editar a mano
\addplot[lineaE, flechaen=0.396] coordinates {(2.600,1.454) (2.571,1.426) (2.542,1.398) (2.513,1.370) (2.484,1.342) (2.455,1.314) (2.425,1.286) (2.395,1.259) (2.365,1.232) (2.335,1.205) (2.305,1.178) (2.274,1.151) (2.244,1.125) (2.213,1.098) (2.182,1.072) (2.151,1.047) (2.120,1.021) (2.089,0.995) (2.057,0.970) (2.025,0.945) (1.993,0.920) (1.961,0.895) (1.929,0.871) (1.897,0.847) (1.864,0.823) (1.831,0.799) (1.799,0.775) (1.766,0.752) (1.732,0.729) (1.699,0.706) (1.666,0.683) (1.632,0.661) (1.598,0.639) (1.564,0.617) (1.530,0.595) (1.496,0.573) (1.461,0.552) (1.427,0.531) (1.392,0.511) (1.357,0.490) (1.322,0.470) (1.287,0.450) (1.251,0.431) (1.216,0.411) (1.180,0.392) (1.144,0.374) (1.108,0.355) (1.072,0.337) (1.036,0.319) (0.999,0.302) (0.963,0.285)};
\addplot[lineaE, flechaen=0.604] coordinates {(-0.963,0.285) (-0.999,0.302) (-1.036,0.319) (-1.072,0.337) (-1.108,0.355) (-1.144,0.374) (-1.180,0.392) (-1.216,0.411) (-1.251,0.431) (-1.287,0.450) (-1.322,0.470) (-1.357,0.490) (-1.392,0.511) (-1.427,0.531) (-1.461,0.552) (-1.496,0.573) (-1.530,0.595) (-1.564,0.617) (-1.598,0.639) (-1.632,0.661) (-1.666,0.683) (-1.699,0.706) (-1.732,0.729) (-1.766,0.752) (-1.799,0.775) (-1.831,0.799) (-1.864,0.823) (-1.897,0.847) (-1.929,0.871) (-1.961,0.895) (-1.993,0.920) (-2.025,0.945) (-2.057,0.970) (-2.089,0.995) (-2.120,1.021) (-2.151,1.047) (-2.182,1.072) (-2.213,1.098) (-2.244,1.125) (-2.274,1.151) (-2.305,1.178) (-2.335,1.205) (-2.365,1.232) (-2.395,1.259) (-2.425,1.286) (-2.455,1.314) (-2.484,1.342) (-2.513,1.370) (-2.542,1.398) (-2.571,1.426) (-2.600,1.454)};
\addplot[lineaE, flechaen=0.074] coordinates {(2.052,2.000) (2.038,1.962) (2.024,1.924) (2.009,1.887) (1.994,1.849) (1.979,1.812) (1.963,1.775) (1.946,1.738) (1.930,1.701) (1.913,1.665) (1.895,1.628) (1.877,1.592) (1.859,1.556) (1.840,1.520) (1.821,1.485) (1.801,1.450) (1.781,1.415) (1.761,1.380) (1.740,1.345) (1.719,1.311) (1.698,1.276) (1.676,1.242) (1.654,1.209) (1.631,1.175) (1.608,1.142) (1.585,1.109) (1.561,1.077) (1.537,1.044) (1.512,1.012) (1.487,0.980) (1.462,0.949) (1.437,0.918) (1.411,0.887) (1.384,0.856) (1.358,0.826) (1.330,0.796) (1.303,0.766) (1.275,0.737) (1.247,0.708) (1.219,0.680) (1.190,0.651) (1.161,0.623) (1.131,0.596) (1.101,0.569) (1.071,0.542) (1.040,0.516) (1.009,0.490) (0.978,0.464) (0.946,0.439) (0.914,0.415)};
\addplot[lineaE, flechaen=0.926] coordinates {(-0.914,0.415) (-0.946,0.439) (-0.978,0.464) (-1.009,0.490) (-1.040,0.516) (-1.071,0.542) (-1.101,0.569) (-1.131,0.596) (-1.161,0.623) (-1.190,0.651) (-1.219,0.680) (-1.247,0.708) (-1.275,0.737) (-1.303,0.766) (-1.330,0.796) (-1.358,0.826) (-1.384,0.856) (-1.411,0.887) (-1.437,0.918) (-1.462,0.949) (-1.487,0.980) (-1.512,1.012) (-1.537,1.044) (-1.561,1.077) (-1.585,1.109) (-1.608,1.142) (-1.631,1.175) (-1.654,1.209) (-1.676,1.242) (-1.698,1.276) (-1.719,1.311) (-1.740,1.345) (-1.761,1.380) (-1.781,1.415) (-1.801,1.450) (-1.821,1.485) (-1.840,1.520) (-1.859,1.556) (-1.877,1.592) (-1.895,1.628) (-1.913,1.665) (-1.930,1.701) (-1.946,1.738) (-1.963,1.775) (-1.979,1.812) (-1.994,1.849) (-2.009,1.887) (-2.024,1.924) (-2.038,1.962) (-2.052,2.000)};
\addplot[lineaE, flechaen=0.500] coordinates {(1.374,2.000) (1.375,1.959) (1.374,1.918) (1.374,1.877) (1.372,1.837) (1.369,1.796) (1.366,1.755) (1.362,1.714) (1.357,1.674) (1.351,1.633) (1.345,1.593) (1.338,1.553) (1.330,1.513) (1.321,1.473) (1.312,1.433) (1.302,1.393) (1.291,1.354) (1.279,1.315) (1.267,1.276) (1.254,1.237) (1.240,1.199) (1.226,1.160) (1.211,1.122) (1.195,1.085) (1.178,1.047) (1.161,1.010) (1.143,0.974) (1.125,0.937) (1.105,0.901) (1.085,0.865) (1.065,0.830) (1.044,0.795) (1.022,0.761) (0.999,0.727) (0.976,0.693) (0.952,0.660) (0.928,0.627) (0.903,0.595) (0.877,0.563) (0.851,0.532)};
\addplot[lineaE, flechaen=0.500] coordinates {(-0.851,0.532) (-0.877,0.563) (-0.903,0.595) (-0.928,0.627) (-0.952,0.660) (-0.976,0.693) (-0.999,0.727) (-1.022,0.761) (-1.044,0.795) (-1.065,0.830) (-1.085,0.865) (-1.105,0.901) (-1.125,0.937) (-1.143,0.974) (-1.161,1.010) (-1.178,1.047) (-1.195,1.085) (-1.211,1.122) (-1.226,1.160) (-1.240,1.199) (-1.254,1.237) (-1.267,1.276) (-1.279,1.315) (-1.291,1.354) (-1.302,1.393) (-1.312,1.433) (-1.321,1.473) (-1.330,1.513) (-1.338,1.553) (-1.345,1.593) (-1.351,1.633) (-1.357,1.674) (-1.362,1.714) (-1.366,1.755) (-1.369,1.796) (-1.372,1.837) (-1.374,1.877) (-1.374,1.918) (-1.375,1.959) (-1.374,2.000)};
\addplot[lineaE, flechaen=0.500] coordinates {(0.763,0.650) (0.780,0.687) (0.797,0.723) (0.812,0.761) (0.827,0.799) (0.840,0.837) (0.853,0.875) (0.864,0.914) (0.874,0.954) (0.883,0.993) (0.891,1.033) (0.898,1.073) (0.904,1.113) (0.909,1.153) (0.913,1.193) (0.915,1.234) (0.916,1.274) (0.916,1.315) (0.915,1.355) (0.913,1.396) (0.909,1.436) (0.904,1.476) (0.898,1.516) (0.891,1.556) (0.882,1.596) (0.873,1.635) (0.861,1.674) (0.849,1.712) (0.835,1.750) (0.820,1.788) (0.804,1.825) (0.786,1.862) (0.767,1.897) (0.747,1.932) (0.725,1.967) (0.702,2.000)};
\addplot[lineaE, flechaen=0.500] coordinates {(-0.702,2.000) (-0.725,1.967) (-0.747,1.932) (-0.767,1.897) (-0.786,1.862) (-0.804,1.825) (-0.820,1.788) (-0.835,1.750) (-0.849,1.712) (-0.861,1.674) (-0.873,1.635) (-0.882,1.596) (-0.891,1.556) (-0.898,1.516) (-0.904,1.476) (-0.909,1.436) (-0.913,1.396) (-0.915,1.355) (-0.916,1.315) (-0.916,1.274) (-0.915,1.234) (-0.913,1.193) (-0.909,1.153) (-0.904,1.113) (-0.898,1.073) (-0.891,1.033) (-0.883,0.993) (-0.874,0.954) (-0.864,0.914) (-0.853,0.875) (-0.840,0.837) (-0.827,0.799) (-0.812,0.761) (-0.797,0.723) (-0.780,0.687) (-0.763,0.650)};
\addplot[lineaE, flechaen=0.708, flechaen=0.292] coordinates {(0.633,0.775) (0.637,0.816) (0.640,0.856) (0.641,0.896) (0.641,0.937) (0.639,0.978) (0.635,1.018) (0.629,1.058) (0.622,1.098) (0.613,1.137) (0.602,1.176) (0.589,1.215) (0.574,1.253) (0.557,1.289) (0.538,1.325) (0.518,1.360) (0.496,1.394) (0.471,1.427) (0.445,1.458) (0.418,1.487) (0.388,1.515) (0.357,1.541) (0.324,1.565) (0.290,1.587) (0.254,1.606) (0.217,1.623) (0.179,1.637) (0.140,1.649) (0.101,1.657) (0.061,1.663) (0.020,1.666) (-0.020,1.666) (-0.061,1.663) (-0.101,1.657) (-0.140,1.649) (-0.179,1.637) (-0.217,1.623) (-0.254,1.606) (-0.290,1.587) (-0.324,1.565) (-0.357,1.541) (-0.388,1.515) (-0.418,1.487) (-0.445,1.458) (-0.471,1.427) (-0.496,1.394) (-0.518,1.360) (-0.538,1.325) (-0.557,1.289) (-0.574,1.253) (-0.589,1.215) (-0.602,1.176) (-0.613,1.137) (-0.622,1.098) (-0.629,1.058) (-0.635,1.018) (-0.639,0.978) (-0.641,0.937) (-0.641,0.896) (-0.640,0.856) (-0.637,0.816) (-0.633,0.775)};
\addplot[lineaE, flechaen=0.500] coordinates {(0.423,0.908) (0.406,0.946) (0.387,0.982) (0.365,1.017) (0.340,1.050) (0.313,1.081) (0.283,1.110) (0.250,1.136) (0.216,1.159) (0.180,1.178) (0.142,1.194) (0.102,1.206) (0.062,1.215) (0.021,1.219) (-0.021,1.219) (-0.062,1.215) (-0.102,1.206) (-0.142,1.194) (-0.180,1.178) (-0.216,1.159) (-0.250,1.136) (-0.283,1.110) (-0.313,1.081) (-0.340,1.050) (-0.365,1.017) (-0.387,0.982) (-0.406,0.946) (-0.423,0.908)};
\addplot[lineaE, flechaen=0.396] coordinates {(2.600,-1.454) (2.571,-1.426) (2.542,-1.398) (2.513,-1.370) (2.484,-1.342) (2.455,-1.314) (2.425,-1.286) (2.395,-1.259) (2.365,-1.232) (2.335,-1.205) (2.305,-1.178) (2.274,-1.151) (2.244,-1.125) (2.213,-1.098) (2.182,-1.072) (2.151,-1.047) (2.120,-1.021) (2.089,-0.995) (2.057,-0.970) (2.025,-0.945) (1.993,-0.920) (1.961,-0.895) (1.929,-0.871) (1.897,-0.847) (1.864,-0.823) (1.831,-0.799) (1.799,-0.775) (1.766,-0.752) (1.732,-0.729) (1.699,-0.706) (1.666,-0.683) (1.632,-0.661) (1.598,-0.639) (1.564,-0.617) (1.530,-0.595) (1.496,-0.573) (1.461,-0.552) (1.427,-0.531) (1.392,-0.511) (1.357,-0.490) (1.322,-0.470) (1.287,-0.450) (1.251,-0.431) (1.216,-0.411) (1.180,-0.392) (1.144,-0.374) (1.108,-0.355) (1.072,-0.337) (1.036,-0.319) (0.999,-0.302) (0.963,-0.285)};
\addplot[lineaE, flechaen=0.604] coordinates {(-0.963,-0.285) (-0.999,-0.302) (-1.036,-0.319) (-1.072,-0.337) (-1.108,-0.355) (-1.144,-0.374) (-1.180,-0.392) (-1.216,-0.411) (-1.251,-0.431) (-1.287,-0.450) (-1.322,-0.470) (-1.357,-0.490) (-1.392,-0.511) (-1.427,-0.531) (-1.461,-0.552) (-1.496,-0.573) (-1.530,-0.595) (-1.564,-0.617) (-1.598,-0.639) (-1.632,-0.661) (-1.666,-0.683) (-1.699,-0.706) (-1.732,-0.729) (-1.766,-0.752) (-1.799,-0.775) (-1.831,-0.799) (-1.864,-0.823) (-1.897,-0.847) (-1.929,-0.871) (-1.961,-0.895) (-1.993,-0.920) (-2.025,-0.945) (-2.057,-0.970) (-2.089,-0.995) (-2.120,-1.021) (-2.151,-1.047) (-2.182,-1.072) (-2.213,-1.098) (-2.244,-1.125) (-2.274,-1.151) (-2.305,-1.178) (-2.335,-1.205) (-2.365,-1.232) (-2.395,-1.259) (-2.425,-1.286) (-2.455,-1.314) (-2.484,-1.342) (-2.513,-1.370) (-2.542,-1.398) (-2.571,-1.426) (-2.600,-1.454)};
\addplot[lineaE, flechaen=0.074] coordinates {(2.052,-2.000) (2.038,-1.962) (2.024,-1.924) (2.009,-1.887) (1.994,-1.849) (1.979,-1.812) (1.963,-1.775) (1.946,-1.738) (1.930,-1.701) (1.913,-1.665) (1.895,-1.628) (1.877,-1.592) (1.859,-1.556) (1.840,-1.520) (1.821,-1.485) (1.801,-1.450) (1.781,-1.415) (1.761,-1.380) (1.740,-1.345) (1.719,-1.311) (1.698,-1.276) (1.676,-1.242) (1.654,-1.209) (1.631,-1.175) (1.608,-1.142) (1.585,-1.109) (1.561,-1.077) (1.537,-1.044) (1.512,-1.012) (1.487,-0.980) (1.462,-0.949) (1.437,-0.918) (1.411,-0.887) (1.384,-0.856) (1.358,-0.826) (1.330,-0.796) (1.303,-0.766) (1.275,-0.737) (1.247,-0.708) (1.219,-0.680) (1.190,-0.651) (1.161,-0.623) (1.131,-0.596) (1.101,-0.569) (1.071,-0.542) (1.040,-0.516) (1.009,-0.490) (0.978,-0.464) (0.946,-0.439) (0.914,-0.415)};
\addplot[lineaE, flechaen=0.926] coordinates {(-0.914,-0.415) (-0.946,-0.439) (-0.978,-0.464) (-1.009,-0.490) (-1.040,-0.516) (-1.071,-0.542) (-1.101,-0.569) (-1.131,-0.596) (-1.161,-0.623) (-1.190,-0.651) (-1.219,-0.680) (-1.247,-0.708) (-1.275,-0.737) (-1.303,-0.766) (-1.330,-0.796) (-1.358,-0.826) (-1.384,-0.856) (-1.411,-0.887) (-1.437,-0.918) (-1.462,-0.949) (-1.487,-0.980) (-1.512,-1.012) (-1.537,-1.044) (-1.561,-1.077) (-1.585,-1.109) (-1.608,-1.142) (-1.631,-1.175) (-1.654,-1.209) (-1.676,-1.242) (-1.698,-1.276) (-1.719,-1.311) (-1.740,-1.345) (-1.761,-1.380) (-1.781,-1.415) (-1.801,-1.450) (-1.821,-1.485) (-1.840,-1.520) (-1.859,-1.556) (-1.877,-1.592) (-1.895,-1.628) (-1.913,-1.665) (-1.930,-1.701) (-1.946,-1.738) (-1.963,-1.775) (-1.979,-1.812) (-1.994,-1.849) (-2.009,-1.887) (-2.024,-1.924) (-2.038,-1.962) (-2.052,-2.000)};
\addplot[lineaE, flechaen=0.500] coordinates {(1.374,-2.000) (1.375,-1.959) (1.374,-1.918) (1.374,-1.877) (1.372,-1.837) (1.369,-1.796) (1.366,-1.755) (1.362,-1.714) (1.357,-1.674) (1.351,-1.633) (1.345,-1.593) (1.338,-1.553) (1.330,-1.513) (1.321,-1.473) (1.312,-1.433) (1.302,-1.393) (1.291,-1.354) (1.279,-1.315) (1.267,-1.276) (1.254,-1.237) (1.240,-1.199) (1.226,-1.160) (1.211,-1.122) (1.195,-1.085) (1.178,-1.047) (1.161,-1.010) (1.143,-0.974) (1.125,-0.937) (1.105,-0.901) (1.085,-0.865) (1.065,-0.830) (1.044,-0.795) (1.022,-0.761) (0.999,-0.727) (0.976,-0.693) (0.952,-0.660) (0.928,-0.627) (0.903,-0.595) (0.877,-0.563) (0.851,-0.532)};
\addplot[lineaE, flechaen=0.500] coordinates {(-0.851,-0.532) (-0.877,-0.563) (-0.903,-0.595) (-0.928,-0.627) (-0.952,-0.660) (-0.976,-0.693) (-0.999,-0.727) (-1.022,-0.761) (-1.044,-0.795) (-1.065,-0.830) (-1.085,-0.865) (-1.105,-0.901) (-1.125,-0.937) (-1.143,-0.974) (-1.161,-1.010) (-1.178,-1.047) (-1.195,-1.085) (-1.211,-1.122) (-1.226,-1.160) (-1.240,-1.199) (-1.254,-1.237) (-1.267,-1.276) (-1.279,-1.315) (-1.291,-1.354) (-1.302,-1.393) (-1.312,-1.433) (-1.321,-1.473) (-1.330,-1.513) (-1.338,-1.553) (-1.345,-1.593) (-1.351,-1.633) (-1.357,-1.674) (-1.362,-1.714) (-1.366,-1.755) (-1.369,-1.796) (-1.372,-1.837) (-1.374,-1.877) (-1.374,-1.918) (-1.375,-1.959) (-1.374,-2.000)};
\addplot[lineaE, flechaen=0.500] coordinates {(0.763,-0.650) (0.780,-0.687) (0.797,-0.723) (0.812,-0.761) (0.827,-0.799) (0.840,-0.837) (0.853,-0.875) (0.864,-0.914) (0.874,-0.954) (0.883,-0.993) (0.891,-1.033) (0.898,-1.073) (0.904,-1.113) (0.909,-1.153) (0.913,-1.193) (0.915,-1.234) (0.916,-1.274) (0.916,-1.315) (0.915,-1.355) (0.913,-1.396) (0.909,-1.436) (0.904,-1.476) (0.898,-1.516) (0.891,-1.556) (0.882,-1.596) (0.873,-1.635) (0.861,-1.674) (0.849,-1.712) (0.835,-1.750) (0.820,-1.788) (0.804,-1.825) (0.786,-1.862) (0.767,-1.897) (0.747,-1.932) (0.725,-1.967) (0.702,-2.000)};
\addplot[lineaE, flechaen=0.500] coordinates {(-0.702,-2.000) (-0.725,-1.967) (-0.747,-1.932) (-0.767,-1.897) (-0.786,-1.862) (-0.804,-1.825) (-0.820,-1.788) (-0.835,-1.750) (-0.849,-1.712) (-0.861,-1.674) (-0.873,-1.635) (-0.882,-1.596) (-0.891,-1.556) (-0.898,-1.516) (-0.904,-1.476) (-0.909,-1.436) (-0.913,-1.396) (-0.915,-1.355) (-0.916,-1.315) (-0.916,-1.274) (-0.915,-1.234) (-0.913,-1.193) (-0.909,-1.153) (-0.904,-1.113) (-0.898,-1.073) (-0.891,-1.033) (-0.883,-0.993) (-0.874,-0.954) (-0.864,-0.914) (-0.853,-0.875) (-0.840,-0.837) (-0.827,-0.799) (-0.812,-0.761) (-0.797,-0.723) (-0.780,-0.687) (-0.763,-0.650)};
\addplot[lineaE, flechaen=0.708, flechaen=0.292] coordinates {(0.633,-0.775) (0.637,-0.816) (0.640,-0.856) (0.641,-0.896) (0.641,-0.937) (0.639,-0.978) (0.635,-1.018) (0.629,-1.058) (0.622,-1.098) (0.613,-1.137) (0.602,-1.176) (0.589,-1.215) (0.574,-1.253) (0.557,-1.289) (0.538,-1.325) (0.518,-1.360) (0.496,-1.394) (0.471,-1.427) (0.445,-1.458) (0.418,-1.487) (0.388,-1.515) (0.357,-1.541) (0.324,-1.565) (0.290,-1.587) (0.254,-1.606) (0.217,-1.623) (0.179,-1.637) (0.140,-1.649) (0.101,-1.657) (0.061,-1.663) (0.020,-1.666) (-0.020,-1.666) (-0.061,-1.663) (-0.101,-1.657) (-0.140,-1.649) (-0.179,-1.637) (-0.217,-1.623) (-0.254,-1.606) (-0.290,-1.587) (-0.324,-1.565) (-0.357,-1.541) (-0.388,-1.515) (-0.418,-1.487) (-0.445,-1.458) (-0.471,-1.427) (-0.496,-1.394) (-0.518,-1.360) (-0.538,-1.325) (-0.557,-1.289) (-0.574,-1.253) (-0.589,-1.215) (-0.602,-1.176) (-0.613,-1.137) (-0.622,-1.098) (-0.629,-1.058) (-0.635,-1.018) (-0.639,-0.978) (-0.641,-0.937) (-0.641,-0.896) (-0.640,-0.856) (-0.637,-0.816) (-0.633,-0.775)};
\addplot[lineaE, flechaen=0.500] coordinates {(0.423,-0.908) (0.406,-0.946) (0.387,-0.982) (0.365,-1.017) (0.340,-1.050) (0.313,-1.081) (0.283,-1.110) (0.250,-1.136) (0.216,-1.159) (0.180,-1.178) (0.142,-1.194) (0.102,-1.206) (0.062,-1.215) (0.021,-1.219) (-0.021,-1.219) (-0.062,-1.215) (-0.102,-1.206) (-0.142,-1.194) (-0.180,-1.178) (-0.216,-1.159) (-0.250,-1.136) (-0.283,-1.110) (-0.313,-1.081) (-0.340,-1.050) (-0.365,-1.017) (-0.387,-0.982) (-0.406,-0.946) (-0.423,-0.908)};
\addplot[lineaE, flechaen=0.625] coordinates {(1.000,0.000) (2.600,0.000)};
\addplot[lineaE, flechaen=0.375] coordinates {(-2.600,0.000) (-1.000,0.000)};

  \draw[texto, line width=0.6pt, dashed] (axis cs:0,0) circle[radius=1];
  \draw[vecP] (axis cs:-0.32,0) -- (axis cs:0.32,0);
  \node[etiqm, font=\footnotesize] at (axis cs:0,0.3) {$\vp$};
  \node[etiqm, font=\footnotesize] at (axis cs:-0.45,-0.55) {$r<a$};
\end{axis}
\end{tikzpicture}
\caption{Campo del dipolo para $r>a$}
\end{subfigure}\hfill
\begin{subfigure}[b]{0.44\textwidth}\centering
\begin{tikzpicture}
\begin{polaraxis}[width=5.6cm, xtick={0,45,90,135,180,225,270,315}, ytick={2,4}, yticklabels={}, ymax=4.3,
   xticklabel={$\pgfmathprintnumber{\tick}^\circ$}, tick label style={font=\tiny},
   grid style={draw=gris, line width=0.3pt}, axis line style={draw=gris}]
  \addplot[curvaE, domain=0:360, samples=181, fill=campoE, fill opacity=0.15] {3*cos(x)^2+1};
\end{polaraxis}
\end{tikzpicture}
\caption{$u_e(r,\theta)\propto 3\cos^2\theta+1$}
\end{subfigure}
\caption{\Cref{prob:dipolo}. (a) Líneas del campo de un dipolo puntual $\vp$ (horizontal) fuera de la esfera de radio $a$; el interior,
donde el campo diverge al acercarse al dipolo, queda excluido del cálculo. (b) A distancia fija, la densidad de energía es cuatro
veces mayor en el eje del dipolo ($\theta=0,\pi$) que en el plano ecuatorial ($\theta=\pi/2$), donde el campo es la mitad.}
\label{fig:p52}
\end{figure}


% -----------------------------------------------------------------------------
\subsection{Problema 5.3. Dos esferas conductoras lejanas}\label{sec:p53}

\begin{problema}[Dos esferas con cargas opuestas]\label{prob:dos-esferas}
Sean dos esferas conductoras de radios $a$ y $b$ separadas una distancia $d$ ($a,b\ll d$) y con cargas $+Q$ y $-Q$, respectivamente. Calcula la energía del
sistema.
\end{problema}

\begin{resolucion}
\idea Sistema de dos conductores (método M3). Como están lejos, cada esfera ve a la otra como una carga puntual y su carga sigue repartida casi
uniformemente: $p_{11}\simeq k/a$, $p_{22}\simeq k/b$, $p_{12}\simeq k/d$ (Tema~3).

\paso{Paso 1: potenciales} Cada esfera tiene el potencial de su propia carga más el de la otra, vista como puntual:
\[
  \phi_1=k\Big(\frac Qa-\frac Qd\Big),\qquad \phi_2=k\Big(-\frac Qb+\frac Qd\Big).
\]
\paso{Paso 2: energía}
\[
  U_e=\frac12\big(Q_1\phi_1+Q_2\phi_2\big)=\frac12Q\big(\phi_1-\phi_2\big)=\frac12\,kQ^2\Big[\frac1a-\frac1d-\Big(-\frac1b+\frac1d\Big)\Big],
\]
\begin{equation}\label{eq:p53}
  \recuadro{U_e=\frac{Q^2}{8\pi\eo}\Big(\frac1a+\frac1b-\frac2d\Big)=\frac{Q^2}{8\pi\eo}\Big(\frac{a+b}{ab}-\frac2d\Big).}
\end{equation}
\end{resolucion}

\begin{observacion}[Energía total y de interacción]\label{obs:p53}
El resultado separa sus dos contribuciones: $\frac{Q^2}{8\pi\eo a}+\frac{Q^2}{8\pi\eo b}$ son las energías de formación de cada esfera (\cref{ej:esfera-aislada}) y
$-\frac{Q^2}{4\pi\eo d}$ es la energía de interacción, la de dos cargas puntuales opuestas, negativa porque se atraen. La energía total es positiva, como
debe ser, porque $d\gg a,b$. Si se hubiera usado la fórmula de cargas puntuales, solo se habría obtenido el último término. Equivalentemente, el sistema
es un condensador de capacidad $C=Q/(\phi_1-\phi_2)=\big[k\big(\frac1a+\frac1b-\frac2d\big)\big]^{-1}$, y $U_e=Q^2/2C$.
\end{observacion}

% -----------------------------------------------------------------------------
\subsection{Problema 5.4. Fuerza entre las placas de un condensador}\label{sec:p54}

\begin{problema}[Fuerza entre placas]\label{prob:fuerza-placas}
Calcula las fuerzas de atracción entre las placas de un condensador plano.
\end{problema}

\begin{resolucion}
\idea Trabajos virtuales con la separación $x$ como coordenada (método M4), a carga constante y a potencial constante; deben coincidir. Sean $S$ el
área, $\pm Q$ las cargas y $C(x)=\eo S/x$.

\paso{A carga constante} $U_e=\dfrac{Q^2}{2C}=\dfrac{Q^2x}{2\eo S}$, y
\[
  F_x=-\frac{\partial U_e}{\partial x}\bigg|_Q=-\frac{Q^2}{2\eo S}.
\]
\paso{A potencial constante} $U_e=\tfrac12CV^2=\dfrac{\eo SV^2}{2x}$, y
\[
  F_x=+\frac{\partial U_e}{\partial x}\bigg|_V=-\frac{\eo SV^2}{2x^2}=-\frac{CV^2}{2x}.
\]
\paso{Coincidencia} Con $Q=CV$ y $C=\eo S/d$ en la separación real $x=d$: $\dfrac{Q^2}{2\eo S}=\dfrac{C^2V^2}{2\eo S}=\dfrac{CV^2}{2d}$. Por tanto,
\begin{equation}\label{eq:p54}
  \recuadro{F_x=-\frac{Q^2}{2\eo S}=-\frac{CV^2}{2d}=-\frac{\sigma^2}{2\eo}\,S .}
\end{equation}
\end{resolucion}

\begin{observacion}[Signo y origen físico]\label{obs:p54}
El signo negativo indica que la fuerza tiende a reducir $x$: las placas se atraen. Es la presión electrostática $\sigma^2/2\eo$ del Tema~3 multiplicada por el
área (\cref{ej:presion}). También puede verse directamente: la carga $Q$ de una placa está en el campo de la otra, $\sigma/2\eo$ (no en el campo total $\sigma/\eo$,
que incluye el de la propia placa), y $F=Q\,\sigma/2\eo$. A carga constante la fuerza no depende de la separación; a potencial constante crece como $1/x^2$
al acercarlas, porque la batería aporta más carga.
\end{observacion}

% Figura: problemas 5.3 (dos esferas lejanas con cargas opuestas) y 5.4 (fuerza
% entre las placas de un condensador plano).
\begin{figure}[t]
\centering
\begin{subfigure}[b]{0.5\textwidth}\centering
\begin{tikzpicture}[scale=0.92]
  \shade[ball color=conductorluz] (0,0) circle (0.45);
  \draw[texto, line width=0.5pt] (0,0) circle (0.45);
  \shade[ball color=conductorluz] (4.6,0) circle (0.3);
  \draw[texto, line width=0.5pt] (4.6,0) circle (0.3);
  \draw[texto!60!fondo, line width=0.4pt, dashed] (0.45,0) -- (4.3,0);
  \node[cpos, font=\small] at (0,0.75) {$+Q$};
  \node[cneg, font=\small] at (4.6,0.6) {$-Q$};
  \draw[texto, line width=0.4pt, -{Stealth[length=1.4mm]}] (0,0) -- (-130:0.45);
  \node[etiq, anchor=north east] at (-130:0.48) {$a$};
  \draw[texto, line width=0.4pt, -{Stealth[length=1.4mm]}] (4.6,0) -- ++(-50:0.3);
  \node[etiq, anchor=north west] at ($(4.6,0)+(-50:0.33)$) {$b$};
  \draw[cota] (0,-0.95) -- (4.6,-0.95) node[midway, below, etiq] {$d\gg a,\,b$};
\end{tikzpicture}
\caption{Problema 5.3}
\end{subfigure}\hfill
\begin{subfigure}[b]{0.46\textwidth}\centering
\begin{tikzpicture}[scale=0.92]
  \fill[conductor] (0,1.6) rectangle (3.6,1.75);
  \fill[conductor] (0,-0.15) rectangle (3.6,0);
  \draw[texto, line width=0.5pt] (0,1.6) rectangle (3.6,1.75) (0,-0.15) rectangle (3.6,0);
  \node[etiq, anchor=west] at (3.7,1.68) {$+Q$};
  \node[etiq, anchor=west] at (3.7,-0.08) {$-Q$};
  \foreach \x in {0.5,1.3,2.1,2.9} \draw[campoE, line width=0.55pt, flechamedio=0.55] (\x,1.55) -- (\x,0.05);
  \draw[vecD] (1.8,1.8) -- (1.8,1.25);
  \draw[vecD] (1.8,-0.2) -- (1.8,0.35);
  \node[etiqm, font=\footnotesize, text=campoD, anchor=west] at (1.88,1.32) {$\vF$};
  \node[etiqm, font=\footnotesize, text=campoD, anchor=west] at (1.88,0.3) {$-\vF$};
  \draw[cota] (-0.35,0) -- (-0.35,1.6) node[midway, left, etiq] {$d$};
  \node[etiq, anchor=south] at (1.0,1.78) {$S$};
\end{tikzpicture}
\caption{Problema 5.4}
\end{subfigure}
\caption{(a) Dos esferas conductoras de radios $a$ y $b$, con cargas $\pm Q$, separadas una distancia $d$ mucho mayor que sus radios
(\cref{prob:dos-esferas}). (b) Condensador plano de área $S$ y separación $d$: las placas se atraen con una fuerza
$F=Q^2/(2\eo S)$ que no depende de $d$ a carga constante (\cref{prob:fuerza-placas}).}
\label{fig:p53-p54}
\end{figure}


% -----------------------------------------------------------------------------
\subsection{Problema 5.5. Equilibrio de la placa central}\label{sec:p55}

\begin{problema}[Tres placas a potenciales fijos]\label{prob:tres-placas}
Tres placas metálicas paralelas de área $S$ están conectadas a potenciales $\phi_1$, $\phi_2$ y $\phi_3$ tales que $\phi_2>\phi_1,\phi_3$, como se indica en la
\cref{fig:p55}. Calcula la posición de equilibrio de la placa central.
\end{problema}

\begin{resolucion}
\idea Las placas están unidas a baterías: trabajos virtuales a potencial constante (método M4), con $x$ la distancia de la placa central a la superior
(placa 1) y $d$ la separación total. Despreciamos los efectos de borde.

\paso{Paso 1: energía en función de $x$} Entre las placas 1 y 2 hay un campo uniforme $E_1=(\phi_2-\phi_1)/x$ en el volumen $Sx$, y entre la 2 y la 3,
$E_2=(\phi_2-\phi_3)/(d-x)$ en el volumen $S(d-x)$. Con la densidad de energía,
\[
  U_e=\frac{\eo}{2}E_1^2\,Sx+\frac{\eo}{2}E_2^2\,S(d-x)=\frac{\eo S}{2}\left[\frac{(\phi_2-\phi_1)^2}{x}+\frac{(\phi_2-\phi_3)^2}{d-x}\right].
\]
Es la energía de dos condensadores, $\frac12C_1V_1^2+\frac12C_2V_2^2$, con $C_1=\eo S/x$ y $C_2=\eo S/(d-x)$.

\paso{Paso 2: fuerza a potencial constante}
\[
  F_x=+\frac{\partial U_e}{\partial x}\bigg|_\phi=\frac{\eo S}{2}\left[-\frac{(\phi_2-\phi_1)^2}{x^2}+\frac{(\phi_2-\phi_3)^2}{(d-x)^2}\right].
\]
El primer término tira de la placa hacia la placa 1 (reduce $x$) y el segundo hacia la placa 3: cada placa atrae a la central.

\paso{Paso 3: equilibrio} $F_x=0$ exige $\dfrac{\phi_2-\phi_1}{x}=\dfrac{\phi_2-\phi_3}{d-x}$ (ambas diferencias son positivas), es decir, que los campos a ambos lados
sean iguales. Despejando,
\begin{equation}\label{eq:p55}
  \recuadro{x_{\mathrm{eq}}=\frac{d\,(\phi_2-\phi_1)}{2\phi_2-\phi_1-\phi_3}.}
\end{equation}
\end{resolucion}

\begin{observacion}[Estabilidad]\label{obs:p55}
Derivando otra vez,
\[
  \frac{\partial F_x}{\partial x}=\eo S\left[\frac{(\phi_2-\phi_1)^2}{x^3}+\frac{(\phi_2-\phi_3)^2}{(d-x)^3}\right]>0 :
\]
si la placa se desplaza hacia abajo, la fuerza la empuja
aún más hacia abajo, y al revés. El equilibrio es inestable, como exige el teorema de Earnshaw: a potencial constante, $-U_e$ hace de energía potencial
mecánica y presenta un máximo en $x_{\mathrm{eq}}$ (\cref{fig:p55}). En un acelerómetro, unos muelles proporcionan la estabilidad (\cref{ej:acelerometro}). Comprobación: si
$\phi_1=\phi_3$, $x_{\mathrm{eq}}=d/2$, como exige la simetría.
\end{observacion}

% Figura: problema 5.5. Tres placas paralelas a potenciales fijos; la central puede
% moverse.  Energía potencial mecánica U_mec = -U_e (a potencial constante) y fuerza.
% Valores de la gráfica: phi1 = 0, phi3 = 0,4, phi2 = 1 (unidades V0), d = 1.
\begin{figure}[t]
\centering
\begin{subfigure}[b]{0.4\textwidth}\centering
\begin{tikzpicture}[scale=0.85, x={(1cm,0cm)}, y={(0.42cm,0.3cm)}, z={(0cm,1cm)}]
  \def\L{3.0}\def\W{2.2}
  \foreach \h/\c/\n in {2.0/conductor!70!fondo/1, 1.15/acentoC!45!conductor/2, 0/conductor!70!fondo/3} {
    \fill[\c, opacity=0.92] (0,0,\h) -- (\L,0,\h) -- (\L,\W,\h) -- (0,\W,\h) -- cycle;
    \draw[texto, line width=0.5pt] (0,0,\h) -- (\L,0,\h) -- (\L,\W,\h) -- (0,\W,\h) -- cycle;
    \node[etiq, anchor=west] at (\L+0.15,\W,\h) {$\phi_{\n}$};
  }
  \draw[cota] (-0.3,0,2.0) -- (-0.3,0,1.15) node[midway, left, etiq] {$x$};
  \draw[cota] (-0.3,0,1.15) -- (-0.3,0,0) node[midway, left, etiq] {$d-x$};
  \draw[cota] (-1.85,0,2.0) -- (-1.85,0,0) node[midway, left, etiq] {$d$};
  \draw[vecD] (1.5,1.1,1.15) -- (1.5,1.1,1.65);
  \node[etiqm, font=\footnotesize, text=campoD, anchor=west] at (1.55,1.1,1.5) {$F$};
  \node[etiqt, anchor=north] at (\L/2,0,-0.15) {área $S$};
\end{tikzpicture}
\caption{Geometría}
\end{subfigure}\hfill
\begin{subfigure}[b]{0.56\textwidth}\centering
\begin{tikzpicture}
\begin{axis}[width=7.4cm, height=5.0cm, xmin=0.18, xmax=0.92, ymin=-3.2, ymax=-0.85, grid=major,
   xlabel={$x/d$}, ylabel={$U_{\mathrm{mec}}=-U_e$ \ (unid. $\eo SV_0^2/d$)}, xtick={0.2,0.4,0.625,0.8},
   xticklabels={$0{,}2$,$0{,}4$,$x_{\mathrm{eq}}$,$0{,}8$}, ytick={-3,-2,-1}, clip=false]
  \addplot[curvaD, domain=0.18:0.92, samples=120] {-0.5*(1/x + 0.36/(1-x))};
  \draw[equipot] (axis cs:0.625,-3.2) -- (axis cs:0.625,-1.28);
  \addplot[only marks, mark=*, mark size=1.7pt, color=acentoC] coordinates {(0.625,-1.28)};
  \node[etiqf, font=\scriptsize, anchor=south] at (axis cs:0.625,-1.2) {máximo: equilibrio inestable};
  \draw[vecE] (axis cs:0.52,-1.55) -- (axis cs:0.38,-1.8);
  \draw[vecE] (axis cs:0.73,-1.55) -- (axis cs:0.82,-1.85);
\end{axis}
\end{tikzpicture}
\caption{Energía potencial de la placa central}
\end{subfigure}
\caption{\Cref{prob:tres-placas}. (a) Tres placas de área $S$ a potenciales $\phi_1$, $\phi_2$ y $\phi_3$, con $\phi_2>\phi_1,\phi_3$; la central está a
distancia $x$ de la superior. (b) A potencial constante, la fuerza es $F=+\partial U_e/\partial x$, de modo que el papel de energía
potencial mecánica lo desempeña $-U_e$. Con $\phi_1=0$, $\phi_3=0{,}4V_0$ y $\phi_2=V_0$, $-U_e$ tiene un máximo en
$x_{\mathrm{eq}}=d(\phi_2-\phi_1)/(2\phi_2-\phi_1-\phi_3)=0{,}625\,d$: cualquier desplazamiento aleja la placa del equilibrio (flechas), como
predice el teorema de Earnshaw.}
\label{fig:p55}
\end{figure}


% -----------------------------------------------------------------------------
\setcounter{problema}{6}
\subsection{Problema 5.7. Barra cargada en el eje de un anillo}\label{sec:p57}

\begin{problema}[Barra y anillo]\label{prob:barra-anillo}
Una barra cilíndrica delgada está uniformemente cargada con una carga $Q$ y tiene una longitud $2\ell$. A una distancia $z_0$ del centro de la barra y
perpendicular a esta se coloca un anillo de radio $a$ con carga $Q'$ uniformemente repartida. ¿Qué fuerza experimenta la barra?
\end{problema}

\begin{resolucion}
\idea La barra está sobre el eje del anillo, donde conocemos su campo. La fuerza sobre la barra es la suma de las fuerzas sobre cada elemento de carga en el
campo del anillo (método M5); la comprobamos con la energía de interacción. Tomamos el anillo en el plano $z=0$ y la barra sobre el eje $z$, entre
$z_0-\ell$ y $z_0+\ell$ (\cref{fig:p57}).

\paso{Paso 1: campo del anillo en su eje}~\cite{T1}
\[
  \vE_{\mathrm{anillo}}(z)=\frac{Q'}{4\pi\eo}\,\frac{z}{(a^2+z^2)^{3/2}}\,\uv z .
\]
\paso{Paso 2: fuerza sobre la barra} Con $\lambda=Q/2\ell$ y $\dd q=\lambda\,\dd z$,
\[
  \vF=\int_{z_0-\ell}^{z_0+\ell}\lambda\,\vE_{\mathrm{anillo}}(z)\,\dd z=\frac{\lambda Q'}{4\pi\eo}\left[-\frac{1}{\sqrt{a^2+z^2}}\right]_{z_0-\ell}^{z_0+\ell}\uv z ,
\]
\begin{equation}\label{eq:p57}
  \recuadro{\vF=\frac{QQ'}{8\pi\eo\ell}\left[\frac{1}{\sqrt{a^2+(z_0-\ell)^2}}-\frac{1}{\sqrt{a^2+(z_0+\ell)^2}}\right]\uv z .}
\end{equation}
\paso{Comprobación con la energía de interacción} Con $\phi_{\mathrm{anillo}}(z)=\dfrac{Q'}{4\pi\eo\sqrt{a^2+z^2}}$, la energía de interacción es
$\Uint(z_0)=\int_{z_0-\ell}^{z_0+\ell}\lambda\,\phi_{\mathrm{anillo}}\,\dd z$, y por la regla de Leibniz,
\[
  F_z=-\frac{\dd\Uint}{\dd z_0}=-\lambda\big[\phi_{\mathrm{anillo}}(z_0+\ell)-\phi_{\mathrm{anillo}}(z_0-\ell)\big],
\]
que es \eqref{eq:p57}.
\end{resolucion}

\begin{observacion}[Límites]\label{obs:p57}
Si $z_0=0$, la fuerza es nula por simetría. Si $z_0\gg\ell,a$, con $f(z)=(a^2+z^2)^{-1/2}$, $f(z_0-\ell)-f(z_0+\ell)\simeq-2\ell f'(z_0)\simeq2\ell/z_0^2$, y $F\simeq QQ'/(4\pi\eo z_0^2)$: la ley de
Coulomb entre dos cargas puntuales. La fuerza sobre el anillo es la opuesta, por la tercera ley de Newton.
\end{observacion}

% Figura: problema 5.7. Barra cargada (Q, longitud 2l) a lo largo del eje de un
% anillo (Q', radio a); fuerza sobre la barra en función de la altura z0 de su centro.
\begin{figure}[t]
\centering
\begin{subfigure}[b]{0.42\textwidth}\centering
\begin{tikzpicture}[scale=1.0]
  % anillo en perspectiva
  \draw[texto!45!fondo, line width=1.4pt] (0,0) ellipse (1.6 and 0.42);
  \draw[acentoB, line width=1.6pt] (-1.6,0) arc (180:360:1.6 and 0.42);
  % ejes
  \draw[tenue, line width=0.45pt, -{Stealth[length=1.6mm]}] (0,0) -- (2.4,0) node[below, etiqt] {$y$};
  \draw[tenue, line width=0.45pt, -{Stealth[length=1.6mm]}] (0,0) -- (-1.1,-0.85) node[left, etiqt] {$x$};
  % barra
  \fill[campoP!80!fondo] (-0.06,-0.55) rectangle (0.06,2.75);
  \draw[texto, line width=0.4pt] (-0.06,-0.55) rectangle (0.06,2.75);
  \draw[tenue, line width=0.45pt, -{Stealth[length=1.6mm]}] (0,2.75) -- (0,3.25) node[left, etiqt] {$z$};
  \fill[texto] (0,1.1) circle (0.035);
  \draw[cota] (0.45,-0.55) -- (0.45,2.75) node[pos=0.78, right, etiq] {$2\ell$};
  \draw[cota] (-0.4,0) -- (-0.4,1.1) node[midway, left, etiq] {$z_0$};
  \draw[texto, line width=0.4pt, -{Stealth[length=1.4mm]}] (0,0) -- (1.05,-0.33);
  \node[etiq, anchor=north] at (0.75,-0.38) {$a$};
  \node[etiq, text=campoP, anchor=west] at (0.1,2.35) {$Q$};
  \node[etiq, text=acentoB, anchor=north west] at (1.25,-0.32) {$Q'$};
\end{tikzpicture}
\caption{Geometría}
\end{subfigure}\hfill
\begin{subfigure}[b]{0.54\textwidth}\centering
\begin{tikzpicture}
\begin{axis}[width=7.2cm, height=5.0cm, xmin=-4, xmax=4, ymin=-1.75, ymax=1.75, grid=major,
   xlabel={$z_0/\ell$}, ylabel={$F_z\big/\frac{QQ'}{8\pi\eo\ell^2}$}, xtick={-4,-2,0,2,4}, ytick={-1.5,-1,-0.5,0,0.5,1,1.5},
   legend style={at={(0.02,0.98)}, anchor=north west, font=\scriptsize}]
  \addplot[curvaE, domain=-4:4, samples=241] {1/sqrt(0.25+(x-1)^2)-1/sqrt(0.25+(x+1)^2)};
  \addlegendentry{$a=\ell/2$}
  \addplot[curvaD, domain=-4:4, samples=241] {1/sqrt(1+(x-1)^2)-1/sqrt(1+(x+1)^2)};
  \addlegendentry{$a=\ell$}
\end{axis}
\end{tikzpicture}
\caption{Fuerza sobre la barra}
\end{subfigure}
\caption{\Cref{prob:barra-anillo}. (a) Barra de longitud $2\ell$ y carga $Q$ sobre el eje de un anillo de radio $a$ y carga $Q'$, con su centro a altura
$z_0$. (b) Fuerza axial sobre la barra (cargas del mismo signo): es repulsiva, impar en $z_0$ y nula cuando la barra está centrada
en el anillo. Para $\abs{z_0}\gg\ell,a$ decae como $QQ'/(4\pi\eo z_0^2)$, la ley de Coulomb entre dos cargas puntuales.}
\label{fig:p57}
\end{figure}


% -----------------------------------------------------------------------------
\subsection{Problema 5.8. Fuerza sobre un dipolo}\label{sec:p58}
% Figura: problemas 5.8 (dipolo rígido en un campo no uniforme) y 5.9 (cilindro
% interior de un condensador coaxial parcialmente introducido, a potencial fijo).
\begin{figure}[t]
\centering
\begin{subfigure}[b]{0.4\textwidth}\centering
\begin{tikzpicture}[scale=1.0]
  % líneas de un campo no uniforme (divergentes desde la izquierda)
  \foreach \a in {-24,-12,0,12,24}
    \draw[campoE!55!fondo, line width=0.45pt, flechamedio=0.92] (-1.9,0) -- ++(\a:4.3);
  \coordinate (m) at (0.2,-0.25); \coordinate (p) at (1.25,0.35);
  \draw[texto, line width=0.6pt] (m) -- (p);
  \fill[carganeg] (m) circle (0.075); \fill[cargapos] (p) circle (0.075);
  \draw[vecP, line width=1.1pt] (m) -- ($(m)!0.62!(p)$);
  \node[etiqf, font=\footnotesize, anchor=south east] at ($(m)!0.42!(p)+(0,0.08)$) {$\vp=q\vec d$};
  \draw[draw=cargapos, line width=1.0pt, -{Stealth[length=2mm]}] (p) -- ++(14:0.85);
  \draw[draw=carganeg, line width=1.0pt, -{Stealth[length=2mm]}] (m) -- ++(186:1.05);
  \node[etiqf, font=\footnotesize, anchor=south] at ($(p)+(0.55,0.25)$) {$q\vE(\vr+\vec d)$};
  \node[etiqf, font=\footnotesize, anchor=north] at ($(m)+(-0.5,-0.22)$) {$-q\vE(\vr)$};
\end{tikzpicture}
\caption{Problema 5.8}
\end{subfigure}\hfill
\begin{subfigure}[b]{0.56\textwidth}\centering
\begin{tikzpicture}[scale=0.95]
  % cilindro exterior (radio R1, longitud l)
  \fill[dielec] (0,-1.0) -- (4.2,-1.0) arc (-90:90:0.3 and 1.0) -- (0,1.0) arc (90:270:0.3 and 1.0);
  \draw[borde] (0,-1.0) -- (4.2,-1.0) (0,1.0) -- (4.2,1.0);
  \draw[borde] (0,0) ellipse (0.3 and 1.0);
  \draw[borde] (4.2,0) ellipse (0.3 and 1.0);
  % cilindro interior (radio R2), introducido una longitud x
  \fill[conductor] (2.4,-0.45) rectangle (5.9,0.45);
  \fill[conductor] (5.9,0) ellipse (0.135 and 0.45);
  \draw[texto, line width=0.6pt] (2.4,-0.45) -- (5.9,-0.45) (2.4,0.45) -- (5.9,0.45);
  \draw[texto, line width=0.6pt] (5.9,0) ellipse (0.135 and 0.45);
  \draw[texto, line width=0.6pt] (2.4,0) ellipse (0.135 and 0.45);
  \draw[texto!60!fondo, line width=0.4pt, dashed] (-0.6,0) -- (6.4,0);
  % cotas
  \draw[cota] (2.4,-1.3) -- (4.2,-1.3) node[midway, below, etiq] {$x$};
  \draw[cota] (0,-1.85) -- (4.2,-1.85) node[midway, below, etiq] {$\ell$};
  \draw[texto, line width=0.4pt, -{Stealth[length=1.4mm]}] (0.5,0) -- (0.5,0.98);
  \node[etiq, anchor=west] at (0.55,0.6) {$R_1$};
  \draw[texto, line width=0.4pt, -{Stealth[length=1.4mm]}] (5.0,0) -- (5.0,0.43);
  \node[font=\small, text=black, anchor=west] at (5.05,0.2) {$R_2$};
  % fuerza
  \draw[vecD] (5.25,-0.75) -- (4.45,-0.75);
  \node[etiq, text=campoD, anchor=north] at (4.85,-0.8) {$F$};
  % batería
  \draw[texto, line width=0.55pt] (4.2,1.0) -- (4.2,1.55) -- (5.0,1.55);
  \draw[texto, line width=1.1pt] (5.05,1.33) -- (5.05,1.77);
  \draw[texto, line width=0.6pt] (5.25,1.42) -- (5.25,1.68);
  \draw[texto, line width=0.55pt] (5.25,1.55) -- (6.3,1.55) -- (6.3,0) -- (6.04,0);
  \node[etiq, anchor=south] at (5.15,1.8) {$V_0$};
\end{tikzpicture}
\caption{Problema 5.9}
\end{subfigure}
\caption{(a) Dipolo rígido $\vp=q\vec d$ en un campo no uniforme: las fuerzas sobre sus cargas no se compensan y la resultante es
$(\vp\cdot\grad)\vE$ (\cref{prob:dipolo-fuerza}). (b) Condensador cilíndrico de longitud $\ell$ y radios $R_1>R_2$ con el cilindro interior
introducido una longitud $x$; las armaduras están a diferencia de potencial $V_0$. El cilindro interior es atraído hacia dentro con una
fuerza constante (\cref{prob:coaxial}).}
\label{fig:p58-p59}
\end{figure}


\begin{problema}[Fuerza sobre un dipolo rígido]\label{prob:dipolo-fuerza}
Demuestra que la fuerza que actúa sobre un dipolo rígido de momento $\vp$ sometido a un campo $\vE$ es de la forma
\[
  \vF=\big(\vp\cdot\grad\big)\vE=\Big(p_x\frac{\partial}{\partial x}+p_y\frac{\partial}{\partial y}+p_z\frac{\partial}{\partial z}\Big)\big(E_x\uv x+E_y\uv y+E_z\uv z\big).
\]
\end{problema}

\begin{resolucion}
\idea Dos demostraciones: sumando las fuerzas sobre las dos cargas, y derivando la energía de interacción del dipolo con el campo (método M5).

\paso{Demostración 1: fuerzas sobre las cargas} Sea el dipolo formado por $-q$ en $\vr$ y $+q$ en $\vr+\vec d$, con $\vp=q\vec d$ fijo (\cref{fig:p58-p59}\,(a)). La fuerza total
es $\vF=q\big[\vE(\vr+\vec d)-\vE(\vr)\big]$. Desarrollando cada componente en serie de Taylor hasta primer orden,
\[
  E_i(\vr+\vec d)=E_i(\vr)+\vec d\cdot\grad E_i+\mathcal O(d^2)\quad\Longrightarrow\quad F_i=q\,\vec d\cdot\grad E_i=\big(\vp\cdot\grad\big)E_i ,
\]
y en el límite de dipolo puntual ($d\to0$ con $p$ constante) el resto se anula.

\paso{Demostración 2: energía} La energía de interacción del dipolo con el campo externo es $\Uint=q\phi(\vr+\vec d)-q\phi(\vr)\simeq\vp\cdot\grad\phi=-\vp\cdot\vE$ (Tema~4), y
$\vF=-\grad\Uint=\grad(\vp\cdot\vE)$. Con $\vp$ constante, la identidad del gradiente de un producto escalar da
\[
  \grad(\vp\cdot\vE)=\big(\vp\cdot\grad\big)\vE+\vp\times\big(\rot\vE\big)=\big(\vp\cdot\grad\big)\vE ,
\]
pues el campo electrostático es irrotacional. \qedhere
\end{resolucion}

\begin{observacion}[Campo uniforme y no uniforme]\label{obs:p58}
En un campo uniforme la fuerza es nula y solo queda el par $\vtau=\vp\times\vE$ del Tema~4. En un campo no uniforme, un dipolo alineado con el campo es atraído hacia
las regiones de campo más intenso: es el principio de la dielectroforesis, que se tratará en la segunda parte.
\end{observacion}

% -----------------------------------------------------------------------------
\subsection{Problema 5.9. Condensador cilíndrico con el cilindro interior móvil}\label{sec:p59}

\begin{problema}[Condensador coaxial]\label{prob:coaxial}
En un condensador cilíndrico como el de la \cref{fig:p58-p59}\,(b) se conectan las armaduras a una diferencia de potencial $V_0$. Despreciando el efecto de bordes,
calcula la fuerza que experimenta el cilindro de menor radio.
\end{problema}

\begin{resolucion}
\idea Batería: trabajos virtuales a potencial constante (método M4) con la longitud introducida $x$ como coordenada. Basta conocer $C(x)$.

\paso{Paso 1: capacidad} Despreciando los bordes, la parte solapada, de longitud $x$, es un condensador cilíndrico de radios $R_1>R_2$, y el resto del sistema aporta una
capacidad $C_0$ que no depende de $x$:
\[
  C(x)=\frac{2\pi\eo\,x}{\ln(R_1/R_2)}+C_0 .
\]
\paso{Paso 2: fuerza} Por el \cref{cor:fuerza-condensador},
\begin{equation}\label{eq:p59}
  \recuadro{F_x=\frac12V_0^2\,\frac{\dd C}{\dd x}=\frac{\pi\eo V_0^2}{\ln(R_1/R_2)} ,}
\end{equation}
positiva: tiende a aumentar $x$, es decir, a introducir el cilindro interior.
\end{resolucion}

\begin{observacion}[Una fuerza constante]\label{obs:p59}
La fuerza no depende de $x$: al introducir el cilindro, solo se añade un tramo más de condensador idéntico a los anteriores; los efectos de borde, que
dependen de la geometría de los extremos, no cambian. A carga constante se obtendría $F=\frac{Q^2}{2C^2}\dd C/\dd x$, que sí depende de $x$ a través de $C$, pero que coincide
con \eqref{eq:p59} en cada configuración con $V_0=Q/C$ (\cref{prop:equivalencia}).
\end{observacion}

% -----------------------------------------------------------------------------
\subsection{Resumen de resultados y lista de control}\label{sec:resumen-seminarios}

\begin{table}[H]
\centering
\caption{Resultados de los seminarios del Tema 5 (primera parte).}\label{tab:resultados}
\small
\renewcommand{\arraystretch}{1.95}
\begin{tabularx}{\textwidth}{@{}c>{\raggedright\arraybackslash}Xc@{}}
\toprule
Problema & Resultado & Método\\
\midrule
5.1(a) & $U_e=\dfrac{3Q^2}{20\pi\eo a}$ \quad ($\tfrac16$ dentro, $\tfrac56$ fuera) & M1, M2\\
5.2 & $U_e(r>a)=\dfrac{p^2}{12\pi\eo a^3}$ & M2\\
5.3 & $U_e=\dfrac{Q^2}{8\pi\eo}\Big(\dfrac{a+b}{ab}-\dfrac2d\Big)$ & M3\\
5.4 & $F=-\dfrac{Q^2}{2\eo S}=-\dfrac{CV^2}{2d}$ (atractiva) & M4\\
5.5 & $x_{\mathrm{eq}}=\dfrac{d(\phi_2-\phi_1)}{2\phi_2-\phi_1-\phi_3}$ (inestable) & M4\\
5.7 & $\vF=\dfrac{QQ'}{8\pi\eo\ell}\Big[\dfrac{1}{\sqrt{a^2+(z_0-\ell)^2}}-\dfrac{1}{\sqrt{a^2+(z_0+\ell)^2}}\Big]\uv z$ & M5\\
5.8 & $\vF=(\vp\cdot\grad)\vE=\grad(\vp\cdot\vE)$ & M5\\
5.9 & $F=\dfrac{\pi\eo V_0^2}{\ln(R_1/R_2)}$ (hacia dentro) & M4\\
\bottomrule
\end{tabularx}
\end{table}

\begin{metodo}{Lista de control}
\begin{itemize}
  \item ¿Es energía de interacción ($\frac12\sum q_i\phi_i$, sin términos propios) o total ($\frac12\int\rho\phi$, $\frac{\eo}{2}\int E^2$, $\frac12\sum Q_i\phi_i$)? ¿Tiene el signo esperado?
  \item Con $\frac{\eo}{2}\int E^2$, ¿he integrado en \emph{todo} el espacio donde hay campo, dentro y fuera?
  \item ¿He usado el campo \emph{total} al cuadrado, sin sumar energías por separado?
  \item En $\frac12\sum Q_i\phi_i$, ¿es $\phi_i$ el potencial total del conductor, incluida su propia carga?
  \item En las fuerzas: ¿está el sistema aislado ($-\partial_\xi U_e|_Q$) o conectado a baterías ($+\partial_\xi U_e|_\phi$)? ¿He derivado con la variable correcta constante?
  \item ¿Coincide la fuerza por las dos vías y tiene el signo esperado ($C$ tiende a aumentar)? ¿Recupero los límites (cargas lejanas, simetría)?
\end{itemize}
\end{metodo}

% =============================================================================
\clearpage
\appendix
\renewcommand{\thesection}{\Alph{section}}
\section{Formulario del Tema 5 (primera parte)}\label{sec:formulario}
% =============================================================================

La \cref{tab:formulario} reúne las expresiones centrales de esta parte del tema, con la referencia al enunciado en que se demuestran.

\begingroup
\centering
\captionof{table}{Formulario del Tema 5. Energía y fuerzas, primera parte ($k=1/4\pi\eo$, $\phi(\infty)=0$).}
\label{tab:formulario}
\small
\renewcommand{\arraystretch}{1.75}
\begin{tabularx}{\textwidth}{@{}>{\raggedright\arraybackslash}p{4.2cm}>{\raggedright\arraybackslash}X>{\raggedright\arraybackslash}p{2.9cm}@{}}
\toprule
Magnitud o resultado & Expresión & Referencia\\
\midrule
\multicolumn{3}{@{}l}{\textit{Energía de cargas y distribuciones}}\\
Trabajo reversible & $\Wext(A\to B)=q\big[\phi(B)-\phi(A)\big]$ & \cref{lem:trabajo}\\
Cargas puntuales (interacción) & $U_e=\dfrac12\displaystyle\sum_{i\neq j}\frac{kq_iq_j}{r_{ij}}=\frac12\sum_iq_i\phi_i$ & \cref{prop:puntuales}\\
Distribución continua (total) & $U_e=\dfrac12\displaystyle\int_V\rho\,\phi\,\dd V$,\quad $\dfrac12\displaystyle\int_S\sigma\,\phi\,\dd S$ & \cref{prop:continua}\\
\multicolumn{3}{@{}l}{\textit{Energía del campo}}\\
Energía y densidad & $U_e=\dfrac{\eo}{2}\displaystyle\int_{\mathbb R^3}E^2\,\dd V$,\quad $u_e=\dfrac{\eo}{2}E^2$,\quad $U_e\geq0$ & \cref{teo:energia-campo}, \cref{cor:positiva}\\
No aditividad & $U_e=U_1+U_2+\Uint$,\quad $\Uint=\eo\displaystyle\int\vE_1\cdot\vE_2\,\dd V=\int\rho_2\phi_1\,\dd V$ & \cref{prop:no-aditiva}\\
Corteza esférica / esfera maciza & $\dfrac{Q^2}{8\pi\eo a}$\quad / \quad$\dfrac{3Q^2}{20\pi\eo a}$ & \cref{ej:corteza}, \cref{prob:esfera-cargada}\\
\multicolumn{3}{@{}l}{\textit{Conductores}}\\
Sistema de conductores & $U_e=\dfrac12\displaystyle\sum_iQ_i\phi_i=\frac12\sum_{i,j}p_{ij}Q_iQ_j=\frac12\sum_{i,j}c_{ij}\phi_i\phi_j$ & \cref{prop:conductores}\\
Condensador & $U_e=\dfrac12QV=\dfrac12CV^2=\dfrac{Q^2}{2C}$ & \cref{cor:condensador}\\
\multicolumn{3}{@{}l}{\textit{Fuerzas}}\\
Earnshaw & no hay equilibrio estable solo con fuerzas electrostáticas & \cref{teo:earnshaw}\\
Carga constante & $\vFe=-\grad U_e|_Q$,\ \ $\tau_\theta=-\partial_\theta U_e|_Q$,\ \ $F_\xi=-\tfrac12\sum_{i,j}Q_iQ_j\,\partial_\xi p_{ij}$ & \cref{teo:fuerza-Q}\\
Potencial constante & $\vFe=+\grad U_e|_\phi$,\quad $\tau_\theta=+\partial_\theta U_e|_\phi$,\quad $\dd\Wbat=2\,\dd U_e$ & \cref{teo:fuerza-phi}\\
Equivalencia & $-\partial_\xi U_e|_Q=+\partial_\xi U_e|_\phi$ & \cref{prop:equivalencia}\\
Condensador & $F_\xi=\dfrac12V^2\dfrac{\dd C}{\dd\xi}=\dfrac{Q^2}{2C^2}\dfrac{\dd C}{\dd\xi}$ & \cref{cor:fuerza-condensador}\\
Dipolo en un campo & $\Uint=-\vp\cdot\vE$,\quad $\vF=(\vp\cdot\grad)\vE$ & \cref{prob:dipolo-fuerza}\\
\bottomrule
\end{tabularx}
\par\endgroup

% =============================================================================
\begin{thebibliography}{99}
\setlength{\itemsep}{0.25em}
\small
\bibitem{T1} D.~Martín Guerrero. \emph{Leyes básicas de la electrostática en el vacío} (Electromagnetismo I, Tema 1). Apuntes formalizados con
  seminarios prácticos. Grado en Física, Universidad de Salamanca, 2026.
\bibitem{T2} D.~Martín Guerrero. \emph{Desarrollo multipolar} (Electromagnetismo I, Tema 2). Apuntes formalizados con seminarios prácticos. Grado
  en Física, Universidad de Salamanca, 2026.
\bibitem{T3} D.~Martín Guerrero. \emph{Conductores en equilibrio electrostático} (Electromagnetismo I, Tema 3). Apuntes formalizados con seminarios
  prácticos. Grado en Física, Universidad de Salamanca, 2026.
\bibitem{T4} D.~Martín Guerrero. \emph{Materiales dieléctricos} (Electromagnetismo I, Tema 4). Apuntes formalizados con seminarios prácticos. Grado
  en Física, Universidad de Salamanca, 2026.
\bibitem{T0} D.~Martín Guerrero. \emph{Herramientas matemáticas para el electromagnetismo} (Electromagnetismo I, Tema 0). Apuntes formalizados con
  seminarios prácticos. Grado en Física, Universidad de Salamanca, 2026.
\bibitem{DFA-teoria} Departamento de Física Aplicada. \emph{Electromagnetismo I. Tema V: Energía y fuerzas}. Apuntes de teoría, §V.1--V.5. Material de
  la asignatura. Grado en Física, Universidad de Salamanca, 2025.
\bibitem{DFA-figuras} Departamento de Física Aplicada. \emph{Electromagnetismo I. Tema 5: Energía y fuerzas (figuras)}. Material de la asignatura.
  Grado en Física, Universidad de Salamanca, 2025.
\bibitem{DFA-problemas} Departamento de Física Aplicada. \emph{Electromagnetismo I. Tema 5: Energía y fuerzas (relación de problemas)}. Material de la
  asignatura. Grado en Física, Universidad de Salamanca, curso 2025-2026.
\bibitem{Guia-T5} D.~Martín Guerrero. \emph{Tema 5. Energía y fuerzas (I): formulario, métodos y fichas de problemas}. Guías de estudio de
  Electromagnetismo I. Grado en Física, Universidad de Salamanca, 2026.
\bibitem{Griffiths} D.~J. Griffiths. \emph{Introduction to Electrodynamics}. 4.ª ed. §2.4: «Work and energy in electrostatics»; §2.5.4: fuerza sobre
  conductores; §4.4.3: fuerzas sobre dieléctricos. Cambridge University Press, 2017.
\bibitem{Wangsness} R.~K. Wangsness. \emph{Electromagnetic Fields}. 2.ª ed. Cap.~7: «Electrostatic energy». John Wiley \& Sons, 1986.
\bibitem{Purcell} E.~M. Purcell y D.~J. Morin. \emph{Electricity and Magnetism}. 3.ª ed. §1.15: energía de un sistema de cargas; §3.5--3.6: energía y
  fuerzas en condensadores. Cambridge University Press, 2013.
\bibitem{Feynman} R.~P. Feynman, R.~B. Leighton y M.~Sands. \emph{The Feynman Lectures on Physics. Volume II}. Cap.~8: «Electrostatic energy». Addison-Wesley,
  1964.
\bibitem{Jackson} J.~D. Jackson. \emph{Classical Electrodynamics}. 3.ª ed. §1.11: energía potencial electrostática y densidad de energía; §1.13:
  coeficientes de capacidad. John Wiley \& Sons, 1999.
\bibitem{Kellogg} O.~D. Kellogg. \emph{Foundations of Potential Theory}. Cap.~VIII: teoremas del valor medio y principio del máximo. Berlín: Springer, 1929.
\bibitem{Earnshaw} S.~Earnshaw. «On the nature of the molecular forces which regulate the constitution of the luminiferous ether». \emph{Transactions of the
  Cambridge Philosophical Society} 7 (1842), 97--112.
\end{thebibliography}

